\documentclass[11pt,a4paper]{report}

% --- Langue, polices ---------------------------------------------------------------------------
\usepackage[T1]{fontenc}
\usepackage[utf8]{inputenc}
\usepackage{lmodern}
\usepackage[french]{babel}
\usepackage[margin=2.4cm]{geometry}
\usepackage{microtype}
\usepackage{csquotes}
\setlength{\emergencystretch}{3em}

% --- Maths -------------------------------------------------------------------------------------
\usepackage{amsmath,amssymb,mathtools,bm}
\usepackage{siunitx}
\sisetup{locale=FR, group-minimum-digits=4}

% --- Tableaux, listes --------------------------------------------------------------------------
\usepackage{booktabs,tabularx,longtable,array}
\usepackage{enumitem}
\setlist{itemsep=2pt, topsep=4pt}
\newcolumntype{L}{>{\raggedright\arraybackslash}X}

% --- Figures -----------------------------------------------------------------------------------
\usepackage{xcolor}
\usepackage{tikz}
\usetikzlibrary{arrows.meta,positioning,shapes.geometric,shapes.misc,fit,calc,backgrounds,babel}
\usepackage{pgfplots}
\pgfplotsset{compat=1.18}
\usepackage{float}
\usepackage{caption}

% --- Code et algorithmes -----------------------------------------------------------------------
\usepackage{listings}
\usepackage[ruled,vlined,linesnumbered,french]{algorithm2e}
\usepackage[most]{tcolorbox}

\definecolor{codebg}{RGB}{246,247,250}
\definecolor{codekw}{RGB}{0,70,160}
\definecolor{codecom}{RGB}{110,110,110}
\definecolor{codestr}{RGB}{160,60,0}
\definecolor{accent}{RGB}{79,142,247}
\definecolor{doux}{RGB}{120,126,140}

\lstset{
  language=Python, basicstyle=\ttfamily\small, keywordstyle=\color{codekw}\bfseries,
  commentstyle=\color{codecom}\itshape, stringstyle=\color{codestr}, backgroundcolor=\color{codebg},
  frame=single, rulecolor=\color{black!15}, breaklines=true, columns=fullflexible, keepspaces=true,
  showstringspaces=false, xleftmargin=4pt, xrightmargin=4pt, aboveskip=8pt, belowskip=8pt,
  literate={é}{{\'e}}1 {è}{{\`e}}1 {ê}{{\^e}}1 {à}{{\`a}}1 {â}{{\^a}}1 {ç}{{\c{c}}}1
           {ù}{{\`u}}1 {û}{{\^u}}1 {ô}{{\^o}}1 {î}{{\^i}}1 {ï}{{\"i}}1 {É}{{\'E}}1 {À}{{\`A}}1
           {œ}{{\oe}}1 {«}{{\guillemotleft}}1 {»}{{\guillemotright}}1 {→}{{$\to$}}1
           {θ}{{$\theta$}}1 {φ}{{$\varphi$}}1 {·}{{$\cdot$}}1 {×}{{$\times$}}1
}

% --- Encadrés ----------------------------------------------------------------------------------
\newtcolorbox{origine}[1][]{colback=accent!6, colframe=accent!60, boxrule=0.5pt, arc=2pt,
  left=6pt, right=6pt, top=4pt, bottom=4pt, fonttitle=\bfseries\small, title={Origine des valeurs}, #1}
\newtcolorbox{attention}[1][]{colback=orange!7, colframe=orange!70!black, boxrule=0.5pt, arc=2pt,
  left=6pt, right=6pt, top=4pt, bottom=4pt, fonttitle=\bfseries\small, title={Point d'attention}, #1}

% --- Raccourcis --------------------------------------------------------------------------------
% \code{...} : texte machine, coupable après _ . / (sans casser les identifiants ailleurs)
\ExplSyntaxOn
\cs_new_protected:Npn \projet_code:n #1
  {
    \str_set:Nn \l_tmpa_str {#1}
    \str_map_inline:Nn \l_tmpa_str
      { ##1 \str_if_in:nnT { _./ } {##1} { \allowbreak } }
  }
\DeclareRobustCommand{\code}[1]{\texttt{\projet_code:n{#1}}}
\DeclareRobustCommand{\fichier}[1]{\texttt{\projet_code:n{#1}}}
\ExplSyntaxOff
\newcommand{\R}{\mathbb{R}}
\newcommand{\vect}[1]{\bm{#1}}
\newcommand{\T}{^{\top}}
\DeclareMathOperator{\atantwo}{atan2}
\DeclareMathOperator{\med}{med}
\DeclareMathOperator{\MAD}{MAD}
\DeclareMathOperator{\diag}{diag}
\DeclareMathOperator{\tr}{tr}
\DeclareMathOperator{\sign}{sign}
\DeclareMathOperator*{\argmin}{arg\,min}

% Étiquettes d'origine des constantes
\newcommand{\litt}{\textsc{littérature}}
\newcommand{\deriv}{\textsc{dérivée}}
\newcommand{\empi}{\textsc{empirique}}
\newcommand{\conv}{\textsc{convention}}

\usepackage[hidelinks]{hyperref}
\usepackage{bookmark}

\title{\textbf{Eye tracker par webcam pour l'accessibilité}\\[6pt]
  \large Documentation technique complète : modèles, mathématiques, pipeline et implémentation}
\author{Projet vision -- ISAE}
\date{\today}

\begin{document}
\maketitle

\begin{abstract}
Ce document décrit de bout en bout le logiciel développé dans le dépôt : un système qui transforme une webcam ordinaire en dispositif de pointage par le regard, destiné à une personne qui ne peut pas utiliser ses mains, sur PC ou sur Raspberry Pi. L'utilisateur fixe un bouton pendant deux secondes pour déclencher une action (raccourci clavier, texte, programme, broche GPIO, requête HTTP).

Il explique les deux estimateurs de regard successifs : une régression ridge sur des caractéristiques 2D de l'œil (v1), puis un modèle géométrique 3D de la tête et des globes oculaires (v2), avec recalibration continue à chaque clic. Pour chaque algorithme non standard, il donne la dérivation mathématique ; pour chaque constante du code, il indique d'où vient sa valeur (littérature, dérivation, convention ou choix empirique). Il détaille enfin l'organisation en threads, les objectifs de performance, l'interface, la configuration et les actions.
\end{abstract}

\tableofcontents

\chapter{Introduction}
\label{chap:intro}

\section{Objectif du projet}

Le but est de permettre à une personne qui ne peut pas se servir de ses mains de commander un ordinateur, ou des appareils reliés à un Raspberry Pi, \emph{uniquement avec le regard}, en utilisant une webcam ordinaire et sans matériel spécialisé (pas de caméra infrarouge, pas d'éclairage actif, pas de mentonnière).

Le principe d'interaction est celui de la \emph{sélection par fixation} (\emph{dwell selection}) :
\begin{enumerate}
  \item l'écran est découpé en quelques grands boutons (grille $2\times2$ ou $2\times3$) ;
  \item un curseur suit le point regardé ;
  \item lorsque le regard reste sur un bouton pendant \code{duree_clic} secondes (\qty{2}{\second} par défaut), un anneau de progression se remplit puis l'action du bouton se déclenche ;
  \item au centre de l'écran, une \emph{zone morte} ovale permet de reposer les yeux sans rien déclencher.
\end{enumerate}

Les menus, les boutons et leurs actions sont décrits dans un fichier de configuration (YAML, TOML ou JSON) qu'un aidant peut modifier sans toucher au code.

\section{Contraintes de conception}

Le code porte la trace de plusieurs contraintes, qui expliquent la plupart des choix techniques décrits dans ce document.

\begin{description}[style=nextline]
  \item[Matériel bon marché.] Une webcam RGB standard (résolution demandée $640\times480$, \qty{30}{images\per\second}) et un Raspberry Pi 4 ou 5, ou un PC. L'estimation du regard doit donc se contenter de la position de l'iris dans l'image, mesurée à quelques pixels près.
  \item[Calcul limité.] Sur Raspberry Pi, le réseau de neurones de détection du visage occupe déjà une bonne partie du processeur. Le reste du pipeline doit rester léger : calculs vectorisés, rendu graphique incrémental, mise en cache.
  \item[Robustesse aux mouvements de tête.] L'utilisateur ne peut pas garder la tête parfaitement immobile pendant des heures. C'est la raison d'être du tracker v2 (chapitres~\ref{chap:v2geo} et~\ref{chap:v2mod}).
  \item[Tolérance aux erreurs de pointage.] Même bien calibré, un eye tracker webcam a une erreur de plusieurs pour cent de la largeur de l'écran. L'interface utilise donc peu de boutons, de grande taille, et une zone morte.
  \item[Sécurité de l'interaction.] Regarder n'est pas cliquer : c'est le problème du \enquote{toucher de Midas} \cite{jacob1991}. La durée de fixation, le blocage après un clic et la zone morte servent à éviter les déclenchements involontaires.
  \item[Configuration par un non-programmeur.] Les messages d'erreur de configuration sont en français, localisés (chemin de la clé fautive) et accompagnés de suggestions.
\end{description}

\section{Historique du travail réalisé}

Le tableau~\ref{tab:historique} résume l'évolution du dépôt d'après l'historique Git.

\begin{table}[H]
\centering\small
\begin{tabularx}{\textwidth}{@{}l l L@{}}
\toprule
Date & Commit & Contenu \\
\midrule
16/09/2026 & \code{4ca3b57} & Premier prototype \fichier{main.py} : sélection sur une grille $3\times3$ sans calibration, par seuillage de la position de l'iris dans l'orbite ; ajout du modèle MediaPipe. \\
16/09/2026 & \code{7e17919}, \code{a961ffd} & Essais, ajout d'un gain et d'une sensibilité. \\
25/11/2026 & \code{f2ba5a4} & Nouvelle calibration : \fichier{test_xav.py}, tracker \enquote{précision} (repère de l'œil, blendshapes, compensation de la tête, calibration 16 points, ridge, One Euro). \\
26/11/2026 & \code{0cd7e8a} & Interface complète dans \fichier{victor/} : briques réutilisables (\fichier{utils.py}), rendu pygame, configuration YAML/TOML/JSON, actions, mode souris. \\
26/11/2026 & \code{7b26ee3} & Réglage de la disposition de l'écran. \\
26/11/2026 & \code{6d68e68} & Tracker v2 (\fichier{utils_v2.py}, \fichier{main_v2.py}) : modèle géométrique 3D de la tête et des yeux. \\
26/11/2026 & \code{d23ef7e} & Amélioration de la tolérance. \\
28/11/2026 & \code{797601f} & Recalibration continue à chaque clic et paramètres associés. \\
28/11/2026 & \code{2f619bf} & Feuille de route (\fichier{ROADMAP.md}). \\
\bottomrule
\end{tabularx}
\caption{Évolution du dépôt.}
\label{tab:historique}
\end{table}

\section{Organisation du document}

\begin{itemize}
  \item Le chapitre~\ref{chap:archi} présente l'architecture : fichiers, threads, pipeline d'une image et objectifs de performance.
  \item Le chapitre~\ref{chap:mediapipe} décrit ce que fournit MediaPipe et comment le code l'utilise.
  \item Le chapitre~\ref{chap:v1} détaille le tracker v1 (régression ridge).
  \item Les chapitres~\ref{chap:v2geo} et~\ref{chap:v2mod} détaillent le tracker v2 : géométrie 3D, puis modèle de regard et son ajustement.
  \item Le chapitre~\ref{chap:recal} traite de la recalibration continue au clic.
  \item Le chapitre~\ref{chap:traitements} décrit les traitements communs : clignements, filtre One Euro, recentrage.
  \item Les chapitres~\ref{chap:interaction} et~\ref{chap:rendu} décrivent la boucle d'interaction et le rendu graphique.
  \item Les chapitres~\ref{chap:config} et~\ref{chap:actions} décrivent la configuration et les actions.
  \item Le chapitre~\ref{chap:bilan} fait le bilan : limites connues, incohérences relevées dans le code, pistes.
  \item L'annexe~\ref{chap:constantes} récapitule toutes les constantes avec leur origine.
\end{itemize}

\section{Conventions}

\paragraph{Origine des constantes.} Chaque constante numérique du code est classée dans l'une des quatre catégories suivantes :
\begin{description}
  \item[\litt] valeur reprise d'une publication ou d'une mesure physiologique connue ;
  \item[\deriv] valeur qui découle mathématiquement d'autres grandeurs ;
  \item[\conv] valeur conventionnelle ou imposée par une bibliothèque ;
  \item[\empi] choix de conception réglé à l'usage. Le code ne contient ni mesure ni étude qui la justifie : le document donne l'intention et l'ordre de grandeur raisonnable, et le signale comme tel.
\end{description}

\paragraph{Notations.}
\begin{itemize}
  \item Les vecteurs sont en gras ($\vect{x}$), les matrices en majuscules ($R$).
  \item $\vect{a}\cdot\vect{b}$ désigne le produit scalaire et $\|\cdot\|$ la norme euclidienne.
  \item Les coordonnées écran sont \emph{normalisées} : $(0,0)$ est le coin haut-gauche de la fenêtre et $(1,1)$ le coin bas-droit. Une erreur de \qty{3}{\percent} signifie donc \qty{3}{\percent} de la largeur (ou de la hauteur) de la fenêtre.
  \item $\med$ désigne la médiane et $\MAD(x) = \med|x - \med(x)|$ l'écart absolu médian.
  \item Les noms entre \code{guillemets_machine} renvoient aux identifiants du code.
\end{itemize}

\chapter{Architecture et pipeline}
\label{chap:archi}

\section{Organisation des fichiers}

Le code utilisé par l'interface se trouve dans \fichier{victor/}. Le point d'entrée actuel est \fichier{victor/main_v2.py}.

\begin{table}[H]
\centering\small
\begin{tabularx}{\textwidth}{@{}l L l@{}}
\toprule
Fichier & Rôle & Lignes \\
\midrule
\fichier{main_v2.py} & Boucle principale avec le tracker v2 : navigation, sélection par fixation, recalibration au clic & 298 \\
\fichier{main.py} & Même boucle avec le tracker v1 & 278 \\
\fichier{utils_v2.py} & Tracker v2 : caméra, modèle de tête, modèle de regard géométrique, calibration 3D, recalibration & 832 \\
\fichier{utils.py} & Tracker v1 et briques partagées (indices, One Euro, features, téléchargement du modèle) & 612 \\
\fichier{interface.py} & Disposition de l'écran (géométrie des zones) et rendu pygame & 503 \\
\fichier{config.py} & Lecture, fusion et validation de la configuration & 332 \\
\fichier{actions.py} & Registre, analyse et exécution des actions ; clavier, GPIO, HTTP & 423 \\
\fichier{config.yaml} & Configuration par défaut, documentée en commentaires & 140 \\
\bottomrule
\end{tabularx}
\caption{Modules du dossier \fichier{victor/}.}
\end{table}

Les dépendances entre modules sont volontairement à sens unique : \fichier{main_v2.py} utilise tous les autres ; \fichier{utils_v2.py} réutilise \fichier{utils.py} ; \fichier{config.py} utilise \fichier{actions.py} pour valider les actions au chargement. \fichier{interface.py} ne connaît ni le tracker ni les actions : il reçoit à chaque image ce qu'il doit dessiner.

\section{Modèle d'exécution : trois threads}

\begin{figure}[H]
\centering
\begin{tikzpicture}[
  node distance=6mm and 8mm, font=\small,
  boite/.style={draw, rounded corners=3pt, align=center, minimum height=9mm, minimum width=27mm, inner sep=3pt, fill=white},
  groupe/.style={draw=doux, dashed, rounded corners=6pt, inner sep=6pt},
  fl/.style={-{Stealth[length=2.2mm]}, thick},
  lab/.style={font=\scriptsize, align=center, fill=white, inner sep=1pt}]
  % Thread tracker
  \node[boite] (cam) {Webcam\\\code{cap.read}};
  \node[boite, below=of cam] (mp) {MediaPipe\\478 landmarks};
  \node[boite, below=of mp] (mod) {Modèle de regard\\v1 ou v2};
  \node[boite, below=of mod] (oe) {Filtre One Euro};
  \draw[fl] (cam) -- (mp); \draw[fl] (mp) -- (mod); \draw[fl] (mod) -- (oe);
  % Thread principal
  \node[boite, right=24mm of cam] (boucle) {Boucle \code{main}};
  \node[boite, below=of boucle] (zone) {Zone visée\\\code{zone_at}};
  \node[boite, below=of zone] (sel) {\code{Selecteur}\\chrono de fixation};
  \node[boite, below=of sel] (rendu) {Rendu pygame};
  \draw[fl] (boucle) -- (zone); \draw[fl] (zone) -- (sel); \draw[fl] (sel) -- (rendu);
  % Thread actions
  \node[boite, right=24mm of sel] (exe) {\code{Executeur}};
  \node[boite, below=of exe] (out) {Clavier, programme,\\GPIO, HTTP};
  \draw[fl] (exe) -- (out);
  \begin{scope}[on background layer]
    \node[groupe, fit=(cam)(oe), label={[font=\small\bfseries]above:Thread \code{gaze-tracker}}] {};
    \node[groupe, fit=(boucle)(rendu), label={[font=\small\bfseries]above:Thread principal}] {};
    \node[groupe, fit=(exe)(out), label={[font=\small\bfseries]above:Thread \code{actions}}] {};
  \end{scope}
  % Échanges
  \draw[fl, accent] (oe.east) -- ++(6mm,0) |- node[lab, pos=0.75, above] {\code{gaze}} (zone.west);
  \draw[fl, accent, dashed] (boucle.west) -- node[lab, above] {\code{_cmds}} (cam.east);
  \draw[fl, accent] (sel.east) -- node[lab, above] {\code{lancer}} (exe.west);
  \draw[fl, accent, dashed] (exe.north) |- node[lab, pos=0.75, above] {\code{vers_interface}} (boucle.east);
\end{tikzpicture}
\caption{Les trois threads et leurs échanges. Traits pleins : valeur lue directement (\code{tracker.gaze}, \code{last_valid}) ou appel de méthode. Tirets : files \code{queue.Queue}.}
\label{fig:threads}
\end{figure}

\subsection{Thread du tracker}

\code{GazeTracker.start()} lance un thread démon nommé \code{gaze-tracker}, qui appelle \code{step()} en boucle jusqu'à ce que l'évènement \code{_stop} soit levé ou que la caméra cesse de répondre ; \code{alive} passe alors à \code{False}, ce qui arrête aussi la boucle principale. Le rythme est imposé par \code{cap.read()}, qui bloque jusqu'à l'image suivante. Le tracker tourne donc à la cadence de la caméra, ou moins vite si MediaPipe ne suit pas.

Un appel à \code{step()} :
\begin{enumerate}
  \item exécute toutes les commandes en attente dans \code{_cmds} (section~\ref{sec:commandes}) ;
  \item lit une image, la retourne en miroir et la convertit en RGB ;
  \item calcule un horodatage strictement croissant et met à jour la mesure de cadence (\code{fps}) ;
  \item appelle MediaPipe, extrait les mesures et teste les clignements ;
  \item pendant une calibration, transmet l'échantillon à la calibration ; sinon, calcule et filtre le point regardé ;
  \item publie \code{frame}, \code{frame_id}, \code{usable} et, si la mesure est exploitable, \code{gaze} et \code{last_valid}.
\end{enumerate}

\subsection{Thread principal}

Il exécute la boucle de \code{main()} (chapitre~\ref{chap:interaction}) à la cadence \code{images_par_seconde} (\num{30} par défaut), imposée par \code{pygame.time.Clock.tick}. Il lit \code{tracker.gaze} \emph{sans attendre} le tracker : l'interface reste fluide même si le tracker ralentit, et inversement.

\subsection{Thread des actions}

\code{actions.Executeur} crée un thread démon \code{actions} qui exécute les séquences d'actions des boutons, une à la fois, dans l'ordre d'arrivée (chapitre~\ref{chap:actions}). Une action lente (\code{attendre 1.5}, une requête HTTP qui met \qty{3}{\second} à échouer) ne fige donc jamais l'interface.

\subsection{Communication et absence de verrous}
\label{sec:commandes}

Le code n'utilise aucun verrou entre ces trois threads. La cohérence repose sur trois règles :

\begin{enumerate}
  \item \textbf{Publication par remplacement.} Le tracker publie le regard en \emph{remplaçant} l'attribut par un nouveau tableau (\code{self.gaze = np.array([...])}), sans jamais le modifier sur place. En CPython, l'affectation d'un attribut est atomique : le lecteur voit soit l'ancien tableau, soit le nouveau, jamais un mélange des deux. Même principe pour \code{last_valid} (un flottant) et \code{frame}.
  \item \textbf{Commandes sérialisées.} Toutes les méthodes publiques qui modifient l'état du tracker (\code{start_calibration}, \code{cancel_calibration}, \code{recenter}, \code{reset}, \code{add_calibration_point}, \code{save_calibration}, \code{load_calibration}) ne font que déposer une fermeture dans la file \code{_cmds}. Le thread tracker les exécute au début de l'image suivante, donc au plus une période caméra (\qty{\approx 33}{\milli\second}) plus tard. L'état interne (modèle, filtres, calibration) n'est ainsi modifié que par un seul thread.
  \item \textbf{Files dans les deux sens pour les actions.} \code{Executeur._travaux} (thread principal $\to$ actions) et \code{Executeur.vers_interface} (actions $\to$ thread principal) sont des \code{queue.Queue}, sûres entre threads.
\end{enumerate}

\begin{attention}
\code{add_calibration_point(cible, t)} reçoit l'instant \code{t} du clic \emph{au moment de l'appel}, et non à l'exécution de la commande. Les images retenues sont donc bien celles qui précèdent le clic, même si la commande s'exécute une image plus tard.
\end{attention}

\paragraph{Parallélisme réel.} Python a un verrou global (GIL), mais le travail lourd se fait dans du code natif (inférence MediaPipe en C++, OpenCV, algèbre linéaire NumPy) qui peut en général relâcher ce verrou. La séparation en threads sert surtout à \emph{découpler les cadences} : la caméra et le rendu tournent chacun à leur rythme, et aucune attente d'un côté ne bloque l'autre.

\section{Pipeline d'une image}

\begin{figure}[H]
\centering
\begin{tikzpicture}[
  node distance=4.5mm, font=\small,
  etape/.style={draw, rounded corners=3pt, align=center, minimum width=62mm, minimum height=8mm, fill=white, inner sep=3pt},
  test/.style={draw, diamond, aspect=2.6, align=center, inner sep=1pt, font=\scriptsize},
  rejet/.style={draw=orange!70!black, rounded corners=3pt, align=center, fill=orange!8, font=\scriptsize, inner sep=3pt},
  fl/.style={-{Stealth[length=2mm]}, thick}]
  \node[etape] (a) {\code{cap.read()} : image BGR $640\times480$};
  \node[etape, below=of a] (b) {\code{cv2.flip} miroir, conversion RGB, horodatage};
  \node[etape, below=of b] (c) {MediaPipe \code{detect_for_video}\\landmarks, blendshapes, matrice de pose};
  \node[etape, below=of c] (d) {\code{extract_features} : 10 features v1, score de clignement};
  \node[etape, below=of d] (e) {\code{HeadModel.observe} : pose $(R, \vect{O})$, angles $(\theta, \varphi)$};
  \node[etape, below=of e] (f) {Test de clignement (seuil \num{0.6}, maintien \qty{150}{\milli\second})};
  \node[etape, below=of f] (g) {\code{GeoGazeModel.predict} : point $(x, y)$ brut};
  \node[etape, below=of g] (h) {$+$ décalage de recentrage, filtre One Euro};
  \node[etape, below=of h, fill=accent!10] (i) {Publication : \code{gaze}, \code{last_valid}};
  \foreach \u/\v in {a/b,b/c,c/d,d/e,e/f,f/g,g/h,h/i} \draw[fl] (\u) -- (\v);
  \node[rejet, right=10mm of c] (r1) {pas de visage\\$\Rightarrow$ image ignorée};
  \node[rejet, right=10mm of e] (r2) {$|\text{yaw}|$ ou $|\text{pitch}| > \ang{50}$\\$\Rightarrow$ image ignorée};
  \node[rejet, right=10mm of f] (r3) {clignement\\$\Rightarrow$ image ignorée};
  \draw[fl, orange!70!black] (c) -- (r1); \draw[fl, orange!70!black] (e) -- (r2); \draw[fl, orange!70!black] (f) -- (r3);
  \node[etape, left=10mm of g, minimum width=30mm, fill=black!4] (cal) {Pendant la calibration :\\landmarks bruts\\$\to$ \code{Calibration3D}};
  \draw[fl, dashed] (f.west) -| (cal.north);
\end{tikzpicture}
\caption{Traitement d'une image par le tracker v2 (thread \code{gaze-tracker}). Une image ignorée ne modifie pas \code{gaze}.}
\label{fig:pipeline}
\end{figure}

Dans le thread principal, à chaque tour :
\begin{enumerate}
  \item \code{tracker.gaze} est borné à $[0, 1]^2$ puis converti en pixels ;
  \item la zone visée est déterminée sur le regard \emph{brut}, et le curseur affiché est lissé à part ;
  \item le chrono de fixation avance ou repart à zéro ; en cas de clic, l'action est envoyée au thread des actions ;
  \item seuls les rectangles modifiés de l'écran sont redessinés.
\end{enumerate}

\section{Objectifs de performance}
\label{sec:perf}

\subsection{Ce que le code fixe explicitement}

Le dépôt ne contient pas de banc de mesure. Les objectifs se lisent dans les constantes et les commentaires du code :

\begin{itemize}
  \item \textbf{Cadence caméra} : \code{CAM_FPS = 30}, résolution $640\times480$. Le commentaire de \fichier{utils.py} justifie cette résolution par son coût : elle est \enquote{bien moins coûteuse sur Raspberry Pi}. Le coût de MediaPipe croît avec la surface de l'image ; or l'iris occupe déjà une dizaine de pixels de diamètre à \qty{60}{\centi\metre} en $640\times480$, ce qui suffit au réseau.
  \item \textbf{Latence de capture minimale} : \code{CAP_PROP_BUFFERSIZE = 1}. Par défaut, le pilote V4L2 garde plusieurs images en file ; si le traitement est plus lent que la caméra, on lirait des images de plus en plus anciennes. Avec une file d'une seule image, \code{read()} renvoie toujours la plus récente disponible.
  \item \textbf{Cadence d'affichage} : \code{images_par_seconde = 30} (bornes \numrange{5}{120}).
  \item \textbf{Rendu incrémental} : les commentaires d'\fichier{interface.py} indiquent que tout est \enquote{pensé pour un Raspberry Pi} (chapitre~\ref{chap:rendu}).
  \item \textbf{Ajustement de calibration supportable} : le tracker v2 chronomètre l'ajustement et affiche sa durée dans la console (\code{ajustement x.xs}).
\end{itemize}

\subsection{Budget de latence}

Le temps entre le moment où l'œil se pose sur un bouton et le déclenchement de l'action est la somme de plusieurs termes :

\begin{equation}
T_\text{clic} \approx \underbrace{T_\text{saccade}}_{\approx\,\qty{0.2}{\second}}
 + \underbrace{T_\text{capture} + T_\text{MediaPipe} + T_\text{calcul}}_{\text{latence du tracker}}
 + \underbrace{T_\text{filtre}}_{\text{One Euro}}
 + \underbrace{T_\text{boucle}}_{\leq\,\qty{33}{\milli\second}}
 + \underbrace{T_\text{fixation}}_{\code{duree_clic} = \qty{2}{\second}} .
\end{equation}

\begin{itemize}
  \item $T_\text{saccade}$ : temps de réaction d'une saccade visuelle, environ \qtyrange{150}{250}{\milli\second}, auquel s'ajoute sa durée (quelques dizaines de millisecondes).
  \item $T_\text{capture}$ : au plus une période caméra, \qty{33}{\milli\second}, grâce à la file d'une image.
  \item $T_\text{MediaPipe}$ : non mesuré dans le dépôt. Il est visible en direct : la vignette de debug (touche D) affiche les images par seconde du tracker (moyenne exponentielle $\text{fps} \leftarrow 0.9\,\text{fps} + 0.1/\Delta t$, soit une fenêtre effective d'une dizaine d'images).
  \item $T_\text{calcul}$ : négligeable devant MediaPipe. Le tracker v2 fait trois alignements de Procrustes sur 69 points (trois SVD $3\times3$) et quelques produits matrice-vecteur.
  \item $T_\text{filtre}$ : au repos, le filtre One Euro a une constante de temps $\tau = 1/(2\pi\cdot\num{0.6}) \approx \qty{0.27}{\second}$. Pendant une saccade, la fréquence de coupure monte et ce retard diminue fortement (section~\ref{sec:oneeuro}).
  \item $T_\text{fixation}$ domine volontairement : c'est la garantie contre les clics involontaires.
\end{itemize}

Le chrono de fixation démarre dès que le regard \emph{brut} entre dans la case. Le lissage d'affichage du curseur (\qty{50}{\milli\second}) n'ajoute donc rien à $T_\text{clic}$.

\subsection{Coût de la calibration v2}

L'ajustement final est le seul calcul lourd. On peut l'estimer :
\begin{itemize}
  \item nombre d'échantillons $N \leq 16 \times 60 + 3 \times 135 = \num{1365}$ (\qty{2}{\second} et \qty{4.5}{\second} de collecte à \qty{30}{images\per\second}), moins les images perdues ;
  \item vecteur de résidus de taille $2N + 21$ (deux coordonnées par échantillon, plus 21 termes d'a priori) ;
  \item par itération de Levenberg-Marquardt, $23 + 1 = 24$ évaluations du modèle pour la jacobienne par différences finies, plus une évaluation par essai de pas ;
  \item 60 itérations au plus pour l'ajustement principal, puis $16 \times 3$ itérations pour la validation croisée.
\end{itemize}
Cela fait quelques milliers d'évaluations vectorisées sur $\sim\num{1000}$ échantillons, et la résolution de systèmes $23\times23$. L'ensemble se compte en secondes sur un PC ; la durée réelle est affichée à chaque calibration. L'ajustement s'exécute dans le thread tracker, qui ne publie plus de regard pendant ce temps, mais l'interface continue d'être dessinée.

\subsection{Objectifs proposés}

Pour transformer ces intentions en critères vérifiables, le tableau~\ref{tab:objectifs} propose des objectifs. Ce sont des \emph{propositions} : elles ne sont pas encore mesurées.

\begin{table}[H]
\centering\small
\begin{tabularx}{\textwidth}{@{}l l L@{}}
\toprule
Grandeur & Objectif & Comment la mesurer \\
\midrule
Cadence du tracker (PC) & \qty{30}{images\per\second} & champ \code{fps} affiché en mode debug \\
Cadence du tracker (Pi) & $\geq$ \qty{15}{images\per\second} & idem, sur la cible \\
Cadence de l'interface & \qty{30}{images\per\second} stables & \code{horloge.get_fps()} de pygame \\
Erreur validée (LOPO) & $<$ \qty{5}{\percent} en $x$ et $y$ & affichée en fin de calibration \\
Erreur tête en mouvement & $<$ \qty{8}{\percent} & affichée en fin de calibration (v2) \\
Durée d'ajustement & $<$ \qty{5}{\second} sur Pi & affichée en fin de calibration (v2) \\
\bottomrule
\end{tabularx}
\caption{Objectifs de performance proposés (non mesurés à ce jour).}
\label{tab:objectifs}
\end{table}

Ces seuils viennent de la taille des cibles. En grille $2\times2$ avec des bandes latérales de \qty{12}{\percent}, une case fait environ \qty{38}{\percent} de la largeur et \qty{50}{\percent} de la hauteur de l'écran. Une erreur de \qty{5}{\percent} laisse une marge confortable autour du centre d'une case, même en tenant compte de la zone morte.

\chapter{Détection du visage avec MediaPipe}
\label{chap:mediapipe}

\section{Le modèle Face Landmarker}

Les deux trackers reposent sur \emph{MediaPipe Face Landmarker}, un réseau de neurones de Google qui estime un maillage 3D du visage à partir d'une seule image RGB \cite{kartynnik2019}. Il est complété par un raffinement de la région des yeux qui localise l'iris \cite{ablavatski2020}. Le fichier \fichier{models/face_landmarker.task} contient le détecteur de visage, le maillage, l'estimateur d'iris et le modèle de \emph{blendshapes}. Il s'agit de la version \code{float16} : les poids sont stockés sur 16 bits, ce qui divise par deux la taille du fichier, sans effet notable sur la précision.

\paragraph{Téléchargement.} \code{ensure_model()} télécharge le fichier depuis le stockage public de MediaPipe s'il est absent. Il écrit d'abord dans \fichier{face_landmarker.task.part}, puis renomme ce fichier avec \code{os.replace}, qui est atomique. Une coupure pendant le téléchargement ne laisse donc jamais un modèle tronqué, qui ferait échouer tous les lancements suivants. Le chemin est calculé par rapport au fichier source (\code{os.path.dirname(os.path.dirname(__file__))}) : le modèle est trouvé quel que soit le dossier d'où l'on lance le programme.

\section{Options utilisées}

\begin{table}[H]
\centering\small
\begin{tabularx}{\textwidth}{@{}l l L@{}}
\toprule
Option & Valeur & Raison \\
\midrule
\code{running_mode} & \code{VIDEO} & Le modèle réutilise le visage trouvé à l'image précédente au lieu de relancer le détecteur à chaque image, ce qui est plus rapide. Il exige des horodatages strictement croissants. \\
\code{num_faces} & 1 & Un seul utilisateur. Le commentaire du code note que le lissage temporel interne de MediaPipe n'est actif qu'avec un seul visage. \\
\code{min_face_detection_confidence} & \num{0.5} & Valeur par défaut de MediaPipe (\conv). \\
\code{min_face_presence_confidence} & \num{0.5} & Idem : sous ce score, le visage est considéré comme perdu et le détecteur est relancé. \\
\code{min_tracking_confidence} & \num{0.5} & Idem. \\
\code{output_face_blendshapes} & \code{True} & Scores de clignement et de regard haut/bas. \\
\code{output_facial_transformation_matrixes} & \code{True} & Pose de la tête utilisée par la v1. \\
\bottomrule
\end{tabularx}
\caption{Options de \code{FaceLandmarkerOptions}.}
\end{table}

\section{Ce que renvoie MediaPipe}

\subsection{Landmarks}

Pour chaque image, \code{result.face_landmarks[0]} contient 478 points $(x, y, z)$ : 468 pour le maillage du visage et $2\times5$ pour les iris (un centre et quatre points du contour, pour chaque œil).
\begin{itemize}
  \item $x$ et $y$ sont normalisés par la largeur et la hauteur de l'image : $u = x\,w$ et $v = y\,h$ donnent les pixels.
  \item $z$ est une profondeur \emph{relative}, d'origine environ au centre de la tête et négative vers la caméra, exprimée \emph{à la même échelle que $x$}. Multiplier $z$ par $w$ la convertit donc en pixels : c'est ce que fait \code{landmarks_array} (\code{* np.array([w, h, w])}).
\end{itemize}

L'interprétation précise de $z$ compte pour la v2 (section~\ref{sec:retroprojection}). Le code suppose que $z\,w$ est un écart de profondeur $\Delta Z$ projeté avec le même facteur que les coordonnées latérales, c'est-à-dire $z\,w \approx \Delta Z \cdot f / Z_0$, où $Z_0$ est la distance de la tête et $f$ la focale en pixels. Cette hypothèse est cohérente avec la façon dont le maillage est appris (un modèle de projection faiblement perspective), mais $z$ reste une \emph{estimation} du réseau, plus bruitée que $x$ et $y$.

\subsection{Blendshapes}

\code{result.face_blendshapes[0]} donne 52 scores entre 0 et 1, nommés selon la convention ARKit. Le code en utilise six :
\begin{itemize}
  \item \code{eyeBlinkLeft}, \code{eyeBlinkRight} : fermeture des paupières ;
  \item \code{eyeLookUpLeft/Right}, \code{eyeLookDownLeft/Right} : regard vers le haut et vers le bas.
\end{itemize}
Ces scores sont appris globalement sur l'apparence de l'œil : ils captent des indices que la seule position de l'iris rend mal, comme l'abaissement de la paupière supérieure quand on regarde en bas.

\subsection{Matrice de transformation}

\code{facial_transformation_matrixes[0]} est une matrice homogène $4\times4$ qui place un visage canonique de MediaPipe dans le repère de la caméra. La v1 n'en garde que le bloc $3\times3$ de rotation, dont elle normalise les colonnes pour retirer l'échelle (fonction \code{head_pose}). Les angles d'Euler en sont extraits par :
\begin{equation}
\text{pitch} = \atantwo(r_{21}, r_{22}), \quad
\text{yaw} = \atantwo\!\left(-r_{20}, \sqrt{r_{21}^2 + r_{22}^2}\right), \quad
\text{roll} = \atantwo(r_{10}, r_{00}),
\end{equation}
avec repliement du pitch dans $[-\ang{90}, \ang{90}]$. C'est la décomposition classique \[ R = R_z(\text{roll})\,R_y(\text{yaw})\,R_x(\text{pitch}). \] La v2 n'utilise pas cette matrice : elle estime sa propre pose avec un modèle de tête propre à l'utilisateur (section~\ref{sec:pose}).

\section{Points du visage utilisés}

\begin{table}[H]
\centering\small
\begin{tabular}{@{}lcc@{}}
\toprule
Point & Œil A (gauche de l'image) & Œil B (droite de l'image) \\
\midrule
Coin extérieur (\code{outer}) & 33 & 263 \\
Coin intérieur (\code{inner}) & 133 & 362 \\
Paupière haute (\code{top}) & 159 & 386 \\
Paupière basse (\code{bottom}) & 145 & 374 \\
Centre de l'iris (\code{iris}) & 468 & 473 \\
Contour de l'iris (\code{contour}) & 469--472 & 474--477 \\
\bottomrule
\end{tabular}
\caption{Indices du maillage MediaPipe (dictionnaires \code{EYE_A}, \code{EYE_B} de \fichier{utils.py}).}
\end{table}

Ce sont les indices canoniques du maillage MediaPipe (\conv) : les paupières 159/145 et 386/374 sont les points médians des bords palpébraux. Le bout du nez (\code{NOSE_TIP = 1}) sert à la position de la tête dans la v1. Le premier prototype (\fichier{main.py} à la racine) utilisait d'autres points de paupière (27/23 et 258/253), plus éloignés de l'œil ; un commentaire \code{TODO} de ce fichier en notait déjà l'imprécision.

\paragraph{Miroir et identité des yeux.} L'image est retournée horizontalement avant l'analyse (\code{cv2.flip(image, 1)}), pour que le curseur se déplace comme dans un miroir. Ce retournement échange l'œil gauche et l'œil droit du sujet. Il n'a pourtant aucune conséquence : \code{eye_features} trie les deux coins de chaque œil par abscisse, et la v2 associe chaque iris au globe le plus proche. \enquote{A} désigne donc toujours l'œil situé à gauche de l'image traitée.

\section{Horodatage}

En mode \code{VIDEO}, MediaPipe exige des horodatages en millisecondes strictement croissants. Le tracker calcule
\[
  t_\text{ms} = \big\lfloor 1000\,(t_\text{monotone} - t_\text{démarrage}) \big\rfloor
\]
et, si deux images tombent sur la même milliseconde, impose $t_\text{ms} = t_\text{ms}^\text{précédent} + 1$. L'horloge \code{time.monotonic()} ne recule jamais, contrairement à \code{time.time()} qui peut être modifiée par une synchronisation NTP. Le premier prototype utilisait \code{time.time()}.

\chapter{Tracker v1 : régression sur des caractéristiques 2D}
\label{chap:v1}

Le tracker v1 (\fichier{victor/utils.py}, repris de \fichier{test_xav.py}) est une approche \emph{par apprentissage} : il mesure à chaque image un vecteur de caractéristiques (\emph{features}) de l'œil et de la tête, et apprend pendant la calibration une fonction \enquote{features $\mapsto$ point de l'écran}. Il est encore utilisé par \fichier{main.py}, et plusieurs de ses briques servent à la v2 : features, clignements, One Euro, repli avant calibration.

\section{Caractéristiques de l'œil}
\label{sec:features}

\subsection{Repère propre à chaque œil}

Pour un œil, on note en pixels $\vect{c}_g$ et $\vect{c}_d$ les deux coins triés par abscisse, $\vect{p}$ le centre de l'iris, et $\vect{l}_h$, $\vect{l}_b$ les milieux des paupières haute et basse. La fonction \code{eye_features} construit :
\begin{align}
  \ell &= \|\vect{c}_d - \vect{c}_g\| && \text{(largeur de l'œil)}\\
  \vect{u} &= (\vect{c}_d - \vect{c}_g)/\ell && \text{(axe des coins)}\\
  \vect{v} &= (-u_y,\ u_x) && \text{(perpendiculaire, vers le bas de l'image)}\\
  \vect{m} &= (\vect{c}_g + \vect{c}_d)/2 && \text{(milieu des coins)}
\end{align}
puis les quatre mesures, toutes \emph{sans dimension}, exprimées en largeurs d'œil :
\begin{equation}
  f_x = \frac{(\vect{p} - \vect{m})\cdot\vect{u}}{\ell}, \qquad
  f_y = \frac{(\vect{p} - \vect{m})\cdot\vect{v}}{\ell}, \qquad
  f_l = \frac{\big(\vect{p} - \tfrac{\vect{l}_h + \vect{l}_b}{2}\big)\cdot\vect{v}}{\ell}, \qquad
  a_p = \frac{\|\vect{l}_h - \vect{l}_b\|}{\ell}.
\end{equation}

Ces choix ont chacun une raison :
\begin{itemize}
  \item \textbf{Les calculs se font en pixels} (et non en coordonnées normalisées). Avec une image $640\times480$, les échelles horizontale et verticale des coordonnées normalisées diffèrent d'un facteur $4/3$ ; en pixels, le repère est orthonormé et les angles sont justes.
  \item \textbf{Le repère suit l'axe des coins} : une inclinaison de la tête (roll) fait tourner $\vect{u}$ et $\vect{v}$ avec l'œil, si bien que $f_x$ et $f_y$ n'en dépendent pas.
  \item \textbf{La normalisation se fait par la largeur} $\ell$ : la distance entre les coins ne dépend pas de la direction du regard ni des clignements, contrairement à la hauteur de l'œil. Le premier prototype divisait l'écart vertical par la hauteur paupière haute -- paupière basse ; or, quand on regarde vers le bas, la paupière supérieure descend avec l'iris et ce rapport bouge peu, puis devient instable au moindre clignement. La division par $\ell$ rend aussi les mesures indépendantes de la distance à la caméra.
  \item \textbf{Deux mesures verticales} : $f_y$ (par rapport aux coins) et $f_l$ (par rapport aux paupières) se complètent ; la régression choisit leur combinaison.
\end{itemize}
Si $\ell < 10^{-6}$ (œil dégénéré), la fonction renvoie \code{None}.

\subsection{Vecteur complet}

\code{extract_features} moyenne les mesures des deux yeux et y ajoute des grandeurs de tête :

\begin{table}[H]
\centering\small
\begin{tabularx}{\textwidth}{@{}c l L@{}}
\toprule
Indice & Nom & Définition \\
\midrule
0 & $f_x$ & iris horizontal (moyenne des deux yeux) \\
1 & $f_y$ & iris vertical par rapport aux coins \\
2 & $f_l$ & iris vertical par rapport aux paupières \\
3 & $b_v$ & $\tfrac12(\text{lookDown}_L + \text{lookDown}_R) - \tfrac12(\text{lookUp}_L + \text{lookUp}_R)$ \\
4 & $a_p$ & ouverture de l'œil \\
5 & yaw & angle de la tête (degrés), matrice MediaPipe \\
6 & pitch & idem \\
7 & $h_x$ & abscisse normalisée du bout du nez \\
8 & $h_y$ & ordonnée normalisée du bout du nez \\
9 & $h_s$ & $\|\vect{c}^A_\text{ext} - \vect{c}^B_\text{ext}\|_\text{px} / w$ : écart des coins extérieurs, $\propto 1/\text{distance}$ \\
\bottomrule
\end{tabularx}
\caption{Les 10 features de la v1 (\code{FEATURE_NAMES}).}
\end{table}

La position $(h_x, h_y)$ et la taille apparente $h_s$ complètent la pose : déplacer la tête latéralement ou s'approcher de l'écran change le point regardé même à direction de l'œil constante.

\paragraph{Score de clignement.} $s_\text{blink} = \tfrac12(\text{eyeBlinkLeft} + \text{eyeBlinkRight})$. Si les blendshapes sont absents, le repli est $s_\text{blink} = 1$ si $a_p < \num{0.08}$, et $0$ sinon. Un œil ouvert a typiquement $a_p \approx \numrange{0.25}{0.35}$ (hauteur d'environ un tiers de la largeur) ; \num{0.08} correspond à un œil presque fermé (\empi).

\section{Modèle de régression ridge}

\subsection{Matrice de design}

\code{design_matrix} ajoute trois termes quadratiques en $(f_x, f_y)$ aux 10 features :
\begin{equation}
  \vect{x} = \big(f_x, f_y, f_l, b_v, a_p, \text{yaw}, \text{pitch}, h_x, h_y, h_s,\ f_x^2,\ f_y^2,\ f_x f_y\big) \in \R^{13}.
\end{equation}
La relation entre position de l'iris et point de l'écran n'est pas linéaire : l'iris décrit une sphère, et l'écran plat est vu sous un angle variable. Un polynôme d'ordre 2 en $(f_x, f_y)$ capte l'essentiel de cette courbure. La tête n'entre que linéairement : c'est une correction au premier ordre.

\subsection{Standardisation pondérée avec planchers}

Avec des poids d'échantillon $\omega_i$ normalisés à $\sum_i \omega_i = n$ (moyenne 1), chaque colonne $j$ est centrée et réduite :
\begin{equation}
  \mu_j = \frac{\sum_i \omega_i x_{ij}}{n}, \qquad
  \sigma_j = \max\!\Big(\sqrt{\tfrac{1}{n}\textstyle\sum_i \omega_i (x_{ij} - \mu_j)^2},\ \sigma_j^\text{min}\Big), \qquad
  z_{ij} = \frac{x_{ij} - \mu_j}{\sigma_j}.
\end{equation}

\paragraph{Pourquoi un plancher $\sigma_j^\text{min}$ (\code{SD_FLOORS}) ?} La pénalité ridge porte sur les coefficients \emph{standardisés} $\beta_j$. Le coefficient dans les unités d'origine vaut $\beta_j / \sigma_j$. Si une feature a très peu varié pendant la calibration (la tête presque immobile, par exemple), $\sigma_j$ est minuscule : un $\beta_j$ modeste, donc à peine pénalisé, devient alors un coefficient énorme dans les unités d'origine. Le jour où la tête bouge vraiment, la prédiction explose. Le plancher impose de traiter une feature peu variable \emph{comme si} elle avait au moins l'écart-type $\sigma_j^\text{min}$, ce qui borne son coefficient réel.

\begin{origine}
\code{SD_FLOORS} $= (\num{0.01}, \num{0.005}, \num{0.005}, \num{0.05}, \num{0.02}, \ang{1}, \ang{1}, \num{0.005}, \num{0.005}, \num{0.003}, \num{0.002}, \num{0.0005}, \num{0.0005})$ (\empi). Ces valeurs sont de l'ordre du bruit de mesure de chaque feature. Par exemple, à \qty{60}{\centi\metre}, un œil de \qty{30}{\milli\metre} fait environ \qty{28}{px} de large ; \num{0.01} largeur d'œil représente donc environ \qty{0.3}{px}, l'ordre de grandeur du bruit de localisation de l'iris par MediaPipe. Un degré pour yaw et pitch correspond au bruit typique de la matrice de pose.
\end{origine}

\subsection{Résolution}

Il s'agit d'une régression ridge \cite{hoerl1970}. Les sorties $\vect{y}_i = (x_i, y_i)$ sont centrées par leur moyenne pondérée $\bar{\vect{y}}$. On résout, pour les deux sorties à la fois ($B \in \R^{13\times2}$) :
\begin{equation}
  \min_{B}\ \sum_i \omega_i \big\|\vect{y}_i - \bar{\vect{y}} - \vect{z}_i\T B\big\|^2 + \lambda\, n \sum_j P_j \|\vect{\beta}_j\|^2,
\end{equation}
dont la solution est
\begin{equation}
  B = \big(Z\T \Omega Z + \lambda n \diag(P)\big)^{-1} Z\T \Omega\, (Y - \bar{\vect{y}}),
  \qquad \hat{\vect{y}}(\vect{x}) = \bar{\vect{y}} + \vect{z}(\vect{x})\T B .
\end{equation}

\begin{itemize}
  \item L'ordonnée à l'origine n'est pas pénalisée : elle est absorbée par le centrage ($\bar{\vect{y}}$).
  \item Le facteur $n$ devant $\lambda$ rend la régularisation \enquote{par échantillon}. Le terme de données croît comme $n$ ; sans ce facteur, le même $\lambda$ régulariserait beaucoup plus une calibration courte qu'une longue.
  \item $P_j = 1$ pour les features linéaires et $P_j = 4$ pour les termes quadratiques (\code{PENALTY}, \empi) : les courbures ne sont acceptées que si les données les justifient nettement. Comme la pénalité est quadratique, un facteur 4 réduit leur amplitude typique d'un facteur 2 environ.
\end{itemize}

\subsection{Écrêtage des features de tête}

À la prédiction, les colonnes standardisées de la tête (yaw, pitch, $h_x$, $h_y$, $h_s$) sont bornées à $\pm3$ (\code{HEAD_CLIP_Z}, \deriv) : une pose de tête jamais rencontrée pendant la calibration ne peut pas être extrapolée au-delà de 3 écarts-types appris. Pour une distribution normale, $\pm3\sigma$ couvre \qty{99.7}{\percent} des valeurs ; l'écrêtage ne touche donc que des poses réellement nouvelles.

\section{Calibration v1}

\subsection{Protocole}

\code{Calibration} est une machine à états qui présente une cible à la fois.

\begin{table}[H]
\centering\small
\begin{tabularx}{\textwidth}{@{}l l l L@{}}
\toprule
Constante & Valeur & Origine & Justification \\
\midrule
\code{CAL_COLS}, \code{CAL_ROWS} & $4\times4$ & \empi & 16 points : assez pour un polynôme d'ordre 2 en 2D (6 coefficients par sortie) avec une large redondance, en restant sous la minute. \\
\code{CAL_MARGIN} & \num{0.05} & \empi & Les cibles extrêmes sont à \qty{5}{\percent} des bords. Le commentaire indique que cela couvre les bandes latérales (alors de \qty{7}{\percent}), qu'il faut savoir viser. \\
\code{CAL_LEAD_S} & \qty{2.5}{\second} & \empi & Temps de lecture de la consigne avant la première cible. \\
\code{CAL_SETTLE_S} & \qty{0.8}{\second} & \empi / \litt & Temps laissé à l'œil pour atteindre la cible : latence d'une saccade (\qtyrange{150}{250}{\milli\second}), saccade correctrice éventuelle, retard du filtrage interne de MediaPipe. \qty{0.8}{\second} est une marge confortable. \\
\code{CAL_COLLECT_S} & \qty{2.0}{\second} & \empi & Environ 60 images à \qty{30}{images\per\second} : assez pour moyenner le bruit et laisser la tête bouger un peu. \\
\code{CAL_MIN_SAMPLES} & 15 & \empi & \qty{25}{\percent} des images possibles : tolère clignements et pertes du visage. \\
\code{CAL_MIN_POINTS} & 10 & \empi & Au moins 10 cibles sur 16, pour garder une couverture spatiale suffisante. \\
Reprises & 2 & \empi & Une cible ratée est recommencée deux fois, puis ignorée (\code{skipped}). \\
\bottomrule
\end{tabularx}
\caption{Paramètres du protocole de calibration v1.}
\end{table}

Les 16 cibles sont placées en
\begin{equation}
  (x_c, y_r) = \Big(m + (1-2m)\tfrac{c}{3},\ m + (1-2m)\tfrac{r}{3}\Big), \qquad c, r \in \{0,1,2,3\},\ m = \num{0.05},
\end{equation}
puis \emph{mélangées aléatoirement} (\code{random.shuffle}). Avec un ordre fixe, l'utilisateur anticipe la cible suivante et y porte le regard avant qu'elle n'apparaisse, ce qui fausse les premiers échantillons.

Durée totale : $\num{2.5} + 16 \times (\num{0.8} + \num{2.0}) = \qty{47.3}{\second}$ sans reprise.

La consigne de la v1 est de \emph{bouger doucement la tête} (petits cercles de quelques centimètres) tout en fixant la cible. C'est indispensable : sans variation de la tête pendant la calibration, la régression ne peut pas apprendre comment les features de tête modifient le point regardé.

\subsection{Rejet des aberrants par cible}

À la fin de chaque cible (\code{reject_outliers}), on écarte les échantillons dont $f_x$ ou $f_y$ s'éloigne de la médiane de plus de 5 écarts-types robustes :
\begin{equation}
  \hat\sigma = \max\big(\num{1.4826}\,\MAD,\ \num{0.01}\big), \qquad \text{garder } i \iff |f_i - \med(f)| \leq 5\,\hat\sigma .
\end{equation}

\paragraph{D'où vient \num{1.4826} ?} Pour une variable normale $X \sim \mathcal N(\mu, \sigma^2)$, la médiane de $|X - \mu|$ vaut $\sigma\,\Phi^{-1}(3/4) = \num{0.67449}\,\sigma$, où $\Phi$ est la fonction de répartition de la loi normale centrée réduite. Donc $\hat\sigma = \MAD / \num{0.67449} = \num{1.4826}\,\MAD$ est un estimateur \emph{cohérent} de $\sigma$ pour des données normales, mais insensible aux valeurs aberrantes : il supporte jusqu'à \qty{50}{\percent} de contamination \cite{rousseeuw1993} (\deriv).

Le seuil de 5 est volontairement large (\empi) : la tête bouge pendant la collecte, si bien que $f_x$ et $f_y$ varient légitimement au sein d'une même cible. Seuls les échantillons franchement aberrants (clignement partiel, saccade) sont retirés ici ; le tri fin vient ensuite, sur les résidus du modèle. Le plancher \num{0.01} évite de tout rejeter quand l'œil est très stable et que le MAD tend vers 0.

\subsection{Ajustement}

\begin{algorithm}[H]
\caption{Ajustement de la calibration v1 (\code{Calibration.fit})}
\KwIn{échantillons $\{(\vect{x}_i, \vect{y}_i, g_i)\}$, $g_i$ = indice de la cible}
\If{moins de 10 cibles validées}{\Return échec}
\ForEach{$\lambda \in \{\num{0.002}, \num{0.01}, \num{0.05}, \num{0.2}, \num{1}\}$}{
  $e(\lambda) \leftarrow$ erreur LOPO, somme des erreurs moyennes en $x$ et $y$\;
}
$\lambda^\star \leftarrow \argmin_\lambda e(\lambda)$\;
$\omega_i \leftarrow 1/|\{k : g_k = g_i\}|$ \tcp*{chaque cible pèse autant}
ajuster le modèle ridge $(\lambda^\star, \omega)$\;
$r_i \leftarrow \|\hat{\vect{y}}_i - \vect{y}_i\|$\;
garder $i$ si $r_i \leq \med(r) + 3{,}5 \max(\num{1.4826}\MAD(r), \num{0.01})$\;
\If{au moins \qty{70}{\percent} sont gardés et au moins un est rejeté}{réajuster sur les échantillons gardés}
calculer l'erreur d'entraînement et l'erreur LOPO finales\;
\end{algorithm}

\paragraph{Validation croisée \enquote{une cible laissée de côté} (LOPO).} Pour chaque cible $g$, on entraîne sur toutes les autres cibles et on mesure l'erreur absolue moyenne sur les échantillons de $g$ ; l'erreur LOPO est la moyenne sur les cibles \cite{stone1974}. On laisse de côté une cible \emph{entière}, et non un échantillon : les images d'une même cible sont fortement corrélées (même point, images consécutives), et en retirer une seule sous-estimerait grossièrement l'erreur réelle. Une cible n'est évaluée que s'il reste au moins 20 échantillons pour l'entraînement.

\paragraph{Grille de $\lambda$.} Les cinq valeurs couvrent près de trois ordres de grandeur avec un pas d'environ $\times5$ (\empi). Cela suffit pour une régression ridge : l'erreur varie lentement avec $\log\lambda$.

\paragraph{Poids par cible.} Une cible qui a perdu des images (clignements) ne doit pas compter moins qu'une autre. Le poids $1/n_g$ donne à chaque cible un poids total de 1.

\paragraph{Second rejet.} Le seuil \num{3.5} sur les résidus est celui recommandé par Iglewicz et Hoaglin pour le \enquote{z-score modifié} \cite{iglewicz1993} (\litt). La condition \qty{70}{\percent} (\empi) empêche une calibration ratée de se \enquote{réparer} en jetant la majorité des données.

\subsection{Diagnostics}

À la fin, le statut affiche l'erreur LOPO et l'erreur d'entraînement en pourcentage de la fenêtre, le $\lambda$ choisi et l'écart-type moyen de la tête au sein d'une cible (yaw, pitch). Deux avertissements sont prévus :
\begin{itemize}
  \item tête quasi immobile : $\max(\text{sd}_\text{yaw}, \text{sd}_\text{pitch}) < \ang{0.8}$ (\code{CAL_HEAD_MIN_SD}, \empi), la tolérance aux mouvements de tête sera faible ;
  \item erreur LOPO $> \qty{12}{\percent}$ (\empi) : environ un tiers de la largeur d'une case en grille $2\times2$, l'utilisation devient pénible.
\end{itemize}

\section{Garde-fous et repli}

\begin{itemize}
  \item \textbf{Pose trop forte} : $|\text{yaw}| > \ang{40}$ ou $|\text{pitch}| > \ang{35}$ (\code{MAX_YAW_DEG}, \code{MAX_PITCH_DEG}, \empi). Au-delà, l'œil le plus éloigné est fortement raccourci par la perspective, voire masqué par le nez, et la localisation de l'iris n'est plus fiable. Le pitch a une limite plus basse parce que les paupières masquent vite l'iris quand la tête se baisse ou se lève.
  \item \textbf{Avant calibration} : \code{fallback_gaze} donne $x = \num{0.5} + 3 f_x$ et $y = \num{0.5} + 6 f_y$ (\code{FALLBACK_GAIN_X}, \code{FALLBACK_GAIN_Y}, \empi). Quand le regard balaie l'écran, $f_x$ varie d'environ $\pm\num{0.15}$ et $f_y$ d'environ moitié moins, d'où le gain vertical double. Ce mapping n'a aucune prétention de précision : il permet seulement au curseur de bouger dans le bon sens pendant l'installation.
\end{itemize}

\section{Limite de la v1}

Le modèle traite la tête comme une correction \emph{linéaire} apprise. Pour bien l'apprendre, il faudrait que la tête explore pendant la calibration toutes les poses de l'utilisation, ce qui est irréaliste. Quand la tête a peu bougé, la régularisation écrase les coefficients de tête vers 0, et le curseur dérive dès que la tête se déplace en utilisation. C'est ce constat qui a motivé le modèle géométrique de la v2 : la compensation de la tête y est \emph{calculée}, pas apprise.

\chapter{Tracker v2 : géométrie 3D de la tête et des yeux}
\label{chap:v2geo}

\section{Principe}

La v2 (\fichier{victor/utils_v2.py}) part de l'observation qui clôt le chapitre précédent : la v1 essaie d'\emph{apprendre} l'effet de la tête, alors que cet effet est purement géométrique. Si l'on connaît la position et l'orientation de la tête, et la direction des yeux par rapport à la tête, le point regardé se \emph{calcule} : c'est l'intersection d'un rayon 3D avec le plan de l'écran. La calibration n'a plus qu'à estimer quelques paramètres physiques propres à l'utilisateur et à l'installation (décalage de l'axe visuel, gains, position de l'écran), qui ne dépendent pas de la pose.

Le chapitre décrit la partie géométrique : reconstruction 3D, pose de la tête, direction des yeux. Le chapitre~\ref{chap:v2mod} décrit le passage au point de l'écran et l'ajustement des paramètres.

\begin{figure}[H]
\centering
\begin{tikzpicture}[scale=1.05, font=\small, >={Stealth[length=2.2mm]}]
  % plan écran / caméra
  \draw[very thick, doux] (0,-2.6) -- (0,2.2) node[above, black] {plan de l'écran};
  \fill[black] (-0.25,-0.18) rectangle (0.05,0.18);
  \node[left] at (-0.25,0) {caméra};
  % tête
  \draw[rounded corners=18pt, black!50] (6.2,-2.3) -- (5.2,-2.3) -- (4.6,-0.9) -- (4.75,1.4) -- (6.0,2.2) -- (8.6,2.2) -- (9.4,0.6) -- (9.2,-1.8) -- (8.4,-2.4) -- cycle;
  \node[black!60] at (8.1,-1.8) {tête};
  % globe
  \coordinate (E) at (5.9,0.45);
  \draw[thick] (E) circle (0.75);
  \fill (E) circle (1.5pt) node[right=2pt] {$\vect{e}$};
  % iris
  \coordinate (I) at ($(E)+(172:0.75)$);
  \fill[accent] (I) circle (2.4pt);
  \node[accent, below right=1pt and -2pt] at (I) {\scriptsize iris};
  % rayon caméra
  \draw[dashed, ->] (0.05,0) -- ($(0.05,0)!1.12!(I)$) node[pos=0.55, below, sloped] {rayon caméra $\lambda\vect{d}$};
  % axe optique prolongé vers l'écran
  \draw[->, thick, accent] (E) -- ($(E)!7.9!(I)$) coordinate (P);
  \fill[accent] (P) circle (2pt) node[left, accent] {point regardé};
  \node[accent, above, sloped] at ($(I)!0.45!(P)$) {regard $\vect{g}$};
  % rayon
  \draw[<->] (E) -- ($(E)+(-30:0.75)$) node[midway, below right=-2pt] {\scriptsize $r$};
  % profondeur
  \draw[|<->|, doux] (5.9,-0.65) -- (0.05,-0.65) node[midway, below, black] {$\approx Z_0$ (\qtyrange{40}{80}{\centi\metre})};
\end{tikzpicture}
\caption{Vue de profil du modèle. Le rayon caméra passant par le pixel de l'iris est intersecté avec la sphère du globe ; le vecteur centre $\to$ iris donne l'axe de l'œil, prolongé jusqu'au plan de l'écran, qui contient la caméra.}
\label{fig:geometrie}
\end{figure}

\section{Repères}

\begin{description}
  \item[Repère caméra] (cm) : origine au centre optique, $x$ vers la droite de l'image, $y$ vers le bas, $z$ le long de l'axe optique, vers la scène (donc vers l'utilisateur).
  \item[Repère tête] (cm) : $x$ de l'œil A vers l'œil B, $y$ vers le bas du visage, $z$ vers l'arrière du crâne. Un regard droit devant, vers la caméra, a la direction $(0, 0, -1)$ dans ce repère.
  \item[Pose] : une rotation $R$ (repère tête $\to$ repère caméra) et une translation $\vect{T}$ en cm, telles que $\vect{X}_\text{cam} = R\,\vect{X}_\text{tête} + \vect{T}$.
\end{description}

\section{Caméra sténopé}

\code{Camera} modélise la webcam par un sténopé de focale $f$ (en pixels), sans distorsion, avec le centre optique au milieu de l'image, $\vect{c} = (w/2, h/2)$. Les webcams grand public ont une distorsion radiale faible au centre de l'image, là où se trouve le visage ; la négliger est une approximation raisonnable.

\paragraph{Focale.} Sans calibration de la caméra, la focale se déduit d'un champ de vision horizontal supposé :
\begin{equation}
  f = \frac{w/2}{\tan(\text{HFOV}/2)}, \qquad \text{HFOV} = \ang{60} \;\Rightarrow\; f = \frac{320}{\tan \ang{30}} \approx \qty{554}{px} \text{ pour } w = 640 .
\end{equation}

\begin{origine}
\code{CAM_HFOV_DEG = 60} (\litt{} / \empi). Les webcams d'ordinateurs portables et les webcams USB courantes ont un champ horizontal de l'ordre de \ang{55} à \ang{70}. Une erreur sur cette valeur fausse surtout l'estimation de la \emph{distance} de la tête ; c'est pourquoi le modèle de regard a un paramètre de correction de focale, $k_f$ (section~\ref{sec:focale}). Si la résolution réelle de la caméra change en cours de route, la focale est remise à l'échelle proportionnellement à la largeur (\code{self.cam.f * w / self.cam.w}).
\end{origine}

\paragraph{Rayon d'un pixel.} Le rayon issu du centre optique qui passe par le pixel $(u, v)$ a pour direction, non normalisée,
\begin{equation}
  \vect{d}(u, v) = \Big(\frac{u - c_x}{f},\ \frac{v - c_y}{f},\ 1\Big).
\end{equation}

\section{Rétroprojection des landmarks}
\label{sec:retroprojection}

\code{Camera.backproject} reconstruit en 3D les points rigides du visage \emph{à un facteur d'échelle près}.

Soit un point du visage à la profondeur $Z = Z_0 + \Delta Z$, où $Z_0$ est la profondeur de référence de la tête. La projection perspective donne exactement $u = c_x + f X/Z$ et $v = c_y + f Y/Z$, donc
\begin{equation}
  \vect{X} = Z\,\vect{d}(u, v).
\end{equation}
MediaPipe fournit, après conversion en pixels, $z_\text{px} \approx f\,\Delta Z / Z_0$ (chapitre~\ref{chap:mediapipe}), d'où $Z = Z_0 (1 + z_\text{px}/f)$ et
\begin{equation}
  \boxed{\ \vect{X} = Z_0\,\Big(1 + \frac{z_\text{px}}{f}\Big)\,\vect{d}(u, v)\ }
  \qquad\Longrightarrow\qquad
  \vect{q} \coloneqq \frac{\vect{X}}{Z_0} = \Big(1 + \frac{z_\text{px}}{f}\Big)\vect{d}(u, v).
\end{equation}

Le code calcule $\vect{q}$, qui ne dépend pas de l'inconnue $Z_0$. Tous les points sont simplement divisés par la même distance, et ce facteur commun sera retrouvé comme facteur d'échelle de l'alignement de Procrustes. Ce calcul n'utilise \emph{aucune} approximation orthographique : une tête décalée sur le côté de l'image, ou proche de la caméra, est traitée avec la perspective exacte (aux erreurs de $z$ près).

\section{Alignement de Procrustes (méthode d'Umeyama)}
\label{sec:umeyama}

\subsection{Problème}

\code{umeyama(x, y, w)} cherche la similarité (échelle $s > 0$, rotation $R$, translation $\vect{t}$) qui envoie au mieux un nuage $\{\vect{x}_i\}$ sur un nuage $\{\vect{y}_i\}$ mis en correspondance, au sens des moindres carrés pondérés :
\begin{equation}
  \min_{s,\,R \in SO(3),\,\vect{t}} \ \sum_i w_i \,\big\|\vect{y}_i - (s R \vect{x}_i + \vect{t})\big\|^2 .
\end{equation}
Le problème est non linéaire à cause de la contrainte $R \in SO(3)$, mais il a une solution exacte, sans itération \cite{umeyama1991}.

\subsection{Solution}

Avec les poids normalisés $\tilde w_i = w_i / \sum_k w_k$ :
\begin{enumerate}
  \item centres pondérés : $\bar{\vect{x}} = \sum \tilde w_i \vect{x}_i$, $\bar{\vect{y}} = \sum \tilde w_i \vect{y}_i$ ; nuages centrés $\vect{x}'_i$, $\vect{y}'_i$ ;
  \item matrice de covariance croisée : $\Sigma = \sum_i \tilde w_i\, \vect{y}'_i \vect{x}'^{\top}_i \in \R^{3\times3}$ ;
  \item décomposition en valeurs singulières : $\Sigma = U D V\T$ ;
  \item correction de signe : $S = \diag(1, 1, \sign(\det U \det V))$ ;
  \item $R = U S V\T$, \quad $s = \dfrac{\tr(D S)}{\sum_i \tilde w_i \|\vect{x}'_i\|^2}$, \quad $\vect{t} = \bar{\vect{y}} - s R \bar{\vect{x}}$.
\end{enumerate}

\paragraph{Justification.} À $s$ et $R$ fixés, la translation optimale aligne les centres, d'où $\vect{t}$. Il reste à maximiser $\tr(R\T \Sigma)$ sur les rotations. En écrivant $\Sigma = U D V\T$, on a $\tr(R\T U D V\T) = \tr(D\,M)$ avec $M = U\T R V$ orthogonale. Ce produit est maximal pour $M = I$, soit $R = U V\T$. Si cette matrice est une réflexion ($\det = -1$), la meilleure rotation inverse l'axe associé à la plus petite valeur singulière : c'est le rôle de $S$. L'échelle optimale s'obtient ensuite en annulant la dérivée par rapport à $s$, qui est un problème quadratique en une variable.

\paragraph{Quelques mots sur la SVD.} Toute matrice réelle $\Sigma$ s'écrit $U D V\T$, avec $U$ et $V$ orthogonales et $D$ diagonale positive. Géométriquement, $\Sigma$ est une rotation ($V\T$), suivie d'étirements le long des axes ($D$), suivie d'une autre rotation ($U$). Pour une matrice $3\times3$, le calcul est quasi instantané. La SVD sert deux fois dans la v2 : ici, pour l'alignement, et pour trouver la normale au visage (section~\ref{sec:repere-visage}).

\section{Modèle de tête propre à l'utilisateur}
\label{sec:headmodel}

\code{HeadModel.build} construit, une seule fois à la fin de la calibration, la forme 3D rigide du visage de l'utilisateur en centimètres, ainsi que la position de ses deux globes oculaires.

\subsection{Points rigides}

Seuls 69 points du maillage (\code{RIGID_IDX}) sont utilisés : nez (15 points), front (14), coins des yeux et orbites (12), pommettes (14), tempes et côtés du visage (14). La bouche, la mâchoire, les sourcils et les paupières sont exclus : ils bougent avec la parole et les expressions et fausseraient l'estimation de la pose. Ces points ont été choisis à la main parmi les régions osseuses du visage (\empi).

\subsection{Procrustes généralisé robuste}

\begin{algorithm}[H]
\caption{Construction de la forme moyenne (\code{HeadModel.build})}
\KwIn{landmarks de $F$ images de la calibration}
\If{$F > 150$}{garder 150 images régulièrement espacées dans le temps}
rétroprojeter : nuages $Q_1, \dots, Q_F$ (69 points chacun)\;
$M \leftarrow Q_1$ centré et normalisé (RMS $= 1$)\;
\For{4 itérations}{
  \For{$k = 1..F$}{$(s_k, R_k, \vect{t}_k) \leftarrow$ \code{umeyama}$(Q_k, M)$ ; \quad $\tilde Q_k \leftarrow s_k Q_k R_k\T + \vect{t}_k$}
  $M \leftarrow$ médiane point par point des $\tilde Q_k$\;
  $M \leftarrow$ centré, puis normalisé à une distance RMS au centre de 1\;
}
\Return{$M$}
\end{algorithm}

C'est l'analyse de Procrustes généralisée \cite{gower1975}, avec deux adaptations :
\begin{itemize}
  \item \textbf{Médiane au lieu de moyenne.} Une image où MediaPipe a mal placé le visage (début de clignement, flou de mouvement) ne déforme pas la forme moyenne : la médiane résiste à une minorité d'images aberrantes.
  \item \textbf{Normalisation à chaque itération.} Sans elle, la procédure peut dériver vers la solution triviale $M \to 0$ (tous les nuages réduits à un point, erreur nulle). Fixer la taille RMS à 1 rend la forme bien définie à une similarité près.
\end{itemize}

\begin{origine}
\code{MODEL_MAX_FRAMES = 150} (\empi) : limite le coût ($150 \times 4$ alignements) tout en gardant des images de toute la calibration, y compris de la phase tête en mouvement. Le sous-échantillonnage est régulier (\code{np.linspace}), et non aléatoire, pour couvrir uniformément toutes les poses. \textbf{4 itérations} (\empi) : en pratique, l'analyse de Procrustes généralisée converge en deux ou trois itérations.
\end{origine}

\subsection{Échelle métrique}

La forme $M$ n'a pas d'unité. On la met en centimètres en imposant la distance entre les yeux. On note $\vect{m}_A$ et $\vect{m}_B$ les milieux des coins de chaque œil :
\begin{equation}
  k = \frac{\text{IPD}}{\|\vect{m}_B - \vect{m}_A\|}, \qquad M_\text{cm} = k\,M, \qquad \text{IPD} = \qty{6.3}{\centi\metre}.
\end{equation}

\begin{origine}
\code{IPD_CM = 6.3} (\litt). La distance interpupillaire moyenne d'un adulte est d'environ \qty{63}{\milli\metre} : \qty{\approx 64.7}{\milli\metre} chez les hommes, \qty{\approx 62.3}{\milli\metre} chez les femmes, avec un écart-type d'environ \qty{3.5}{\milli\metre} \cite{dodgson2004}. Le milieu des coins de l'œil est proche du centre de la pupille en regard droit devant ; l'approximation $\|\vect{m}_B - \vect{m}_A\| \approx \text{IPD}$ est donc raisonnable. Une IPD réelle différente de \qty{5}{\percent} fausse toutes les longueurs de \qty{5}{\percent}, distance comprise. Cette erreur est absorbée par la correction de distance $k_f$ et par l'application affine vers l'écran (chapitre~\ref{chap:v2mod}).
\end{origine}

\subsection{Repère du visage}
\label{sec:repere-visage}

\begin{enumerate}
  \item \textbf{Axe $z$} : on fait la SVD du nuage centré $M_\text{cm} - \bar M$. Le troisième vecteur singulier droit est la direction de \emph{moindre variance}, c'est-à-dire la normale au plan qui approche le mieux le visage. Il est orienté pour avoir une composante $z$ positive. La forme a été construite dans des coordonnées issues de la caméra, où $+z$ s'éloigne de la caméra : $z$ pointe donc vers l'arrière du crâne.
  \item \textbf{Axe $x$} : direction $\vect{m}_B - \vect{m}_A$, rendue orthogonale à $z$ par Gram-Schmidt, puis normalisée.
  \item \textbf{Axe $y$} $= z \times x$ : pour $z$ vers l'arrière et $x$ vers la droite de l'image, $y$ pointe vers le bas.
\end{enumerate}
Les points et les milieux des yeux sont exprimés dans cette base, $P = M_\text{cm}\,[\vect{x}\ \vect{y}\ \vect{z}]$.

\subsection{Centres des globes oculaires}

\begin{equation}
  \vect{e}_A = \vect{m}_A + (0, 0, \delta), \qquad \vect{e}_B = \vect{m}_B + (0, 0, \delta), \qquad \delta = \qty{1.0}{\centi\metre}.
\end{equation}

\begin{origine}
\code{EYE_DEPTH_CM = 1.0} et \code{EYE_RADIUS_CM = 1.15} (\litt{} pour l'ordre de grandeur, \empi{} pour la valeur exacte). Le globe oculaire humain a un diamètre axial d'environ \qty{24}{\milli\metre}, et son centre de rotation se situe environ \qtyrange{13}{14}{\milli\metre} derrière le sommet de la cornée. La pupille apparente, celle que l'on voit à travers la cornée, se trouve à quelques millimètres derrière ce sommet : la distance centre de rotation -- pupille apparente vaut donc de l'ordre de \qtyrange{10}{12}{\milli\metre}. Les modèles géométriques classiques de l'œil donnent des valeurs du même ordre \cite{guestrin2006,hansen2010}. De même, le plan des coins de l'œil passe à peu près au niveau de l'iris : placer le centre du globe \qty{1}{\centi\metre} derrière ce plan est cohérent. Ces deux valeurs restent \emph{approximatives}. Une erreur sur le rayon $r$ change l'angle déduit d'un déplacement de l'iris d'un facteur à peu près constant ($\Delta\text{iris} \approx r \sin\theta$), et elle est absorbée par les \emph{gains} angulaires appris à la calibration (section~\ref{sec:correction-angles}).
\end{origine}

\section{Pose de la tête à chaque image}
\label{sec:pose}

\code{HeadModel.pose} estime à chaque image la pose de la tête à partir des 69 points rétroprojetés $\vect{q}_k$ (section~\ref{sec:retroprojection}) et du modèle $P_k$ (en cm).

\subsection{Retrouver la distance}

On aligne le modèle sur l'observation, $\vect{q}_k \approx s R P_k + \vect{t}$. La vraie pose vérifie $\vect{X}_k = R P_k + \vect{T}$ et $\vect{q}_k = \vect{X}_k / Z_0$, d'où par identification
\begin{equation}
  s = \frac{1}{Z_0}, \qquad \vect{t} = \frac{\vect{T}}{Z_0} \quad\Longrightarrow\quad \vect{T} = \frac{\vect{t}}{s} \ \text{(en cm)}.
\end{equation}
L'échelle de Procrustes donne donc directement la distance de la tête, en centimètres grâce à l'IPD du modèle. \code{pose} renvoie $(R, \vect{t}/s, s)$, et la mesure est rejetée si $s \leq 0$ ou si $R$ contient des valeurs non finies.

\subsection{Moindres carrés repondérés}

Un landmark mal placé (cheveux, main devant le visage, reflet de lunettes) ne doit pas fausser la pose. Après un premier alignement non pondéré, deux passes de moindres carrés itérativement repondérés (IRLS \cite{holland1977}) attribuent à chaque point le poids de Cauchy :
\begin{equation}
  \rho_k = \big\|\vect{q}_k - (sRP_k + \vect{t})\big\|, \qquad
  \hat\sigma = \num{1.4826}\,\med_k(\rho_k) + 10^{-9}, \qquad
  w_k = \frac{1}{1 + \big(\rho_k / (\num{2.5}\,\hat\sigma)\big)^2} .
\end{equation}

\begin{origine}
\textbf{Poids de Cauchy} (\litt) : c'est le poids de l'estimateur M de Cauchy, $\rho(r) = \tfrac{c^2}{2}\log(1 + (r/c)^2)$. Il réduit continûment l'influence des grands résidus sans jamais l'annuler tout à fait, ce qui le rend stable même si l'alignement initial est imparfait. La constante classique $c = \num{2.385}\,\sigma$ donne une efficacité de \qty{95}{\percent} sous bruit gaussien \cite{holland1977} ; le code prend $c = \num{2.5}\,\hat\sigma$, très proche (\empi). L'échelle $\hat\sigma$ est la médiane des \emph{normes} de résidus multipliée par la constante de cohérence \num{1.4826}. Le terme $10^{-9}$ évite une division par zéro quand l'ajustement est parfait. \textbf{Deux passes} (\empi) suffisent : le premier alignement est déjà bon, et les poids se stabilisent aussitôt.
\end{origine}

\subsection{Angles d'Euler et garde-fou}

\code{euler_deg(R)} convertit la rotation en angles lisibles. La direction du visage est $\vect{a} = -R_{:,3}$ (opposée à l'axe $z$ de la tête, qui pointe vers l'arrière) :
\begin{equation}
  \text{yaw} = \atantwo(a_x, -a_z), \qquad \text{pitch} = \atantwo\!\big(a_y, \sqrt{a_x^2 + a_z^2}\big), \qquad \text{roll} = \atantwo(R_{10}, R_{00}).
\end{equation}
Un yaw positif signifie un visage tourné vers la droite de l'image, un pitch positif une tête baissée. Ces angles servent à l'affichage de debug, à l'attribut \code{head_pose} (avec la distance $O_z$) et au garde-fou : si $|\text{yaw}|$ ou $|\text{pitch}|$ dépasse \ang{50} (\code{MAX_HEAD_ANGLE_DEG}, \empi), la mesure est jetée. La v2 tolère des angles plus grands que la v1 (\ang{40}/\ang{35}) parce que le modèle 3D gère correctement la perspective ; au-delà de \ang{50}, un œil est en pratique masqué par le nez.

\section{Direction des yeux par rapport à la tête}
\label{sec:observe}

\code{HeadModel.observe} est le cœur géométrique de la v2. Pour chaque image, elle renvoie l'origine du regard et sa direction \emph{dans le repère de la tête}.

\subsection{Centres des globes dans le repère caméra}

\begin{equation}
  \vect{e}^\text{cam}_j = \vect{T} + R\,\vect{e}_j, \qquad j \in \{A, B\}.
\end{equation}
Par sécurité, on projette ces centres dans l'image, $\pi(\vect{e}) = f\,(e_x/e_z, e_y/e_z) + \vect{c}$, et on associe chaque iris détecté au centre projeté le plus proche. Les deux iris sont échangés si nécessaire.

\subsection{Intersection rayon -- sphère}

L'iris est un point de la surface du globe : il se trouve sur le rayon caméra $\lambda\,\vect{d}$ de son pixel \emph{et} sur la sphère de centre $\vect{e}$ et de rayon $r$ :
\begin{equation}
  \|\lambda\vect{d} - \vect{e}\|^2 = r^2
  \iff a\lambda^2 - 2b\lambda + c = 0, \qquad a = \vect{d}\cdot\vect{d},\ \ b = \vect{d}\cdot\vect{e},\ \ c = \vect{e}\cdot\vect{e} - r^2 .
\end{equation}
Si le discriminant réduit $\Delta = b^2 - ac$ est positif, la racine la plus petite,
\begin{equation}
  \lambda = \frac{b - \sqrt{\Delta}}{a},
\end{equation}
correspond au point de la sphère le plus proche de la caméra, donc à la face visible du globe. Si $\Delta \leq 0$, le rayon manque la sphère (bruit, ou modèle imparfait) ; on prend alors le point du rayon le plus proche du centre, $\lambda = b/a$, qui minimise $\|\lambda\vect{d} - \vect{e}\|$.

\subsection{Axe optique et angles œil/tête}

L'axe de l'œil est le vecteur unitaire du centre du globe vers l'iris, ramené dans le repère tête :
\begin{equation}
  \vect{g}_j = R\T\,\frac{\lambda\vect{d} - \vect{e}^\text{cam}_j}{\|\lambda\vect{d} - \vect{e}^\text{cam}_j\|} .
\end{equation}
Les deux yeux regardent le même point et sont à peu près parallèles à distance d'écran. On moyenne donc leurs axes, $\vect{v} = (\vect{g}_A + \vect{g}_B)/\|\vect{g}_A + \vect{g}_B\|$, ce qui réduit le bruit d'un facteur $\approx\sqrt2$. On en tire deux angles :
\begin{equation}
  \boxed{\ \theta = \atantwo(v_x,\ -v_z), \qquad \varphi = \atantwo\!\Big(v_y,\ \sqrt{v_x^2 + v_z^2}\Big)\ }
\end{equation}
$\theta$ est l'angle horizontal (positif vers l'œil B), $\varphi$ l'angle vertical (positif vers le bas), et le regard droit devant ($\vect{v} = (0,0,-1)$) donne $\theta = \varphi = 0$.

\paragraph{Invariance à la pose.} Comme $\vect{v}$ est exprimé dans le repère tête, $(\theta, \varphi)$ décrit la rotation de l'œil \emph{dans son orbite}. Si l'on tourne la tête en gardant les yeux sur le même point, $R$ change, et $(\theta, \varphi)$ change exactement de ce qu'il faut pour compenser. C'est cette séparation entre pose de la tête et rotation de l'œil qui rend la v2 robuste : le modèle de regard n'a plus à apprendre l'effet de la tête.

\paragraph{Sortie.} \code{observe} renvoie
\begin{itemize}
  \item \code{O} : le milieu des deux centres de globe, dans le repère caméra (cm), origine du rayon de regard ;
  \item \code{R} : la rotation de la tête ;
  \item \code{ang} : $(\theta, \varphi)$ ;
  \item \code{dirs} : les deux axes séparés (non utilisés en aval).
\end{itemize}
En cas d'erreur numérique (\code{LinAlgError}, \code{ValueError}, \code{ZeroDivisionError}), l'image est ignorée.

\chapter{Tracker v2 : modèle de regard et calibration}
\label{chap:v2mod}

Le chapitre~\ref{chap:v2geo} a produit, pour chaque image, une origine $\vect{O}$ (cm, repère caméra), une rotation de la tête $R$ et deux angles œil/tête $(\theta, \varphi)$. \code{GeoGazeModel} transforme ces mesures en un point $(x, y)$ de la fenêtre, à l'aide de 23 paramètres appris à la calibration.

\section{Paramètres}

\begin{table}[H]
\centering\small
\begin{tabularx}{\textwidth}{@{}c l L@{}}
\toprule
Indices & Nom du bloc & Rôle \\
\midrule
0--5 & \code{P_SCREEN} & application affine du plan de l'écran (cm) vers la fenêtre ($[0,1]^2$) \\
6--12 & \code{P_TH} & 7 coefficients de correction de l'angle horizontal $\theta$ \\
13--19 & \code{P_PH} & 7 coefficients de correction de l'angle vertical $\varphi$ \\
20--21 & \code{P_TILT} & inclinaison du plan de l'écran autour des axes $x$ et $y$ de la caméra (rad) \\
22 & \code{P_KF} & $k_f$, logarithme du facteur de correction de focale / distance \\
\bottomrule
\end{tabularx}
\caption{Les 23 paramètres de \code{GeoGazeModel} (\code{N_PARAMS = 23}).}
\end{table}

\section{Prédiction}

\code{predict_with(p, obs)} est entièrement vectorisée : elle traite d'un coup $N$ observations empilées par \code{stack_obs}, ce qui permet d'évaluer rapidement le modèle sur toute la calibration.

\subsection{Correction des angles}
\label{sec:correction-angles}

Avec la base de 7 fonctions (\code{_basis}, \code{N_BASIS = 7})
\begin{equation}
  \vect{f}(\theta, \varphi, b_v) = \big(1,\ \theta,\ \varphi,\ \theta^2,\ \varphi^2,\ \theta\varphi,\ b_v\big),
\end{equation}
les angles corrigés sont
\begin{equation}
  \theta' = \theta + \vect{f}\cdot\vect{p}_{6..12}, \qquad \varphi' = \varphi + \vect{f}\cdot\vect{p}_{13..19}.
\end{equation}

Chaque terme a une interprétation physique :
\begin{description}
  \item[Constante ($p_6$, $p_{13}$) : angle kappa.] L'axe optique de l'œil (centre du globe $\to$ pupille), que mesure le modèle, ne coïncide pas avec l'axe \emph{visuel} (fovéa $\to$ point fixé). L'écart, appelé angle kappa, vaut typiquement \ang{4} à \ang{5} horizontalement et environ \ang{1} à \ang{2} verticalement, et varie d'une personne à l'autre \cite{guestrin2006,hansen2010}. S'y ajoutent les biais du modèle géométrique (centre du globe approximatif).
  \item[Gains ($p_7$ pour $\theta$, $p_{15}$ pour $\varphi$).] $\theta' \approx (1 + p_7)\,\theta$ : corrige l'erreur sur le rayon du globe (section~\ref{sec:headmodel}) et le biais de localisation de l'iris par MediaPipe, qui a tendance à \enquote{sous-estimer} les grandes excursions.
  \item[Termes croisés ($p_8$ : $\varphi \to \theta$, $p_{14}$ : $\theta \to \varphi$).] Une petite rotation (roll) entre le repère de la tête estimé et le repère réel des yeux mélange les deux axes.
  \item[Termes quadratiques ($p_{9..11}$, $p_{16..18}$).] Ils absorbent les non-linéarités résiduelles : forme non sphérique de la cornée, biais de MediaPipe en bord d'orbite.
  \item[Blendshapes $b_v$ ($p_{12}$, $p_{19}$).] Le score $b_v = \text{lookDown} - \text{lookUp}$ apporte une information verticale indépendante de la position de l'iris. Celle-ci est peu précise verticalement, car l'iris est en partie masqué par les paupières.
\end{description}

\subsection{Direction du regard}

La direction unitaire correspondant à $(\theta', \varphi')$ dans le repère tête, puis dans le repère caméra :
\begin{equation}
  \vect{g}_h = \big(\cos\varphi'\sin\theta',\ \sin\varphi',\ -\cos\varphi'\cos\theta'\big), \qquad \vect{g} = R\,\vect{g}_h .
\end{equation}
C'est l'inverse exact de la définition des angles de la section~\ref{sec:observe}. Pour $\theta' = \varphi' = 0$, $\vect{g}_h = (0, 0, -1)$ : regard droit devant.

\subsection{Correction de distance}
\label{sec:focale}

\begin{equation}
  \vect{O}' = \big(O_x,\ O_y,\ e^{k_f} O_z\big).
\end{equation}

\paragraph{Pourquoi corriger seulement $O_z$ ?} Supposons que la focale réelle vaille $\kappa$ fois la focale supposée. Pour une tête de taille fixe (l'IPD), la distance estimée est $Z \propto f\cdot\text{IPD}/\text{IPD}_\text{px}$ : elle est fausse du même facteur $\kappa$. En revanche, les coordonnées latérales $X = (u - c_x)\,Z / f$ ne changent pas, puisque le rapport $Z/f$ est conservé. Au premier ordre, une erreur de focale revient donc exactement à multiplier la seule composante $z$ de l'origine par un facteur constant. Le logarithme $k_f$ rend ce facteur positif par construction et symétrique ($k_f = \pm\num{0.3}$ correspond à $\times\num{1.35}$ ou $\times\num{0.74}$).

\subsection{Plan de l'écran}

La webcam est fixée au bord de l'écran. On modélise donc l'écran par un plan \emph{passant par le centre optique} de la caméra, de normale
\begin{equation}
  \vect{n} = R_s\,\vect{e}_z, \qquad R_s = R_x(\tau_1)\,R_y(\tau_2), \qquad (\tau_1, \tau_2) = (p_{20}, p_{21}).
\end{equation}
Sans inclinaison, le plan de l'écran est le plan $z = 0$ de la caméra, perpendiculaire à l'axe optique. C'est le cas d'une caméra intégrée qui regarde droit devant. Les deux angles d'inclinaison modélisent une webcam posée en biais, ou un écran de portable incliné par rapport à sa caméra.

Le point d'impact du rayon $\vect{O}' + t\,\vect{g}$ sur le plan $\vect{n}\cdot\vect{X} = 0$ est
\begin{equation}
  t = -\frac{\vect{O}'\cdot\vect{n}}{\vect{g}\cdot\vect{n}}, \qquad \vect{X}_\text{écran} = \vect{O}' + t\,\vect{g}.
\end{equation}
L'utilisateur est devant l'écran ($\vect{O}'\cdot\vect{n} > 0$) et regarde vers lui : $\vect{g}\cdot\vect{n} < 0$. Si le regard est presque parallèle à l'écran, ou s'en éloigne (mesure aberrante, ou paramètres en cours d'ajustement), le dénominateur est borné : $\vect{g}\cdot\vect{n} \leftarrow \min(\vect{g}\cdot\vect{n}, \num{-0.01})$. Cela évite une division par zéro, ou un point projeté \emph{derrière} l'utilisateur, qui déstabiliserait l'ajustement.

Les coordonnées dans le plan, en cm, sont les projections sur les deux autres axes de $R_s$ :
\begin{equation}
  a = \vect{X}_\text{écran}\cdot R_s\vect{e}_x, \qquad b = \vect{X}_\text{écran}\cdot R_s\vect{e}_y .
\end{equation}

\subsection{Application affine vers la fenêtre}

\begin{equation}
  x = p_0\,a + p_1\,b + p_2, \qquad y = p_3\,a + p_4\,b + p_5 .
\end{equation}
Cette application encode tout ce que l'on ignore de l'installation : taille physique de la fenêtre (conversion cm $\to$ fraction de fenêtre), position de la fenêtre par rapport à la caméra, petite rotation ou cisaillement résiduel. Elle n'a pas d'a priori : les 16 cibles de calibration, réparties sur toute la fenêtre, la déterminent très bien.

\section{Ajustement}

\subsection{Fonction de coût}

Avec $N$ observations $(\text{obs}_i, \vect{y}_i, \omega_i)$, \code{GeoGazeModel.fit} minimise
\begin{equation}
  J(\vect{p}) = \sum_{i=1}^{N} \omega_i \sum_{c\in\{x,y\}} \rho\!\left(\frac{\hat y_{ic}(\vect{p}) - y_{ic}}{\sigma_\text{fit}}\right)
   + \sum_{j \in \mathcal{P}} \left(\frac{p_j}{\sigma_j}\right)^2,
  \qquad \rho(r) = 2\big(\sqrt{1 + r^2} - 1\big),
\end{equation}
où $\mathcal{P}$ est l'ensemble des 17 paramètres munis d'un a priori : les 14 corrections d'angle, les deux inclinaisons et $k_f$.

\paragraph{Mise en œuvre.} L'algorithme de Levenberg-Marquardt minimise une somme de carrés $\|\vect{r}\|^2$. On fournit donc le vecteur
\begin{equation}
  \vect{r}(\vect{p}) = \Big[\ \sqrt{\omega_i}\,\sign(r_{ic})\sqrt{\rho(r_{ic})}\ \Big]_{i,c} \ \oplus\ \Big[\ p_j/\sigma_j\ \Big]_{j\in\mathcal P},
  \qquad r_{ic} = \frac{\hat y_{ic} - y_{ic}}{\sigma_\text{fit}},
\end{equation}
de sorte que $\|\vect{r}\|^2 = J$. Le facteur $\sign(\cdot)$ préserve le signe du résidu, ce qui garde une jacobienne régulière autour de 0.

\subsection{Perte soft-L1}

\begin{figure}[H]
\centering
\begin{tikzpicture}
\begin{axis}[width=0.62\textwidth, height=5.2cm, xmin=-5, xmax=5, ymin=0, ymax=10, samples=200,
  xlabel={résidu normalisé $r$}, ylabel={coût}, legend style={font=\small, at={(0.5,0.97)}, anchor=north},
  grid=major, grid style={black!8}, tick label style={font=\small}]
  \addplot[thick, doux, dashed, domain=-5:5] {x^2}; \addlegendentry{$r^2$ (moindres carrés)}
  \addplot[very thick, accent, domain=-5:5] {2*(sqrt(1+x^2)-1)}; \addlegendentry{$\rho(r)$ soft-L1}
\end{axis}
\end{tikzpicture}
\caption{La perte soft-L1 est quadratique près de 0 et linéaire au-delà de $|r| \approx 1$, c'est-à-dire d'une erreur d'environ $\sigma_\text{fit}$.}
\end{figure}

Pour $|r| \ll 1$, $\rho(r) \approx r^2$ : on retrouve les moindres carrés. Pour $|r| \gg 1$, $\rho(r) \approx 2|r|$ : la perte devient linéaire, comme une norme $L^1$. La fonction d'influence $\psi(r) = \tfrac12\rho'(r) = r/\sqrt{1 + r^2}$ est bornée par 1 : un échantillon aberrant (saccade non filtrée, mauvaise détection) ne peut pas tirer le modèle plus fort qu'un échantillon à $\approx 1\,\sigma_\text{fit}$. C'est la même perte que l'option \code{loss='soft_l1'} de SciPy, avec une échelle de \num{1} appliquée au résidu déjà normalisé.

\begin{origine}
\code{FIT_SIGMA = 0.03} (\empi) : bruit typique d'un échantillon, en fraction de fenêtre. C'est l'ordre de grandeur de la dispersion du regard brut pendant une fixation avec une webcam. Il fixe à la fois l'échelle de transition de la perte soft-L1 et le poids relatif des données face aux a priori.
\end{origine}

\subsection{A priori et interprétation bayésienne}

Si les erreurs étaient gaussiennes d'écart-type $\sigma_\text{fit}$ et si chaque paramètre suivait une loi a priori $\mathcal N(0, \sigma_j^2)$, minimiser $J$ (avec $\rho(r) = r^2$) reviendrait exactement à trouver le maximum a posteriori. Le coût a la forme $-2\log p(\text{données}\mid\vect{p}) - 2\log p(\vect{p})$ à une constante près, et les deux termes ont le même facteur. Les a priori ramènent vers 0 (aucune correction, géométrie nominale) les paramètres que les données déterminent mal.

\paragraph{Ordre de grandeur.} Les poids sont normalisés à une moyenne de 1 : le terme de données vaut environ $N \approx \num{1000}$ fois un résidu typique, contre 17 termes d'a priori. Ces a priori sont donc \emph{faibles}, comme l'indique la documentation du module. Ils ne pèsent que dans les directions où les données ne contraignent presque rien, par exemple un échange entre inclinaison de l'écran et application affine.

\begin{table}[H]
\centering\small
\begin{tabularx}{\textwidth}{@{}l c c L@{}}
\toprule
Paramètre & $\sigma_j$ ($\theta$) & $\sigma_j$ ($\varphi$) & Interprétation (\empi, ordres de grandeur physiques) \\
\midrule
constante (kappa) & \num{0.15} & \num{0.15} & \qty{0.15}{\radian} $\approx$ \ang{8.6} : couvre largement l'angle kappa (\ang{4}--\ang{5}) et les biais du modèle. \\
gain propre & \num{0.5} & \num{0.6} & \qtyrange{50}{60}{\percent} d'erreur de gain tolérés sans effort : le rayon du globe et le biais de MediaPipe sont incertains. Le vertical est moins bien mesuré, d'où une tolérance plus grande. \\
terme croisé & \num{0.2} & \num{0.2} & couplage $\theta \leftrightarrow \varphi$ de \qty{20}{\percent} : correspond à un roll de $\approx$ \ang{11}. \\
quadratiques & 1 & 1 & quasi libres : les courbures sont mal connues a priori. \\
$b_v$ & \num{0.2} & \num{0.4} & le score vertical ne devrait guère influencer l'horizontal ; il est plus utile en vertical. \\
\midrule
inclinaisons $\tau_1, \tau_2$ & \multicolumn{2}{c}{\num{0.35}} & \qty{0.35}{\radian} $\approx$ \ang{20} : inclinaison plausible d'une webcam ou d'un écran de portable. \\
$k_f$ & \multicolumn{2}{c}{\num{0.3}} & facteur de distance entre $\times\num{0.74}$ et $\times\num{1.35}$ à $1\sigma$ : couvre un HFOV réel de \ang{45} à \ang{75} environ. \\
affine & \multicolumn{2}{c}{$\infty$} & aucun a priori : bien déterminée par les cibles. \\
\bottomrule
\end{tabularx}
\caption{A priori \code{PRIOR_SIGMA}. Les angles sont en radians, les gains sans unité.}
\end{table}

\subsection{Point de départ}

L'algorithme de Levenberg-Marquardt ne trouve que des minima locaux ; il faut un bon point de départ. \code{fit} l'obtient en deux temps :
\begin{enumerate}
  \item toutes les corrections, inclinaisons et $k_f$ sont mis à 0, et l'affine à l'identité. Le modèle prédit alors les coordonnées $(a, b)$ du point dans le plan de l'écran, en cm, selon la géométrie nominale ;
  \item on résout le \emph{sous-problème linéaire} des 6 paramètres affines par moindres carrés pondérés : $[\,a_i\ \ b_i\ \ 1\,]\,C \approx \vect{y}_i$, avec $C \in \R^{3\times2}$.
\end{enumerate}
Le départ correspond donc déjà à la meilleure correspondance affine entre la géométrie nominale et les cibles, et l'optimisation non linéaire n'a plus qu'à affiner.

\subsection{Algorithme de Levenberg-Marquardt}

Le code n'importe pas SciPy ; \code{levenberg_marquardt} est une implémentation autonome, qui évite une dépendance lourde sur Raspberry Pi.

\begin{algorithm}[H]
\caption{\code{levenberg_marquardt(fun, p, iters, tol)}}
$\vect{r} \leftarrow \text{fun}(\vect{p})$ ; \ $C \leftarrow \|\vect{r}\|^2$ ; \ $\mu \leftarrow 10^{-3}$\;
\For{$k = 1..\text{iters}$}{
  \For(\tcp*[f]{jacobienne par différences finies avant}){$j = 1..\dim \vect{p}$}{
    $h_j \leftarrow 10^{-6}\max(1, |p_j|)$ ; \quad $J_{:,j} \leftarrow \big(\text{fun}(\vect{p} + h_j\vect{e}_j) - \vect{r}\big)/h_j$\;
  }
  $A \leftarrow J\T J$ ; \ $\vect{g} \leftarrow J\T\vect{r}$ ; \ $D \leftarrow \diag(A) + 10^{-9}$\;
  \Repeat{$C' < C$ \textbf{et} $C'$ fini}{
    $\vect{p}' \leftarrow \vect{p} - (A + \mu\,\diag(D))^{-1}\vect{g}$ ; \ $C' \leftarrow \|\text{fun}(\vect{p}')\|^2$\;
    \lIf{échec}{$\mu \leftarrow 4\mu$ ; \lIf{$\mu > 10^{10}$}{\Return $\vect{p}$}}
  }
  $\text{fini} \leftarrow (C - C' < \text{tol}\cdot C)$\;
  $\vect{p} \leftarrow \vect{p}'$ ; \ $C \leftarrow C'$ ; \ $\mu \leftarrow \max(\mu/3, 10^{-9})$\;
  \lIf{fini}{\textbf{arrêter}}
}
\Return $\vect{p}$
\end{algorithm}

\paragraph{Principe.} À chaque itération, on linéarise les résidus, $\vect{r}(\vect{p} + \vect\delta) \approx \vect{r} + J\vect\delta$, et on résout le problème de Gauss-Newton amorti
\begin{equation}
  \big(J\T J + \mu\,\diag(J\T J)\big)\,\vect\delta = -J\T\vect{r}.
\end{equation}
Pour $\mu \to 0$, c'est un pas de Gauss-Newton, rapide près du minimum. Pour $\mu$ grand, c'est une petite descente de gradient, sûre loin du minimum. $\mu$ diminue après chaque pas réussi et augmente après chaque échec \cite{levenberg1944,marquardt1963}.

\paragraph{Choix d'implémentation.}
\begin{itemize}
  \item \textbf{Amortissement de Marquardt} $\mu\,\diag(J\T J)$ plutôt que $\mu I$ : le pas est invariant à l'échelle des paramètres. C'est important ici, où des gains sans unité côtoient des coefficients affines de l'ordre de \num{0.03} par cm (l'inverse de la largeur d'un écran d'environ \qty{30}{\centi\metre}). Le terme $10^{-9}$ évite une diagonale nulle pour un paramètre inactif.
  \item \textbf{Différences finies} avec un pas relatif $10^{-6}$, de l'ordre de la racine de la précision machine ($\sqrt{\epsilon} \approx \num{1.5e-8}$) multipliée par une marge pour le bruit d'arrondi. Le modèle est trop imbriqué (rotations, intersections) pour que l'écriture d'une jacobienne analytique en vaille la peine.
  \item \textbf{Facteurs $\times4$ et $/3$} (\empi) : valeurs proches des usages courants (souvent $\times10$ et $/10$). Une remontée plus douce évite des pas trop timides après un échec isolé.
  \item \textbf{Arrêt} quand la baisse relative du coût passe sous $10^{-8}$, ou après \code{iters} itérations (60 pour la calibration, 15 pour la recalibration, 3 pour la validation croisée).
\end{itemize}

\section{Calibration v2}

\subsection{Protocole en deux phases}

\code{Calibration3D} enchaîne deux phases (tableau~\ref{tab:cal3d}).

\begin{table}[H]
\centering\small
\begin{tabularx}{\textwidth}{@{}l L L l@{}}
\toprule
 & Tête fixe & Tête en mouvement & Origine \\
\midrule
Cibles & 16 (grille $4\times4$, ordre aléatoire) & 3 : $(\frac12, \frac12)$, $(\frac14, \num{0.3})$, $(\frac34, \num{0.7})$ & \empi \\
Stabilisation & \qty{0.8}{\second} & \qty{1.0}{\second} (\code{CAL_HEAD_SETTLE_S}) & \empi \\
Collecte & \qty{2.0}{\second} & \qty{4.5}{\second} (\code{CAL_HEAD_COLLECT_S}) & \empi \\
Minimum d'échantillons & 15 & 30 (\code{CAL_HEAD_MIN_SAMPLES}) & \empi \\
Poids d'une cible & 1 & 3 (\code{HEAD_GROUP_WEIGHT}) & \empi \\
Consigne & \enquote{Fixe le point, tête immobile} & \enquote{Garde les yeux sur le point et bouge doucement la tête} & -- \\
\bottomrule
\end{tabularx}
\caption{Protocole de calibration v2.}
\label{tab:cal3d}
\end{table}

Durée totale sans reprise : $16 \times 2.8 + 3 \times 5.5 = \qty{61.3}{\second}$. Il n'y a pas de temps de lecture initial : la consigne est affichée en permanence sous le titre (\code{calibration_hint}, affichée par \code{afficher_consigne} dans \fichier{main_v2.py}). Comme en v1, une cible qui n'a pas réuni assez d'échantillons est reprise jusqu'à deux fois, puis ignorée.

\paragraph{Les cibles de la phase tête en mouvement.} Le centre, puis deux points en diagonale opposée, à mi-chemin des bords : on teste la compensation de la tête dans trois directions de regard différentes, sans viser les bords, où l'iris est le plus mal mesuré. La collecte dure plus longtemps (\qty{4.5}{\second}) pour laisser le temps de faire un mouvement complet (gauche-droite, haut-bas, avant-arrière, selon la consigne affichée dans la console).

\paragraph{Pourquoi cette phase est indispensable.} Tête immobile, la distance $Z_0$ est fixe et l'angle $\theta$ varie seulement avec la cible. Le point de l'écran vaut alors, en horizontal, environ $x \approx p_0\,Z_0\tan\!\big((1 + p_7)\,\theta\big) + \dots$ : le \emph{gain} de l'œil $p_7$, l'échelle de l'affine $p_0$ et la correction de distance $k_f$ (qui modifie $Z_0$) ont des effets presque interchangeables. Les données ne permettent pas de les séparer. Quand la tête bouge en gardant les yeux sur une cible fixe, $\theta$ varie alors que la cible ne bouge pas : seule la bonne combinaison de gain et de distance fait que le rayon continue à frapper le même point. Ces échantillons \emph{lèvent l'ambiguïté} ; ils reçoivent un poids triple parce qu'ils ne concernent que 3 cibles sur 19.

\subsection{Données brutes}

Pendant la calibration, aucun modèle de tête n'existe encore : on ne peut pas calculer $(\theta, \varphi)$. La v2 stocke donc les \emph{landmarks bruts} de chaque échantillon, avec $b_v$, la cible, l'indice de cible et le type (\code{fix} ou \code{head}). Tout est recalculé à la fin avec le modèle de tête construit sur l'ensemble des images.

\subsection{Ajustement}

\begin{algorithm}[H]
\caption{\code{Calibration3D.fit(cam)}}
\lIf{moins de 10 cibles fixes validées}{\Return échec}
construire le \code{HeadModel} sur toutes les images (section~\ref{sec:headmodel})\;
\ForEach{échantillon}{calculer \code{observe} ; ignorer s'il échoue}
\ForEach(\tcp*[f]{rejet des aberrants}){cible fixe $g$}{
  $\hat\sigma_c \leftarrow \max(\num{1.4826}\MAD(\text{ang}_c), \num{0.003})$ pour $c \in \{\theta, \varphi\}$\;
  rejeter les échantillons avec $|\text{ang}_c - \med(\text{ang}_c)| > 3.5\,\hat\sigma_c$\;
}
$\omega_i \leftarrow H_{g_i}/n_{g_i}$ avec $H_g = 3$ (tête en mouvement) ou 1 (fixe) ; \ $\omega \leftarrow \omega/\bar\omega$\;
ajuster \code{GeoGazeModel} (60 itérations au plus)\;
erreurs d'entraînement (cibles fixes) et tête en mouvement : moyenne de $|\hat{\vect{y}} - \vect{y}|$\;
\ForEach(\tcp*[f]{LOPO approchée}){cible fixe $g$}{
  $\omega^{(g)} \leftarrow \omega$ avec $\omega^{(g)}_i = 0$ pour $i \in g$\;
  réajuster en 3 itérations à partir des paramètres complets ; erreur sur $g$\;
}
\Return modèle de tête, modèle de regard, erreurs\;
\end{algorithm}

\paragraph{Rejet sur les angles, et seulement pour les cibles fixes.} Pour une cible fixe et une tête immobile, les angles œil/tête sont presque constants : un échantillon qui s'en écarte est une saccade ou un clignement partiel. Le plancher \qty{0.003}{\radian} (\ang{\approx 0.17}, \empi) évite de rejeter des échantillons sains quand la dispersion est très faible ; \num{3.5} est le seuil d'Iglewicz et Hoaglin \cite{iglewicz1993}. Pour les cibles \enquote{tête en mouvement}, les angles varient au contraire volontairement, puisque l'œil contre-tourne pour compenser la tête : le même test rejetterait les échantillons les plus utiles.

\paragraph{Poids.} Chaque cible fixe pèse au total 1 et chaque cible tête en mouvement pèse 3, quel que soit son nombre d'images. La normalisation à une moyenne de 1 garde le terme de données à la même échelle que dans l'analyse des a priori ci-dessus.

\paragraph{Validation croisée approchée.} Réajuster complètement 16 fois coûterait 16 fois l'ajustement principal. Le code part des paramètres complets et ne fait que 3 itérations sans la cible testée : c'est une approximation de la vraie erreur LOPO, légèrement optimiste, mais beaucoup moins chère.

\subsection{Résultat}

En cas de succès, \code{_finish_calibration} installe le modèle de tête et le modèle de regard, garde une copie du modèle comme \emph{modèle de base} pour la recalibration (chapitre~\ref{chap:recal}), remet à zéro le décalage de recentrage, l'historique et les filtres, et affiche :
\begin{itemize}
  \item l'erreur validée (LOPO) et l'erreur d'entraînement en $x$ et $y$, en \% de la fenêtre ;
  \item l'erreur sur la phase tête en mouvement ;
  \item le nombre de cibles ignorées et la durée de l'ajustement ;
  \item les paramètres physiques appris : gains de l'œil ($1 + p_7$, $1 + p_{15}$), facteur de focale ($e^{k_f}$) et inclinaison de l'écran ($\tau_1$ en degrés).
\end{itemize}
Ces valeurs servent de diagnostic. Un gain très loin de 1 ou une inclinaison de \ang{40} signalent une calibration douteuse, ou une installation inhabituelle.

\chapter{Recalibration continue au clic}
\label{chap:recal}

Une calibration n'est jamais parfaite, et elle se dégrade avec le temps : l'utilisateur change de position, la lumière change, la webcam bouge. Chaque clic est une occasion de corriger le modèle, puisqu'on sait à peu près ce que l'utilisateur regardait juste avant de cliquer. C'était le premier point de la feuille de route (\enquote{recentrer à chaque clic en milieu de la cellule}), réalisé dans le commit \code{797601f}.

\section{Côté interface : la cible du point}

Quand un bouton est déclenché, \fichier{main_v2.py} ajoute un point de calibration si les quatre conditions suivantes sont réunies :
\begin{itemize}
  \item \code{recalibration_clic} est activé ;
  \item on n'est pas en mode souris ;
  \item le tracker est calibré ;
  \item le clic porte sur une case : les bandes de défilement sont exclues, car elles n'ont pas de bouton dont on pourrait viser le centre.
\end{itemize}
La cible transmise est un mélange du centre de la case $\vect{c}$ et de la position du regard au moment du clic $\vect{g}$ (en pixels, puis normalisés par $(w-1, h-1)$) :
\begin{equation}
  \vect{y}^\star = \alpha\,\vect{c} + (1 - \alpha)\,\vect{g}, \qquad \alpha = \code{recalibration_part_centre}.
\end{equation}

\subsection{Interprétation de $\alpha$}

Ce mélange est un estimateur \emph{à rétrécissement}. Supposons que l'utilisateur fixe un point $\vect{y}$ de la case, distribué autour du centre avec un écart-type $\sigma_U$ (on ne vise jamais exactement le centre d'un grand bouton), et que le regard mesuré soit $\vect{g} = \vect{y} + \vect\varepsilon$, où $\vect\varepsilon$ est l'erreur du modèle, d'écart-type $\sigma_\varepsilon$. Si ces deux lois sont gaussiennes et indépendantes, la meilleure estimation du point réellement fixé est
\begin{equation}
  \mathbb E[\vect{y} \mid \vect{g}] = \vect{c} + \frac{\sigma_U^2}{\sigma_U^2 + \sigma_\varepsilon^2}\,(\vect{g} - \vect{c})
  \quad\Longleftrightarrow\quad
  \alpha^\star = \frac{\sigma_\varepsilon^2}{\sigma_U^2 + \sigma_\varepsilon^2}.
\end{equation}
\begin{itemize}
  \item $\alpha = 1$ revient à supposer que l'utilisateur vise toujours exactement le centre ($\sigma_U = 0$). Les points sont alors très informatifs mais biaisés, dès que l'utilisateur vise ailleurs (par exemple le texte du bouton, décentré si la zone morte masque le centre).
  \item $\alpha = 0$ revient à croire le modèle ($\sigma_\varepsilon = 0$) : la cible vaut le regard prédit, et la recalibration ne corrige rien.
  \item La valeur de \fichier{config.yaml}, $\alpha = \num{0.65}$, correspond à $\sigma_\varepsilon \approx 1.36\,\sigma_U$ : on estime l'erreur du modèle un peu supérieure à la dispersion de la visée dans la case. Le défaut de \fichier{config.py} est $\alpha = \num{0.5}$, soit $\sigma_\varepsilon = \sigma_U$.
\end{itemize}
Ces valeurs sont \empi : aucune mesure de $\sigma_U$ ni de $\sigma_\varepsilon$ n'a été faite. La formule indique comment les régler si l'on mesure un jour ces deux dispersions.

\section{Côté tracker : constitution d'un point}

\code{GazeTracker.add_calibration_point(cible, t)} met en file une commande qui, dans le thread tracker :
\begin{enumerate}
  \item vérifie qu'un modèle de base existe et qu'aucune calibration n'est en cours ;
  \item sélectionne dans l'historique \code{_recent} les observations dont l'instant est dans $[t - \Delta, t]$, avec $\Delta = \code{recalibration_fenetre}$ ;
  \item abandonne s'il y en a moins de 3 ;
  \item empile ces observations et ajoute le couple $(\text{observations},\ \vect{y}^\star - \vect{o})$ à la fenêtre glissante \code{recal_points}, où $\vect{o}$ est le décalage de recentrage courant. Le modèle prédit le regard \emph{avant} ce décalage : il faut donc le retirer de la cible pour rester cohérent ;
  \item réajuste la correction (\code{_refit_recal}).
\end{enumerate}

\paragraph{L'historique \code{_recent}.} Pendant le suivi, chaque image exploitable ajoute $(t, \text{obs}, b_v)$ à une file bornée à 300 éléments, soit environ \qty{10}{\second} à \qty{30}{images\per\second} (\empi, largement supérieur au besoin). Ce sont les observations géométriques complètes ($\vect{O}$, $R$, $\theta$, $\varphi$) et non les points prédits : on peut ainsi réévaluer le modèle avec d'autres paramètres.

\begin{origine}
\textbf{Fenêtre} $\Delta$ : \qty{0.4}{\second} dans \fichier{config.yaml}, \qty{0.5}{\second} par défaut dans \fichier{config.py} et \code{RECAL_WINDOW_S} (\empi). Le clic intervient après \code{duree_clic} $= \qty{2}{\second}$ de fixation. Les dernières \qty{0.4}{\second} sont donc prises au cœur de la fixation : loin de la saccade d'entrée dans la case, et assez récentes pour refléter la pose actuelle. Elles contiennent environ 12 images ; le minimum de \textbf{3 images} écarte les cas où le visage a été perdu juste avant le clic.
\end{origine}

\section{Ajustement de la correction}

\subsection{Paramètres corrigés}

Seuls quatre des 23 paramètres sont modifiés (\code{RECAL_PARAMS = [6, 13, 7, 15]}) :
\begin{center}
\begin{tabular}{@{}cll@{}}
\toprule
Indice & Paramètre & Effet \\
\midrule
6 & constante de $\theta$ & décalage horizontal de l'axe visuel (kappa) \\
13 & constante de $\varphi$ & décalage vertical de l'axe visuel \\
7 & coefficient de $\theta$ dans la correction de $\theta$ & gain horizontal de l'œil \\
15 & coefficient de $\varphi$ dans la correction de $\varphi$ & gain vertical de l'œil \\
\bottomrule
\end{tabular}
\end{center}

\paragraph{Pourquoi corriger les angles et pas l'écran ?} Un simple décalage du curseur à l'écran (comme celui du recentrage, ou le terme constant de l'affine) n'est juste qu'\emph{à la pose actuelle}. Une erreur d'angle de \ang{1} produit à l'écran un décalage proportionnel à la distance de la tête, et dont la direction suit la rotation de la tête. En corrigeant l'erreur \emph{là où elle naît}, dans les angles œil/tête, la correction reste valable quand l'utilisateur bouge. Se limiter à quatre paramètres, les mieux déterminés par quelques points, évite de sur-ajuster des clics bruités.

\subsection{Coût}

On note $\vect{p}^\text{base}$ les paramètres de la calibration et $\vect\delta \in \R^4$ la correction, de sorte que $\vect{p} = \vect{p}^\text{base} + E\vect\delta$, où $E$ place $\vect\delta$ aux indices ci-dessus. Pour une fenêtre de $K$ points, numérotés du plus ancien ($k = 0$) au plus récent ($k = K-1$), le point $k$ ayant $n_k$ images :
\begin{equation}
  J(\vect\delta) = \sum_{k=0}^{K-1} \sum_{i=1}^{n_k} \frac{W\,d^{\,K-1-k}}{n_k}
   \sum_{c\in\{x,y\}} \rho\!\left(\frac{\hat y_{kic}(\vect{p}^\text{base} + E\vect\delta) - y^\star_{kc}}{\sigma_\text{fit}}\right)
   + \sum_{m=1}^{4}\left(\frac{\delta_m}{s_m}\right)^2,
\end{equation}
avec $W = \code{recalibration_poids}$, $d = \code{recalibration_oubli}$ et $\vect{s} = \code{RECAL_SIGMA} = (\num{0.15}, \num{0.15}, \num{0.2}, \num{0.2})$. On réutilise la même perte soft-L1 $\rho$ et le même $\sigma_\text{fit}$ qu'à la calibration.

\begin{itemize}
  \item \textbf{Chaque point pèse au total $W d^{\,\text{âge}}$}, quel que soit son nombre d'images, grâce à la division par $n_k$.
  \item \textbf{La calibration reste l'ancre.} Ses échantillons ne sont pas réutilisés ; c'est l'a priori $\delta_m / s_m$ qui ramène la correction vers 0, c'est-à-dire vers le modèle de base. Les écarts-types reprennent ceux de la calibration : \qty{0.15}{\radian} pour les décalages, \num{0.2} pour les gains, un peu plus serré que les \numrange{0.5}{0.6} de la calibration parce qu'on corrige ici un modèle \emph{déjà} ajusté.
  \item \textbf{Optimisation} : Levenberg-Marquardt, 15 itérations au plus, en partant de la correction précédente (démarrage à chaud). Les corrections successives restent ainsi continues et le calcul est rapide : la docstring parle de \enquote{quelques ms}. Le nouveau modèle est $\vect{p}^\text{base} + E\vect\delta$.
\end{itemize}

\subsection{Mémoire et oubli}

La somme des poids d'une longue suite de clics est $W\sum_k d^k \to W/(1-d)$ : la fenêtre se comporte comme une moyenne exponentielle sur un \emph{nombre effectif} de clics $1/(1-d)$. La \emph{demi-vie}, nombre de clics après lequel un point ne pèse plus que la moitié, vaut $\ln\frac12 / \ln d$.

\begin{table}[H]
\centering\small
\begin{tabular}{@{}lcccc@{}}
\toprule
$d$ & Nombre effectif $1/(1-d)$ & Demi-vie (clics) & Poids du 20\textsuperscript{e} point $d^{19}$ & Source \\
\midrule
\num{0.80} & \num{5.0} & \num{3.1} & \num{0.014} & \code{RECAL_DECAY}, \fichier{config.py} \\
\num{0.85} & \num{6.7} & \num{4.3} & \num{0.046} & \fichier{config.yaml} \\
\bottomrule
\end{tabular}
\caption{Effet du facteur d'oubli.}
\end{table}

Avec ces valeurs, les points au-delà du vingtième pèseraient moins de \qty{5}{\percent} : la fenêtre de 20 points (\code{recalibration_points_max}) borne le coût de calcul sans changer le résultat de façon perceptible.

\subsection{Force d'une correction}

Pour sentir l'équilibre entre données et a priori, on prend un cas simple : un décalage horizontal d'angle $\delta$, qui produit à l'écran un décalage $\gamma\delta$. À \qty{60}{\centi\metre} de distance, \ang{1} ($\qty{0.0175}{\radian}$) donne environ \qty{1.05}{\centi\metre}, soit \qty{3.5}{\percent} d'un écran de \qty{30}{\centi\metre} de large : $\gamma \approx 2$ (fraction d'écran par radian). Dans le régime quadratique, avec un seul point de poids $W = 1$ et une erreur $e$, minimiser $(e - \gamma\delta)^2/\sigma_\text{fit}^2 + \delta^2/s^2$ donne une fraction corrigée de
\begin{equation}
  \frac{\gamma\delta^\star}{e} = \frac{\gamma^2/\sigma_\text{fit}^2}{\gamma^2/\sigma_\text{fit}^2 + 1/s^2}
  = \frac{4/\num{0.0009}}{4/\num{0.0009} + 1/\num{0.0225}} \approx \num{0.99} .
\end{equation}
\emph{Un seul clic corrige donc presque entièrement un décalage}. La recalibration est très réactive. Le revers est que le bruit de visée d'un clic passe lui aussi presque entièrement dans la correction. Ce sont le mélange avec le centre ($\alpha$) et la moyenne sur plusieurs clics ($d$) qui filtrent ce bruit. Pour une recalibration plus prudente, il faut baisser \code{recalibration_poids}.

\section{Remise à zéro}

\code{_clear_recal} vide les points, l'historique et la correction, et restaure le modèle de base. Elle est appelée après chaque nouvelle calibration (touche C), au reset (touche R, qui efface aussi le modèle de base) et au chargement d'une calibration.

\section{Sauvegarde et chargement d'une calibration}

\code{save_calibration(path)} écrit un fichier \code{.npz} NumPy contenant :
\begin{itemize}
  \item \code{pts} : la forme 3D rigide (69 points, cm, repère tête) ;
  \item \code{eyes} : les deux centres de globe ;
  \item \code{f_ratio} : $f/w$, pour retrouver la focale à une autre résolution ;
  \item \code{params} : les 23 paramètres du \emph{modèle de base}. Les corrections de recalibration, liées à une session, ne sont pas sauvegardées.
\end{itemize}

\code{load_calibration(path)} reconstruit la caméra si aucune image n'a encore été lue, à partir de la résolution demandée à OpenCV, puis le modèle de tête et le modèle de regard. Elle remet aussi à zéro la recalibration, le décalage et les filtres.

\begin{attention}
Ces deux méthodes existent mais \textbf{ne sont appelées nulle part} : il faut recalibrer à chaque lancement. Les brancher au démarrage et en fin de calibration fait partie de la feuille de route. Une calibration rechargée suppose la même webcam, à la même place, et le même utilisateur ; la recalibration au clic corrige ensuite les petits écarts.
\end{attention}

\chapter{Traitements communs du signal de regard}
\label{chap:traitements}

Les mécanismes de ce chapitre sont identiques dans les deux trackers (méthode \code{step()}). Ils transforment une suite de mesures brutes, bruitées et parfois absentes, en un signal de regard exploitable par l'interface.

\section{Clignements}
\label{sec:clignements}

Pendant un clignement, la paupière masque l'iris et MediaPipe place le centre de l'iris trop bas. Sans traitement, chaque clignement ferait plonger le curseur vers le bas de l'écran, puis le ferait remonter.

\paragraph{Règle.} Une image est déclarée \enquote{clignement} si $s_\text{blink} > \code{seuil_clignement}$ (\num{0.6} par défaut). On mémorise alors $t_\text{fin} = t + \qty{0.15}{\second}$ (\code{BLINK_HOLD_S}), et une image n'est exploitable que si elle n'est pas un clignement \emph{et} si $t \geq t_\text{fin}$. Les images non exploitables ne modifient ni le regard publié, ni les filtres, ni l'historique ; pendant une calibration, elles ne sont pas collectées.

\begin{origine}
\textbf{Seuil \num{0.6}} (\empi, réglable de 0 à 1) : le score \code{eyeBlink} est proche de 0 œil grand ouvert et proche de 1 œil fermé, mais il monte aussi quand on regarde vers le bas, parce que la paupière s'abaisse. Un seuil au-dessus de la moitié évite de jeter les regards vers le bas de l'écran, tout en détectant les vrais clignements. Ces valeurs typiques ne sont pas mesurées dans le dépôt. On peut le baisser pour une personne aux paupières tombantes, ou s'il reste des sauts du curseur.

\textbf{Maintien \qty{0.15}{\second}} (\litt{} / \empi) : un clignement dure de l'ordre de \qtyrange{100}{400}{\milli\second}, et la réouverture de la paupière est nettement plus lente que la fermeture. Pendant cette réouverture, le score repasse sous le seuil alors que l'iris est encore partiellement couvert. \qty{150}{\milli\second}, soit environ 5 images, couvrent cette phase.
\end{origine}

\section{Filtre One Euro}
\label{sec:oneeuro}

\subsection{Le compromis à résoudre}

Le regard brut est bruité : le centre de l'iris est localisé à environ un pixel près, ce qui fait trembler le curseur de quelques pour cent de l'écran en fixation. Un filtre passe-bas fort supprime ce tremblement, mais retarde le curseur lors des saccades. Un filtre faible est réactif, mais tremble. Le filtre One Euro \cite{casiez2012} résout ce compromis en \emph{adaptant} la fréquence de coupure à la vitesse : forte stabilité à l'arrêt, faible retard en mouvement.

\subsection{Dérivation}

Un filtre passe-bas du premier ordre (circuit RC) de constante de temps $\tau = 1/(2\pi f_c)$ obéit à $\tau\,\dot{\hat x} = x - \hat x$. Discrétisé avec une période $T_e$ (schéma implicite), il donne
\begin{equation}
  \hat x_t = \hat x_{t-1} + \frac{T_e}{\tau + T_e}(x_t - \hat x_{t-1}) = \alpha\,x_t + (1 - \alpha)\,\hat x_{t-1},
  \qquad \alpha(T_e, f_c) = \frac{1}{1 + \tau/T_e} .
\end{equation}
Le filtre One Euro applique ce lissage avec une coupure qui dépend de la vitesse estimée :
\begin{align}
  \dot x_t &= (x_t - \hat x_{t-1})/T_e && \text{vitesse brute} \\
  \hat{\dot x}_t &= \alpha(T_e, f_d)\,\dot x_t + \big(1 - \alpha(T_e, f_d)\big)\,\hat{\dot x}_{t-1} && \text{vitesse lissée à } f_d \\
  f_c &= f_{\min} + \beta\,|\hat{\dot x}_t| && \text{coupure adaptative} \\
  \hat x_t &= \alpha(T_e, f_c)\,x_t + \big(1 - \alpha(T_e, f_c)\big)\,\hat x_{t-1} && \text{sortie}
\end{align}
La période $T_e$ est mesurée à chaque échantillon (différence des instants, bornée à \qty{1}{\milli\second}). Le filtre reste donc correct si la cadence de la caméra varie ou si des images sont sautées (clignements). Le premier échantillon après une remise à zéro est renvoyé tel quel.

\subsection{Réglages et ordres de grandeur}

Les coordonnées sont en fractions de fenêtre ; la vitesse est donc en fenêtres par seconde.

\begin{figure}[H]
\centering
\begin{tikzpicture}
\begin{axis}[width=0.62\textwidth, height=5.4cm, xmin=0, xmax=6, ymin=0, ymax=300, samples=200,
  xlabel={vitesse lissée $|\hat{\dot x}|$ (fenêtres/s)}, ylabel={constante de temps $\tau$ (ms)},
  grid=major, grid style={black!8}, tick label style={font=\small}, label style={font=\small}]
  \addplot[very thick, accent, domain=0:6] {1000/(2*pi*(0.6+3*x))};
  \node[font=\scriptsize, anchor=west] at (axis cs:0.15,265) {repos : $\tau \approx \qty{265}{\milli\second}$};
  \node[font=\scriptsize, anchor=west] at (axis cs:1.05,60) {1~fenêtre/s : $\tau \approx \qty{44}{\milli\second}$};
\end{axis}
\end{tikzpicture}
\caption{Constante de temps du filtre en fonction de la vitesse, avec $f_{\min} = \qty{0.6}{\hertz}$ et $\beta = 3$.}
\end{figure}

\begin{origine}
\code{ONE_EURO_MIN_CUTOFF = 0.6} Hz, \code{ONE_EURO_BETA = 3.0}, \code{ONE_EURO_D_CUTOFF = 1.0} Hz. Le premier est réglable par \code{filtre_stabilite}, le deuxième par \code{filtre_reactivite}.
\begin{itemize}
  \item $f_d = \qty{1}{\hertz}$ est la valeur recommandée par les auteurs du filtre \cite{casiez2012} (\litt).
  \item $f_{\min}$ et $\beta$ se règlent selon la procédure de l'article : $\beta = 0$ d'abord, puis $f_{\min}$ baissé jusqu'à ce que le tremblement en fixation disparaisse ; enfin $\beta$ augmenté jusqu'à ce que le retard en mouvement devienne acceptable (\empi).
  \item À l'arrêt, $\tau = 1/(2\pi\cdot0.6) \approx \qty{265}{\milli\second}$ : le curseur est très stable. Pour une vitesse lissée de 1~fenêtre/s, $f_c = \qty{3.6}{\hertz}$ et $\tau \approx \qty{44}{\milli\second}$.
  \item La vitesse étant elle-même lissée à \qty{1}{\hertz} ($\tau_d \approx \qty{160}{\milli\second}$), la coupure ne monte pas instantanément au début d'une saccade : le curseur rattrape le regard en quelques images au lieu d'une seule. C'est le prix de la robustesse au bruit de la vitesse.
\end{itemize}
\end{origine}

Les filtres (un par coordonnée) sont remis à zéro à la fin d'une calibration, au recentrage, au reset et au chargement d'une calibration : sans cela, le filtre ferait glisser lentement le curseur de l'ancienne échelle vers la nouvelle.

\section{Recentrage}

Le tracker garde les 20 dernières prédictions brutes (\code{raw_hist}), soit environ \qty{0.67}{\second} à \qty{30}{images\per\second}. Sur la touche Espace, \code{recenter} calcule
\begin{equation}
  \vect{o} = \begin{pmatrix}0.5\\0.5\end{pmatrix} - \frac{1}{20}\sum_{k} \hat{\vect{y}}^\text{brut}_k ,
\end{equation}
puis ajoute ce décalage à toutes les prédictions suivantes, avant le filtrage. L'utilisateur regarde le centre de l'écran et appuie sur Espace : le regard est translaté pour que ce point devienne le centre. C'est une correction grossière, valable à la pose courante ; dans la v2, la recalibration au clic la complète par une correction des angles (chapitre~\ref{chap:recal}). La moyenne sur 20 images (\empi) réduit le bruit d'un facteur $\sqrt{20} \approx 4.5$.

\section{Repli avant calibration}

Tant qu'aucun modèle n'est ajusté, les deux trackers utilisent \code{fallback_gaze} (chapitre~\ref{chap:v1}) : le curseur suit grossièrement l'iris, ce qui suffit pour s'installer face à l'écran. Le statut \enquote{Non calibré · touche C} reste affiché.

\section{Fraîcheur du regard}

\code{last_valid} mémorise l'instant de la dernière mesure exploitable. \code{gaze_fresh(max_age)} renvoie vrai si elle date de moins de \code{max_age} secondes ; la boucle principale l'appelle avec \code{delai_perte_regard}.

\begin{origine}
\code{delai_perte_regard = 0.5} s (\deriv{} / \empi). Ce délai doit être plus long qu'un clignement ordinaire suivi de son maintien : environ \qtyrange{0.25}{0.35}{\second} plus \qty{0.15}{\second}. Ainsi, un clignement pendant une fixation \emph{n'interrompt pas} le chrono de sélection. En revanche, si le visage disparaît (l'utilisateur tourne la tête, sort du champ), la sélection en cours est annulée en une demi-seconde, et le curseur devient gris.
\end{origine}

Après une calibration, \code{last_valid} est remis à \code{None} : le regard n'est considéré comme frais qu'à partir de la première mesure faite avec le nouveau modèle.

\chapter{Boucle d'interaction}
\label{chap:interaction}

Ce chapitre décrit \fichier{main_v2.py} et la géométrie de l'écran (\code{Disposition} dans \fichier{interface.py}), c'est-à-dire tout ce qui transforme le regard en clics.

\section{Démarrage}

\begin{enumerate}
  \item \textbf{Arguments.} Le premier argument qui ne commence pas par \code{--} est le chemin de la configuration ; \code{--souris} active le mode souris. Sans chemin, \code{config.chercher} prend le premier fichier existant parmi \fichier{config.yaml}, \fichier{config.yml}, \fichier{config.toml} et \fichier{config.json}, à côté du script.
  \item \textbf{Configuration.} \code{config.charger} ; en cas d'\code{ErreurConfig}, la liste complète des erreurs est affichée et le programme s'arrête avec le code 1.
  \item \textbf{Fenêtre.} pygame est initialisé ; la fenêtre est en plein écran (\code{FULLSCREEN}, résolution native) ou fenêtrée et redimensionnable, puis \code{Rendu.construire} calcule la disposition.
  \item \textbf{Tracker.} \code{TrackerSouris}, ou \code{utils_v2.GazeTracker} configuré à partir de la section \code{avance} (caméra, résolution, seuil de clignement, filtre, recalibration). L'import de \code{utils_v2}, et donc de MediaPipe et d'OpenCV, n'a lieu qu'à ce moment-là : il prend plusieurs secondes, et le mode souris s'en passe.
  \item \textbf{Boucle.} Après \code{tracker.start()}, création de \code{Navigation}, \code{Selecteur} et \code{Executeur}, puis boucle tant que \code{tracker.alive}. Un bloc \code{finally} garantit la fermeture du tracker (thread, caméra, MediaPipe) et de pygame, même après une exception.
\end{enumerate}

\section{Navigation dans les menus}

\code{Navigation} gère une \emph{pile} de couples \code{[menu, page]}, qui vaut initialement \code{[["principal", 0]]}.

\begin{itemize}
  \item \textbf{Pagination.} Une page contient $L\times C$ boutons, avec $L\times C$ la taille de la grille. Le nombre de pages d'un menu de $n$ boutons est $\max(1, \lceil n/(LC) \rceil)$, et \code{boutons()} renvoie la page courante complétée par des \code{None} (cases vides). Le bouton Retour automatique des sous-menus compte dans $n$.
  \item \textbf{Défilement circulaire.} \code{defiler(s)} fait $\text{page} \leftarrow (\text{page} + s) \bmod n_\text{pages}$ : depuis la dernière page, la bande droite ramène à la première. L'utilisateur n'est jamais bloqué en bout de menu.
  \item \textbf{Pile.} \code{ouvrir} empile un menu (page 0), \code{retour} dépile (jamais en dessous du menu principal), \code{accueil} ne garde que le menu principal. Chaque niveau mémorise sa page : revenir d'un sous-menu ramène à la page qu'on avait quittée.
  \item \textbf{Fil d'Ariane.} \code{fil()} affiche par exemple \enquote{Accueil › Raccourcis › Navigateur} dans l'en-tête.
\end{itemize}

\section{Sélection par fixation}

\subsection{Le problème du toucher de Midas}

Avec un dispositif de pointage par le regard, on regarde tout ce que l'on lit. Si regarder suffisait à cliquer, chaque coup d'œil déclencherait une action : c'est le problème du \enquote{toucher de Midas} \cite{jacob1991}. La solution classique est la \emph{fixation prolongée} (\emph{dwell time}) : une cible ne se déclenche qu'après un temps de fixation continue. En saisie de texte par le regard, les utilisateurs entraînés descendent à des durées de l'ordre de \qtyrange{0.5}{1}{\second}, et les débutants ont besoin de plus \cite{majaranta2002}.

\begin{origine}
\code{duree_clic = 2.0} s (\empi, bornes \qtyrange{0.2}{30}{\second}). C'est une valeur prudente pour un utilisateur débutant et une précision webcam modeste. Le premier prototype utilisait \qty{0.8}{\second} (\code{DWELL_SECONDS}), plus vif, mais avec un risque nettement plus élevé de clics involontaires. Chaque bouton peut avoir sa propre durée : dans \fichier{config.yaml}, \enquote{Annuler} et \enquote{Recalibrer} demandent \qty{3}{\second}, parce que ce sont des actions à conséquence.
\end{origine}

\subsection{Machine à états du sélecteur}

\begin{figure}[H]
\centering
\begin{tikzpicture}[node distance=14mm and 30mm, font=\small, >={Stealth[length=2.2mm]},
  etat/.style={draw, rounded corners=6pt, minimum width=30mm, minimum height=10mm, align=center, fill=white},
  lab/.style={font=\scriptsize, align=center, fill=white, inner sep=1.5pt}]
  \node[etat] (rien) {Aucune cible\\{\scriptsize zone morte, vide, regard perdu}};
  \node[etat, right=of rien, fill=accent!8] (fix) {Fixation\\{\scriptsize anneau $= (t - t_0)/D$}};
  \node[etat, below=of fix, fill=orange!8] (bloq) {Bloqué\\{\scriptsize anneau masqué}};
  \draw[->] (rien) to[bend left=12] node[lab, above] {entrée sur une cible\\$t_0 \leftarrow t$} (fix);
  \draw[->] (fix) to[bend left=12] node[lab, below] {sortie} (rien);
  \draw[->] (fix) -- node[lab, right] {$t - t_0 \geq D$\\\textbf{clic}} (bloq);
  \draw[->] (bloq) -| node[lab, pos=0.25, below] {changement de cible} (rien);
  \draw[->] (fix) edge[loop right] node[lab, right] {bande : clic,\\puis $t_0 \leftarrow t$\\(répétition)} (fix);
\end{tikzpicture}
\caption{Machine à états de \code{Selecteur}. $D$ est la durée de clic de la cible.}
\label{fig:selecteur}
\end{figure}

\code{Selecteur.update(cible, now, duree, repeter)} :
\begin{itemize}
  \item si la cible change, on la mémorise, $t_0 \leftarrow \text{now}$, et on lève le blocage ;
  \item si la durée vaut \code{None} (cible non cliquable) ou si la cible est bloquée, rien ne se passe ;
  \item si $\text{now} - t_0 \geq D$ : clic. Pour une bande (\code{repeter=True}), le chrono repart aussitôt, et rester sur la bande fait défiler une page toutes les $D$ secondes. Pour un bouton, la cible est bloquée : il faut la quitter puis revenir pour recliquer, ce qui empêche les rafales.
\end{itemize}
\code{progression(now, duree)} renvoie $\min((\text{now} - t_0)/D, 1)$, ou 0 si la cible est bloquée ou absente ; c'est le remplissage de l'anneau.

\code{bloquer()} est appelé après chaque action de navigation. Quand un sous-menu s'ouvre, un nouveau bouton apparaît exactement sous le regard. Sans blocage, il se déclencherait $D$ secondes plus tard sans que l'utilisateur l'ait choisi.

\section{Géométrie de l'écran}
\label{sec:disposition}

\code{Disposition(w, h, ecran)} calcule toutes les zones en pixels, une fois par taille de fenêtre. La \emph{même} géométrie sert au dessin et à la détection : ce qui est dessiné comme un bouton est exactement ce qui se déclenche.

\begin{figure}[H]
\centering
\begin{tikzpicture}[x=0.11mm, y=0.11mm, font=\scriptsize]
  % écran 1280x720
  \draw[thick] (0,0) rectangle (1280,720);
  % en-tête
  \fill[black!5] (0,720) rectangle (1280,677);
  \node at (640,698) {en-tête : fil d'Ariane · pages · statut};
  % bandes (12 %)
  \fill[accent!12] (9,9) rectangle (148,677);
  \fill[accent!12] (1132,9) rectangle (1271,677);
  \node[align=center] at (78,343) {bande\\gauche\\$\langle$};
  \node[align=center] at (1202,343) {bande\\droite\\$\rangle$};
  % cases 2x2
  \draw[rounded corners=3pt] (158,365) rectangle (635,677);
  \draw[rounded corners=3pt] (645,365) rectangle (1122,677);
  \draw[rounded corners=3pt] (158,9) rectangle (635,355);
  \draw[rounded corners=3pt] (645,9) rectangle (1122,355);
  \node at (300,620) {case 0}; \node at (985,620) {case 1};
  \node at (300,60) {case 2}; \node at (985,60) {case 3};
  % zone morte 30% x 40%
  \fill[black!6, draw=black!40] (640,360) ellipse (192 and 144);
  \node[align=center] at (640,360) {zone morte\\$30\,\% \times 40\,\%$};
  % cotes
  \draw[<->] (0,-30) -- (153,-30) node[midway, below] {$b$};
  \draw[<->] (153,-30) -- (640,-30) node[midway, below] {$c_w$};
  \draw[<->] (1310,0) -- (1310,360) node[midway, right] {$c_h$};
\end{tikzpicture}
\caption{Disposition en grille $2\times2$ avec les réglages de \fichier{config.yaml} (bandes de \qty{12}{\percent}, zone morte de $30 \times \qty{40}{\percent}$), pour une fenêtre $1280\times720$.}
\label{fig:disposition}
\end{figure}

\subsection{Dimensions}

\begin{align}
  b &= \operatorname{round}\big(\code{largeur_bandes}/100 \cdot w\big) && \text{largeur d'une bande}\\
  m &= \max\big(4, \operatorname{round}(\num{0.012}\min(w, h))\big) && \text{marge entre éléments}\\
  e &= \max\big(28, \operatorname{round}(\num{0.06}\,h)\big) && \text{hauteur de l'en-tête}\\
  c_w &= (w - 2b)/C, \qquad c_h = h/L && \text{taille d'une case}
\end{align}
La case $(r, c)$ occupe horizontalement $[b + c\,c_w,\ b + (c+1)\,c_w]$ et verticalement $[r\,c_h,\ (r+1)\,c_h]$, moins des demi-marges entre cases. La première ligne commence sous l'en-tête. Les valeurs \num{0.012}, \num{0.06} et les minimums en pixels (\empi) gardent des proportions constantes d'une résolution à l'autre, sans descendre sous des tailles lisibles.

\subsection{Zone morte}

La zone morte est un ovale centré en $(c_x, c_y) = (w/2, h/2)$, de demi-axes
\begin{equation}
  a = \frac{\code{largeur}}{200}\,w, \qquad b_z = \frac{\code{hauteur}}{200}\,h
\end{equation}
(les réglages sont des \emph{largeurs totales} en pourcentage). La moitié basse est une ellipse ; la moitié haute est une \emph{super-ellipse}, ou courbe de Lamé :
\begin{equation}
  \text{bas } (y \geq c_y) : \Big(\frac{x - c_x}{a}\Big)^2 + \Big(\frac{y - c_y}{b_z}\Big)^2 \leq 1,
  \qquad
  \text{haut } (y < c_y) : \Big|\frac{x - c_x}{a}\Big|^n + \Big|\frac{y - c_y}{b_z}\Big|^n \leq 1,
\end{equation}
avec $n = 2 + 8\,\code{aplatissement_haut}$. Pour $n = 2$, on retrouve l'ellipse ; plus $n$ augmente, plus la courbe se rapproche d'un rectangle. Avec un aplatissement de 1 ($n = 10$), le dessus est presque plat.

\paragraph{Pourquoi aplatir le haut ?} En grille $2\times2$, l'ovale mord sur le bas des deux cases du haut. Un dessus plat libère de la hauteur utile pour leur contenu, tout en gardant la forme arrondie en bas, là où le regard \enquote{tombe} naturellement au repos.

\paragraph{Pourquoi une zone morte ?} L'utilisateur doit pouvoir reposer son regard quelque part sans rien déclencher. Le centre de l'écran est l'endroit le plus confortable pour cela, et c'est aussi là que le regard revient spontanément entre deux sélections. Elle sert également de tampon entre les cases, où l'erreur du tracker ferait hésiter entre deux cibles voisines.

Le dessin de la zone morte ajoute autour d'elle un liseré de la couleur du fond, de largeur $m$, qui la sépare visuellement des cartes. La détection, elle, utilise l'ovale exact.

\subsection{\texorpdfstring{Détection : \code{zone_at}}{Détection : zone_at}}

\code{zone_at(x, y)} renvoie, par ordre de priorité :
\begin{enumerate}
  \item \code{MORTE} si $(x, y)$ est dans la zone morte ;
  \item \code{GAUCHE} si $x < b$, \code{DROITE} si $x \geq w - b$ ;
  \item sinon l'indice de case $r\,C + c$, avec $c = \operatorname{clip}(\lfloor (x - b)/c_w \rfloor, 0, C - 1)$ et $r = \operatorname{clip}(\lfloor y/c_h \rfloor, 0, L - 1)$.
\end{enumerate}
Les indices sont bornés pour qu'un regard légèrement hors fenêtre vise la case du bord, plutôt que rien. L'en-tête n'est pas exclu : le regarder vise la case de la première ligne qui se trouve en dessous.

\subsection{Portion visible d'une carte}

Pour centrer le texte d'une carte là où il ne sera pas masqué par l'ovale, \code{_segment_visible} calcule le plus grand intervalle vertical libre sur l'axe central de la carte, d'abscisse $x_c$. Avec $\tilde a = a + m$, $\tilde b = b_z + m$ et $u = |x_c - c_x| / \tilde a$ :
\begin{itemize}
  \item si $u \geq 1$, l'axe ne coupe pas l'ovale et la carte entière est libre ;
  \item sinon, l'ovale couvre $[y_\text{haut}, y_\text{bas}]$ avec
  \begin{equation}
    y_\text{haut} = c_y - \tilde b\,(1 - u^n)^{1/n}, \qquad y_\text{bas} = c_y + \tilde b\sqrt{1 - u^2},
  \end{equation}
  et on garde le plus long des deux morceaux $[y_0, \min(y_1, y_\text{haut})]$ et $[\max(y_0, y_\text{bas}), y_1]$.
\end{itemize}

\section{Une itération de la boucle principale}

\begin{enumerate}
  \item \textbf{Horloge.} $dt \leftarrow$ \code{horloge.tick(images_par_seconde)}/1000, qui attend le temps nécessaire pour tenir la cadence ; $\text{now} \leftarrow$ \code{time.monotonic()}.
  \item \textbf{Évènements.} Fermeture ; redimensionnement (hors plein écran : reconstruction du rendu) ; touches (tableau~\ref{tab:touches}).
  \item \textbf{Commandes du thread actions.} On vide \code{vers_interface} : \code{message} affiche une bulle ; \code{calibrer} lance la calibration ; \code{ouvrir}, \code{retour} et \code{accueil} modifient la pile, puis \code{selecteur.bloquer()}.
  \item \textbf{Calibration.} Si le tracker calibre, on dessine l'écran de calibration et la consigne, on remet le sélecteur à zéro et on passe à l'itération suivante. À la sortie de la calibration, une bulle indique \enquote{Calibration terminée} ou \enquote{Calibration annulée}.
  \item \textbf{Regard.} En mode souris, on met à jour la position de la souris. $\text{ok} \leftarrow$ \code{gaze_fresh(delai_perte_regard)} ; $g_x = \operatorname{clip}(\text{gaze}_x, 0, 1)\,(w-1)$, et de même pour $g_y$.
  \item \textbf{Curseur affiché.} Lissage exponentiel indépendant de la cadence :
  \begin{equation}
    k \leftarrow 1 - e^{-dt/\tau_c},\qquad \vect{c} \leftarrow \vect{c} + k\,(\vect{g} - \vect{c}), \qquad \tau_c = \qty{50}{\milli\second}.
  \end{equation}
  La forme $1 - e^{-dt/\tau}$ donne la même dynamique quelle que soit la cadence de la boucle. L'écran (30 à 60 Hz) se rafraîchit au moins aussi vite que la caméra (30 Hz) : sans ce lissage, le curseur avancerait par sauts d'une image caméra. $\tau_c = \qty{50}{\milli\second}$ (\empi) est d'un ordre comparable à la période caméra (\qty{33}{\milli\second}) : les sauts sont adoucis sans retard perceptible.
  \item \textbf{Cible et durée}, calculées sur le regard \emph{brut} $\vect{g}$ et non sur le curseur lissé :
  \begin{itemize}
    \item bande : $D = \code{duree_clic}$ s'il y a plus d'une page, sinon \code{None} ;
    \item case avec un bouton : $D$ = durée du bouton, ou réglage global ; l'anneau prend la couleur du bouton ;
    \item zone morte, case vide ou regard perdu : \code{None}.
  \end{itemize}
  \item \textbf{Clic.} Si \code{selecteur.update} renvoie vrai : effet \enquote{pop} à la position du regard. Pour une bande : \code{defiler} et animation de glissement. Pour un bouton : point de recalibration (chapitre~\ref{chap:recal}), puis \code{executeur.lancer(etapes)} ; si le bouton ne navigue pas, une bulle affiche son texte ; avec \code{retour_apres_action} dans un sous-menu, une séquence \code{accueil} est mise en file \emph{après} celle du bouton. Le thread des actions traitant les séquences dans l'ordre, le retour à l'accueil a lieu une fois l'action terminée.
  \item \textbf{Dessin} (chapitre~\ref{chap:rendu}), avec le statut \enquote{Non calibré · touche C} tant qu'il n'y a pas de modèle.
\end{enumerate}

\begin{table}[H]
\centering\small
\begin{tabular}{@{}ll@{}}
\toprule
Touche & Effet \\
\midrule
C & lancer la calibration \\
Échap & annuler la calibration en cours, sinon quitter \\
Espace & recentrer sur le centre de la fenêtre \\
R & oublier la calibration \\
D & afficher ou masquer la vignette caméra annotée \\
F & basculer plein écran / fenêtre \\
\bottomrule
\end{tabular}
\caption{Raccourcis clavier, destinés à l'aidant.}
\label{tab:touches}
\end{table}

\section{Mode souris}

\code{TrackerSouris} imite l'interface publique du tracker : il est toujours \enquote{calibré}, \code{gaze} est la position de la souris normalisée, et \code{gaze_fresh} vaut vrai tant que la souris est dans la fenêtre. Les commandes de calibration sont sans effet. Ce mode permet de concevoir et de tester une configuration (disposition, menus, actions) sur n'importe quel ordinateur, sans caméra ni MediaPipe.

\section{Différences avec \texorpdfstring{\fichier{main.py}}{main.py}}

\fichier{main.py} contient la même boucle, presque ligne pour ligne, mais avec le tracker v1. Il n'affiche pas la consigne de calibration de la v2, ne transmet pas les paramètres de recalibration et n'envoie pas de point de recalibration au clic. Les classes \code{Navigation}, \code{Selecteur} et \code{TrackerSouris} y sont dupliquées.

\chapter{Rendu graphique}
\label{chap:rendu}

\fichier{interface.py} dessine l'interface avec pygame (SDL). Le module est \enquote{pensé pour un Raspberry Pi} : sur cette cible, SDL dessine en général en logiciel, et recopier un écran entier à chaque image coûte cher. Tout le module est organisé pour \emph{ne redessiner et ne transmettre que ce qui change}.

\section{Thèmes et dimensions}

Deux thèmes, \code{sombre} et \code{clair}, définissent 10 couleurs : fond, carte, carte de retour, contour, case vide, zone morte, texte, texte doux, piste de l'anneau et alerte. La couleur d'accent (\code{couleur_accent}) sert pour les anneaux, les contours de survol, le halo du curseur et l'indicateur de page.

\code{Rendu.construire(ecran)} recalcule toutes les tailles en proportion de la fenêtre (\empi, choisies pour rester lisibles à distance) :
\begin{center}
\small
\begin{tabular}{@{}ll@{}}
\toprule
Élément & Taille \\
\midrule
rayon des coins & $\max(6, \operatorname{round}(\num{0.022}\min(w,h)))$ \\
rayon de l'anneau de progression & $\max(16, \operatorname{round}(\num{0.032}\,h))$ \\
épaisseur de l'anneau & $\max(3, \operatorname{round}(\num{0.006}\,h))$ \\
texte des cartes (taille de départ) & $\num{0.055}\,h \times \code{taille_texte}/100$ \\
petit texte (en-tête, bulles) & $\max(12, \operatorname{round}(\num{0.42}\,e))$ \\
rayon du curseur & $\max(5, \operatorname{round}(\num{0.009}\,h))$, halo $2r + 4$ \\
\bottomrule
\end{tabular}
\end{center}

\section{Stratégie de rendu incrémental}

\subsection{Caches}

\begin{table}[H]
\centering\small
\begin{tabularx}{\textwidth}{@{}l l l L@{}}
\toprule
Cache & Clé & Limite & Contenu \\
\midrule
\code{_polices} & taille & -- & objets \code{Font} \\
\code{_textes} & (texte, taille, couleur) & 400 & surfaces de texte rendues (anti-aliasing compris) \\
\code{_icones} & (chemin, taille) & -- & icônes chargées et redimensionnées (\code{smoothscale}) \\
\code{_pages} & (menu, page) & 24 & page complète : fond et toutes les cartes au repos \\
\code{_survols} & ((menu, page), case) & 48 & une carte dans son état survolé \\
\bottomrule
\end{tabularx}
\caption{Caches du rendu. Au-delà de la limite, le cache est vidé entièrement : c'est plus simple qu'une politique LRU, et suffisant puisque les menus se reconstituent en quelques millisecondes.}
\end{table}

Rendre un texte anti-aliasé ou redimensionner une image est coûteux. Avec ces caches, une page n'est dessinée qu'une fois ; ensuite, l'afficher se réduit à une copie de surface (\code{blit}).

\subsection{Fond statique}

\code{self.fond} est une surface de la taille de l'écran qui contient tout ce qui ne bouge pas : la page courante, l'en-tête (fil d'Ariane, points de pagination, statut), les bandes dans leur état normal ou inactif, et la zone morte. Elle n'est reconstruite que si la clé (menu, page, nombre de pages, fil, statut) change.

\subsection{Rectangles sales}

\begin{algorithm}[H]
\caption{Une image de \code{Rendu.dessiner} (hors animation)}
\eIf{tout doit être redessiné (nouveau fond, fin d'animation, calibration)}{copier tout le fond à l'écran}{
  \ForEach{rectangle $r$ modifié à l'image précédente (\code{sales})}{copier $\text{fond}[r]$ à l'écran \tcp*{effacement}}
}
dessiner les éléments mobiles et collecter leurs rectangles : carte ou bande survolée, vignette caméra, bulle, anneau, pop, curseur\;
\eIf{tout}{\code{display.flip()}}{\code{display.update(sales} $\cup$ \code{rects)}}
\code{sales} $\leftarrow$ \code{rects}\;
\end{algorithm}

À chaque image, on transmet à l'écran les rectangles de l'image précédente, qu'il faut effacer, et ceux de l'image actuelle. Le curseur et l'anneau occupent quelques milliers de pixels, contre deux millions pour un écran 1080p ; le gain est de plusieurs ordres de grandeur sur la bande passante vers l'écran.

\paragraph{Superposition carte survolée / zone morte.} La carte survolée est redessinée par-dessus le fond. Si elle chevauche la zone morte, celle-ci est redessinée par-dessus, avec un découpage (\code{set_clip}) limité au rectangle de la carte, pour qu'elle reste visible au premier plan.

\section{Éléments animés}

\subsection{Anneau de progression}

L'arc de progression $p \in [0,1]$ est un polygone fermé : $n + 1$ points sur le cercle extérieur de rayon $R$, de l'angle $-\pi/2$ (en haut) à $-\pi/2 + 2\pi p$, puis les mêmes angles parcourus à l'envers sur le cercle intérieur de rayon $R - e_p$, avec $n = \max(2, \lfloor 64\,p\rfloor)$. Le nombre de segments croît avec l'arc : 64 pour un tour complet, assez pour paraître lisse. Le polygone est rempli puis bordé en anti-aliasé (\code{gfxdraw}), au-dessus d'une piste circulaire grise.

\subsection{Courbe d'animation}

Les animations utilisent la courbe \enquote{ease-out cubique}
\begin{equation}
  E(t) = 1 - (1 - t)^3, \qquad t \in [0, 1],
\end{equation}
qui part vite ($E'(0) = 3$) et ralentit jusqu'à l'arrêt ($E'(1) = 0$) : le mouvement paraît naturel et se termine en douceur.

\begin{itemize}
  \item \textbf{Glissement de page} (\qty{0.22}{\second}, \code{DUREE_GLISSEMENT}) : avec $W$ la largeur de la zone centrale et $s = \pm1$ le sens, le décalage vaut $\delta = \operatorname{round}(E(t)\,W)\,s$. L'ancienne page est dessinée en $-\delta$ et la nouvelle en $sW - \delta$, avec un découpage à la zone centrale, de sorte que les bandes et l'en-tête restent immobiles. Désactivable par \code{animations: non}.
  \item \textbf{Pop de clic} (\qty{0.35}{\second}, \code{DUREE_POP}) : un cercle de rayon $R\,(1 + 0.9\,E(t))$ et d'opacité $255\,(1 - t)$, centré sur le point du clic. C'est la confirmation visuelle qu'un clic a eu lieu.
  \item \textbf{Bulle de message} (durée $T = \code{duree_message}$) : opacité $255 \cdot \min\big(1,\ (T - \text{âge})/0.3,\ \text{âge}/0.12\big)$, soit une apparition en \qty{0.12}{\second} et une disparition sur les \qty{0.3}{\second} finales.
\end{itemize}
Ces durées (\empi) sont de l'ordre des durées d'animation habituelles dans les interfaces (\qtyrange{150}{400}{\milli\second}) : assez longues pour être perçues, assez courtes pour ne pas ralentir l'utilisateur.

\section{Cartes}

\begin{itemize}
  \item \textbf{Couleurs.} Le fond d'une carte colorée est $\text{mix}(\text{carte}, \text{couleur}, 0.3)$, avec $\text{mix}(c_1, c_2, t) = c_1 + t(c_2 - c_1)$. Au survol, on mélange encore \qty{18}{\percent} de la couleur vive (celle du bouton, ou l'accent), et un contour épais apparaît. Le bouton Retour a un fond plus sombre et un texte plus doux, pour qu'on ne le confonde pas avec une action.
  \item \textbf{Mise en page du texte.} \code{_couper} coupe le texte au mot près en respectant les \texttt{\textbackslash n}. La taille part de la taille nominale et est réduite de \qty{15}{\percent} à chaque essai, jusqu'à 6 essais, tant que l'ensemble (icône, écart, lignes) dépasse \qty{86}{\percent} de la hauteur visible ou qu'une ligne dépasse la largeur utile. Le contenu est centré verticalement dans la portion visible de la carte (section~\ref{sec:disposition}).
  \item \textbf{Icône.} Elle occupe au plus $\min(\num{0.55}\,h_\text{visible}, \num{0.6}\,\text{largeur})$ ; elle est omise si elle ferait moins de 8 pixels.
\end{itemize}

\begin{attention}
Si aucun des 6 essais ne tient, la boucle sort après une dernière réduction de la taille, mais les lignes et la hauteur totale gardent les valeurs de l'essai précédent. Le texte est alors rendu un peu plus petit que la mise en page calculée. L'effet est sans gravité (le texte reste centré et tient dans la carte), mais c'est une légère incohérence.
\end{attention}

\section{Vignette caméra}

En mode debug (touche D), l'image annotée par le tracker (contours de l'iris, coins des yeux, cadence, clignement, pose, angles) est affichée en bas à gauche, avec une largeur de $\max(160, w/5)$ pixels. La conversion (redimensionnement \code{INTER_AREA}, BGR $\to$ RGB, surface pygame) n'est refaite que si \code{frame_id} a changé : l'interface tourne à 30 images par seconde comme la caméra, mais ne convertit jamais deux fois la même image.

\section{Écran de calibration}

\code{dessiner_calibration} redessine tout l'écran à chaque image, car il n'y a presque rien à afficher :
\begin{itemize}
  \item un cercle-guide de rayon $R_c = R + 8$ autour de la cible ;
  \item un disque plein de rayon $\max(4, \operatorname{round}(0.85\,R_c\,(1 - \text{frac})))$, qui \emph{rétrécit} à mesure que le temps passe sur la cible. Il est gris pendant la stabilisation et prend la couleur d'accent pendant la collecte : l'utilisateur voit quand ses données comptent ;
  \item un point central de 3 pixels, pour un point de fixation précis ;
  \item en haut, \enquote{Calibration $i/N$ -- fixe le point (Échap = annuler)}.
\end{itemize}
Un disque qui rétrécit attire le regard vers son centre, ce qui améliore la précision de la fixation. La v1 fournit une consigne sur deux lignes, affichée au centre de l'écran : aucune cible n'y est jamais placée, puisque les cibles sont aux positions $\{0.05, 0.35, 0.65, 0.95\}$. La v2 affiche sa consigne sous le titre.

\chapter{Configuration}
\label{chap:config}

La configuration (\fichier{config.py}) est pensée pour être écrite par un aidant qui n'est pas programmeur. D'où trois principes :
\begin{itemize}
  \item tous les réglages sont optionnels ;
  \item les erreurs sont détectées \emph{au chargement}, toutes à la fois, plutôt qu'au premier clic ;
  \item chaque message en français désigne l'endroit exact du problème et propose une correction.
\end{itemize}

\section{Formats}

Le format se déduit de l'extension : \fichier{.yaml}/\fichier{.yml} (PyYAML, \code{safe_load}, qui ne construit pas d'objets Python arbitraires), \fichier{.toml} (\code{tomllib}, bibliothèque standard depuis Python 3.11) ou \fichier{.json}. La structure est la même dans les trois formats ; \fichier{victor/exemples/} contient la configuration par défaut en TOML et en JSON.

\begin{itemize}
  \item Le fichier est lu en UTF-8 avec \code{utf-8-sig}, qui ignore la marque d'ordre des octets (BOM) qu'ajoutent certains éditeurs sous Windows.
  \item Une erreur de syntaxe donne la ligne et la colonne (YAML, JSON), ou le message du lecteur TOML, ainsi qu'un conseil propre au format : indentation par espaces en YAML, guillemets en TOML, virgule finale interdite en JSON.
  \item Si le fichier n'existe pas, toutes les valeurs par défaut s'appliquent et un message le signale.
\end{itemize}

\section{Fusion avec les valeurs par défaut}

\code{_fusionner(defaut, donne, chemin, erreurs)} parcourt récursivement le dictionnaire \code{DEFAUTS}. \emph{Le type de la valeur par défaut détermine la validation} :

\begin{table}[H]
\centering\small
\begin{tabularx}{\textwidth}{@{}l L@{}}
\toprule
Type par défaut & Règle \\
\midrule
dictionnaire & fusion récursive ; une valeur non dictionnaire est une erreur \\
booléen & \code{oui}/\code{vrai}/\code{o}/\code{true} ou \code{non}/\code{faux}/\code{n}/\code{false}, sans distinction de casse \\
nombre & doit être un nombre (un booléen n'est pas accepté), dans les bornes \code{BORNES} si elles existent \\
liste & même longueur que la valeur par défaut, uniquement des nombres \\
texte & converti en chaîne ; s'il figure dans \code{CHOIX}, doit être une valeur autorisée \\
\bottomrule
\end{tabularx}
\caption{Règles de validation de \code{_fusionner}.}
\end{table}

Une clé inconnue produit une erreur avec la clé connue la plus proche (\code{difflib.get_close_matches}, fondé sur la similarité de Ratcliff-Obershelp), par exemple \enquote{\code{reglages.duree_clik} : réglage inconnu, voulais-tu dire \code{"duree_clic"} ?}. Une valeur invalide est remplacée par sa valeur par défaut \emph{et} signalée : la vérification continue, ce qui permet d'afficher toutes les erreurs en un seul passage.

\paragraph{Exclusion des booléens.} En Python, \code{True} est un \code{int}. Sans le test \code{isinstance(v, bool)}, \code{duree_clic: oui} serait accepté comme la valeur 1.

\section{Conversions spécifiques}

\begin{itemize}
  \item \code{grille} : l'expression régulière \verb|(\d+)\s*[xX]\s*(\d+)| (ainsi que le signe $\times$) accepte \enquote{2x3}, \enquote{2 X 3} ou \enquote{2×3}. Lignes et colonnes doivent être entre 1 et 4 : au-delà, les cases deviennent trop petites pour la précision d'un eye tracker webcam (\empi).
  \item \code{couleur_accent} et couleurs de boutons : un nom parmi 10 (rouge, vert, bleu, orange, jaune, violet, rose, gris, turquoise, marron) ou \texttt{\#RRGGBB}. Les valeurs hexadécimales des noms correspondent à une palette d'interface aux teintes vives, lisibles sur fond sombre comme sur fond clair (\conv).
\end{itemize}

\section{Menus et boutons}

\code{_lire_menus} construit un dictionnaire \code{nom -> liste de Bouton} :
\begin{enumerate}
  \item Le menu \code{principal} est obligatoire dès qu'il y a des menus ; s'il n'y a aucun menu, un menu principal vide est créé.
  \item Chaque sous-menu reçoit \emph{en première position} un bouton Retour (\code{texte_retour}, action \code{retour}, \code{est_retour=True}). Il apparaît donc toujours dans la première case de la première page, à la même place dans tous les sous-menus.
  \item Clés d'un bouton : \code{texte}, \code{action}, \code{icone}, \code{duree_clic}, \code{couleur}. Toute autre clé est une erreur, avec suggestion.
  \item \code{action} est une chaîne ou une liste de chaînes. Chacune est analysée par \code{actions.analyser}, qui reçoit la liste de \emph{tous} les noms de menus : un \code{ouvrir} vers un menu défini plus bas dans le fichier est accepté, une faute de frappe est détectée.
  \item \code{icone} est un chemin relatif au fichier de configuration ; l'existence du fichier est vérifiée.
  \item \code{duree_clic} doit être entre \qtyrange{0.2}{30}{\second}.
  \item Il faut au moins un texte ou une icône.
\end{enumerate}
Les messages localisent l'erreur, par exemple \enquote{\texttt{menus.Raccourcis, bouton 2 «~Navigateur~» : ...}}.

\paragraph{Le type \code{Bouton}.} C'est une \code{dataclass} avec \code{texte}, \code{etapes} (liste de \code{(nom, [arguments])}), \code{icone} (chemin absolu), \code{duree_clic}, \code{couleur} (RGB), \code{est_retour} et \code{navigue}. Ce dernier champ est calculé : vrai si la séquence contient \code{ouvrir}, \code{retour} ou \code{accueil}. Il sert à ne pas afficher de bulle pour une simple navigation.

\section{Résultat}

S'il y a au moins une erreur, \code{charger} lève une unique \code{ErreurConfig} qui les liste toutes. Sinon, elle renvoie un \code{SimpleNamespace} imbriqué (\code{cfg.ecran.zone_morte.largeur}, \code{cfg.avance.camera}…), dont l'attribut \code{menus} est le dictionnaire des boutons.

\section{Tous les réglages}

\begin{longtable}{@{}p{5.1cm} p{1.9cm} p{1.6cm} p{1.9cm} p{3.8cm}@{}}
\toprule
Clé & \fichier{config.py} & \fichier{config.yaml} & Bornes & Effet \\
\midrule
\endfirsthead
\toprule
Clé & \fichier{config.py} & \fichier{config.yaml} & Bornes & Effet \\
\midrule
\endhead
\bottomrule
\endfoot
\code{ecran.plein_ecran} & non & oui & & plein écran au démarrage \\
\code{ecran.taille_fenetre} & [1280, 720] & & 2 nombres & taille hors plein écran \\
\code{ecran.grille} & 2x2 & & 1--4 & lignes × colonnes \\
\code{ecran.largeur_bandes} & 7 & 12 & 0--30 \% & largeur de chaque bande \\
\code{ecran.zone_morte.largeur} & 22 & 30 & 0--80 \% & largeur totale \\
\code{ecran.zone_morte.hauteur} & 34 & 40 & 0--80 \% & hauteur totale \\
\code{ecran.zone_morte.aplatissement_haut} & \num{0.5} & 0 & 0--1 & forme du haut \\
\code{ecran.theme} & sombre & & sombre, clair & thème \\
\code{ecran.couleur_accent} & \#4F8EF7 & & & accent \\
\code{ecran.taille_texte} & 100 & & 30--300 \% & taille du texte \\
\code{ecran.police} & (système) & & & fichier \fichier{.ttf} \\
\code{ecran.animations} & oui & & & glissement de page \\
\code{reglages.duree_clic} & \qty{2.0}{\second} & & 0,2--30 & durée de fixation \\
\code{reglages.retour_apres_action} & non & & & retour à l'accueil \\
\code{reglages.texte_retour} & $\leftarrow$ Retour & & & texte du bouton retour \\
\code{reglages.duree_message} & \qty{1.5}{\second} & & 0--30 & durée des bulles \\
\code{avance.camera} & 0 & & & index de la webcam \\
\code{avance.resolution_camera} & [640, 480] & & 2 nombres & résolution demandée \\
\code{avance.afficher_camera} & non & & & vignette de debug \\
\code{avance.images_par_seconde} & 30 & & 5--120 & cadence de l'interface \\
\code{avance.seuil_clignement} & \num{0.6} & & 0--1 & clignements \\
\code{avance.filtre_stabilite} & \num{0.6} & & 0,01--10 & $f_{\min}$ One Euro (Hz) \\
\code{avance.filtre_reactivite} & \num{3.0} & & 0--100 & $\beta$ One Euro \\
\code{avance.delai_perte_regard} & \qty{0.5}{\second} & & 0--10 & regard perdu \\
\code{avance.recalibration_clic} & oui & & & recalibration au clic \\
\code{avance.recalibration_part_centre} & \num{0.5} & \num{0.65} & 0--1 & $\alpha$ \\
\code{avance.recalibration_poids} & \num{1.0} & & 0--20 & $W$ \\
\code{avance.recalibration_oubli} & \num{0.8} & \num{0.85} & 0--1 & $d$ \\
\code{avance.recalibration_points_max} & 20 & & 1--200 & fenêtre glissante \\
\code{avance.recalibration_fenetre} & \qty{0.5}{\second} & \qty{0.4}{\second} & 0,05--5 & $\Delta$ \\
\end{longtable}

\begin{attention}
L'en-tête de \fichier{config.yaml} affirme que \enquote{les valeurs écrites ici SONT les valeurs par défaut}. C'est faux pour huit réglages, ceux dont la colonne \fichier{config.yaml} est remplie ci-dessus. Supprimer l'une de ces lignes du YAML ne rétablit donc pas la valeur qui y était écrite. Il faudrait aligner \code{DEFAUTS} sur le YAML, ou corriger le commentaire. De même, les constantes du module \fichier{utils_v2.py} (\code{RECAL_DECAY = 0.8}, \code{RECAL_WINDOW_S = 0.5}) ne correspondent qu'aux défauts de \fichier{config.py}.
\end{attention}

\chapter{Actions}
\label{chap:actions}

\section{Registre}

Une action est une fonction Python décorée par \code{@action(nom, aide, verifier)}. Le décorateur enregistre la fonction dans le dictionnaire global \code{ACTIONS} et lui attache trois attributs :
\begin{itemize}
  \item le nom employé dans la configuration ;
  \item un texte d'aide, affiché quand le nombre d'arguments est faux ;
  \item une fonction de vérification optionnelle, \code{verifier(args, menus)}, qui renvoie un message d'erreur ou \code{None}.
\end{itemize}
La fonction reçoit le contexte d'exécution (l'\code{Executeur}) puis ses arguments, toujours sous forme de chaînes :

\begin{lstlisting}
@action("gpio_impulsion", "gpio_impulsion 17 0.5   (allume pendant 0.5 s puis éteint)",
        verifier=_verifier_gpio(("la durée", 0.01, 60)))
def gpio_impulsion(ctx, broche, secondes):
    ...
\end{lstlisting}

Ajouter une action revient à écrire une fonction : le chargeur de configuration, la validation et l'exécuteur la prennent en compte automatiquement.

\section{Analyse}

\code{analyser(texte, menus)} transforme \code{'touches CTRL+A CTRL+C'} en \code{('touches', ['CTRL+A', 'CTRL+C'])} :

\begin{enumerate}
  \item \textbf{Découpage} par \code{shlex.split}, comme dans un shell : les guillemets regroupent un argument contenant des espaces (\texttt{ecrire "Bonjour, ça va ?"}). Un guillemet non fermé donne une erreur lisible.
  \item \textbf{Nom} : s'il est inconnu, le message propose le nom le plus proche et la liste des actions disponibles.
  \item \textbf{Nombre d'arguments}, vérifié \emph{par introspection} : \code{inspect.signature(f).bind(None, *args)} lève une erreur si les arguments ne correspondent pas à la signature, en tenant compte des arguments optionnels (\code{nom=None}) et variables (\code{*combos}). Il n'y a rien à déclarer en plus de la fonction elle-même.
  \item \textbf{Contenu}, par le vérificateur propre à l'action.
\end{enumerate}

\section{Exécution}

\code{Executeur} possède deux files et un thread démon :
\begin{itemize}
  \item \code{_travaux} : les séquences à exécuter, déposées par \code{lancer(etapes)} depuis le thread principal ;
  \item \code{vers_interface} : les ordres pour l'interface (\code{message}, \code{ouvrir}, \code{retour}, \code{accueil}, \code{calibrer}), lus par la boucle principale à chaque image.
\end{itemize}

Le thread prend les séquences une par une et exécute leurs étapes dans l'ordre :
\begin{itemize}
  \item une \code{ErreurAction}, c'est-à-dire une erreur prévue (appareil injoignable, broche indisponible), affiche son message à l'écran et \emph{interrompt la séquence}. On n'enchaîne pas \enquote{allumer la lumière puis ouvrir la porte} si la première étape a échoué ;
  \item toute autre exception (un bug, un programme introuvable pour \code{lancer}) est capturée : un message \enquote{Erreur : nom} s'affiche et l'interface ne plante jamais. La séquence, elle, \emph{continue} (chapitre~\ref{chap:bilan}).
\end{itemize}

Les actions ne touchent jamais l'interface directement : pygame n'est pas conçu pour être utilisé depuis plusieurs threads. Elles déposent un ordre dans \code{vers_interface}.

\section{Catalogue}

\begin{longtable}{@{}p{3.2cm} p{4.8cm} p{6.9cm}@{}}
\toprule
Action & Syntaxe & Effet et vérification \\
\midrule
\endfirsthead
\toprule
Action & Syntaxe & Effet et vérification \\
\midrule
\endhead
\bottomrule
\endfoot
\code{touches} & \code{touches CTRL+A CTRL+C} & un ou plusieurs raccourcis ; chaque touche doit être un modificateur, une touche nommée ou un caractère unique. \textbf{Envoi non branché} (section~\ref{sec:clavier}). \\
\code{ecrire} & \code{ecrire "texte"} & tape un texte. \textbf{Non branché.} \\
\code{lancer} & \code{lancer firefox} & \code{subprocess.Popen} avec \code{start_new_session=True} et les sorties redirigées vers \code{/dev/null} : le programme est détaché et survit à l'interface. \\
\code{attendre} & \code{attendre 1.5} & pause de 0 à \qty{60}{\second}, dans le thread des actions. \\
\code{afficher} & \code{afficher "message"} & bulle à l'écran. \\
\code{ouvrir} & \code{ouvrir Raccourcis} & sous-menu ; le nom doit exister (avec suggestion sinon). \\
\code{retour}, \code{accueil} & \code{retour} & menu précédent, menu principal. \\
\code{calibrer} & \code{calibrer} & relance la calibration. \\
\code{gpio_allumer} & \code{gpio_allumer 17} & sortie numérique à l'état haut. \\
\code{gpio_eteindre} & \code{gpio_eteindre 17} & sortie à l'état bas. \\
\code{gpio_basculer} & \code{gpio_basculer 17} & inverse l'état. \\
\code{gpio_impulsion} & \code{gpio_impulsion 17 0.5} & haut pendant la durée (0,01 à \qty{60}{\second}), puis bas, en arrière-plan (\code{blink(..., n=1, background=True)}) : ne bloque pas la séquence. \\
\code{gpio_intensite} & \code{gpio_intensite 18 40} & modulation de largeur d'impulsion (PWM) de 0 à \qty{100}{\percent}. \\
\code{gpio_lire} & \code{gpio_lire 4 "Porte"} & affiche \enquote{Porte : haut/bas}. \\
\code{http_get} & \code{http_get 192.168.1.20/relay/0?turn=on} & requête GET. \\
\code{http_post}, \code{http_put} & \code{http_post 192.168.1.20/api etat=on} & POST ou PUT avec un corps JSON ou texte. \\
\code{http_lire} & \code{http_lire 192.168.1.30/temperature "Salon"} & affiche la réponse sur une ligne, tronquée à 60 caractères. \\
\end{longtable}

\section{Clavier}
\label{sec:clavier}

\begin{itemize}
  \item \textbf{Modificateurs} : CTRL, ALT, MAJ, SHIFT, SUPER, ALTGR.
  \item \textbf{Touches nommées} : ENTREE, ECHAP, TAB, ESPACE, RETOUR\_ARRIERE, SUPPR, INSER, HAUT, BAS, GAUCHE, DROITE, DEBUT, FIN, PAGE\_HAUT, PAGE\_BAS, F1 à F12.
  \item Tout caractère unique est aussi accepté.
\end{itemize}
Les noms sont en français pour l'aidant ; leur traduction vers des codes de touches incombe au \enquote{backend} d'envoi.

\begin{attention}
\code{envoyer_combo(touches)} et \code{envoyer_texte(texte)} sont des \textbf{points d'accroche vides} : ils se contentent d'afficher \code{[CLAVIER] ...} dans la console. La validation est complète, mais aucune touche n'est envoyée au système. Le commentaire du code prévoit de les brancher sur \code{wtype} ou \code{ydotool} (Wayland), ou sur \code{uinput}, selon l'environnement graphique du Raspberry Pi ; \code{xdotool} conviendrait sous X11. C'est la première priorité de la feuille de route.
\end{attention}

\section{GPIO}

\paragraph{Numérotation.} Les broches sont désignées par leur numéro BCM, celui du schéma de brochage (\enquote{GPIO17}), et non par leur position physique sur le connecteur. Les numéros valides vont de 0 à 27 : ce sont les 28 GPIO du connecteur 40 broches des Raspberry Pi (\conv). L'analyseur accepte \code{17} comme \code{GPIO17}, sans distinction de casse.

\paragraph{Chargement paresseux et simulation.} \code{gpiozero} n'est importé qu'à la première action GPIO. Si le module est absent, ou si la machine n'a pas de GPIO (exception \code{BadPinFactory}, par exemple sur un PC), on bascule définitivement en \emph{mode simulé} : chaque action est affichée dans la console (\texttt{[GPIO simulé] GPIO17 allumée}), et un dictionnaire mémorise les états pour que \code{gpio_basculer} et \code{gpio_lire} restent cohérents. Une configuration domotique peut ainsi être testée entièrement sur PC. Les avertissements de \code{gpiozero}, qui en émet un pour chaque backend essayé, sont masqués pendant la création.

\paragraph{Gestion des broches.} Un objet \code{gpiozero} est créé à la demande pour chaque broche, puis réutilisé. Si une broche change d'usage (sortie numérique, PWM, entrée), l'ancien objet est fermé avant d'en créer un nouveau : \code{gpiozero} refuse d'ouvrir deux fois la même broche. Un verrou (\code{_gpio_verrou}) protège ces structures, même si, en pratique, seul le thread des actions y accède. Sur Pi 4 et Pi 5, \code{gpiozero} utilise le backend \code{lgpio}.

\section{HTTP}

Ces actions servent à piloter des appareils du réseau local : relais Wi-Fi, prises connectées, ESP32, Home Assistant.

\begin{itemize}
  \item \textbf{Adresse} : \code{http://} est ajouté si aucun schéma n'est donné, pour qu'on puisse écrire simplement \code{192.168.1.20/on}. La vérification exige un nom d'hôte et aucune espace.
  \item \textbf{Proxy désactivé} (\code{ProxyHandler({})}) : un proxy d'entreprise configuré dans l'environnement ne doit pas intercepter des requêtes destinées au réseau local.
  \item \textbf{Délai} de \qty{3}{\second} (\code{HTTP_DELAI_S}, \empi) : un appareil local répond en quelques dizaines de millisecondes ; au-delà de \qty{3}{\second}, il est considéré comme injoignable et l'utilisateur est prévenu rapidement.
  \item \textbf{Taille lue} : au plus \qty{64}{\kilo\byte} (\code{HTTP_TAILLE_MAX}). Le but est d'afficher une valeur, pas de télécharger une page.
  \item \textbf{Corps de la requête} (\code{_corps}) :
  \begin{itemize}
    \item un seul argument commençant par \texttt{\{} ou \texttt{[} : JSON envoyé tel quel, après validation ;
    \item des paires \code{cle=valeur} : JSON construit, chaque valeur étant interprétée en JSON si possible (\code{3} devient un nombre, \code{true} un booléen) et gardée comme texte sinon ;
    \item tout autre texte : envoyé en \code{text/plain}.
  \end{itemize}
  \item \textbf{Erreurs} : un code HTTP d'erreur donne \enquote{\textit{hôte} a répondu erreur \textit{code}}, une absence de réponse donne \enquote{\textit{hôte} injoignable}. Ce sont des \code{ErreurAction} : elles s'affichent et interrompent la séquence.
  \item \textbf{En-tête} \code{User-Agent: eye-tracker}, générique.
\end{itemize}

\chapter{Prototypes, bilan et perspectives}
\label{chap:bilan}

\section{Les prototypes à la racine du dépôt}

Trois fichiers à la racine sont les étapes qui ont précédé \fichier{victor/}. L'interface ne les utilise plus.

\subsection{\texorpdfstring{\fichier{main.py}}{main.py} : sélection sans calibration}

Le premier prototype affiche une grille $3\times3$ dans une fenêtre OpenCV de $480\times480$ pixels, et suppose une tête immobile. Pour chaque œil, \code{eye_ratio} calcule :
\begin{itemize}
  \item un rapport horizontal $r_x = (p_x - x_\text{min})/(x_\text{max} - x_\text{min})$ entre les deux coins ;
  \item un rapport vertical $r_y = 0.5 + (p_y - y_\text{ligne})/h_\text{œil}$, où $y_\text{ligne}$ est l'ordonnée de la droite des coins à l'abscisse de l'iris, et $h_\text{œil}$ la hauteur paupière haute -- paupière basse ; $r_y$ est borné à $[0, 1]$.
\end{itemize}
Les rapports des deux yeux sont moyennés, puis lissés par une moyenne exponentielle de poids \num{0.35} par image (constante de temps d'environ 2,3 images). La case est obtenue par \emph{seuillage direct} : colonnes séparées en $r_x = \num{0.42}$ et \num{0.58}, lignes en $r_y = \num{0.54}$ et \num{0.58}. Une fixation de \qty{0.8}{\second} sélectionne la case, qui est mise en évidence pendant \qty{0.4}{\second}.

Ce prototype a montré les limites qui ont guidé la suite :
\begin{itemize}
  \item seuils fixes, réglés pour une personne et une position. La bande verticale centrale ne fait que \num{0.04} de large, ce qui rend la ligne du milieu presque impossible à tenir ;
  \item instabilité verticale due à la normalisation par la hauteur de l'œil (chapitre~\ref{chap:v1}) ;
  \item aucune compensation de la tête.
\end{itemize}

\subsection{\texorpdfstring{\fichier{test_xav.py}}{test_xav.py} : le tracker \enquote{précision}}

Ce programme autonome introduit tout ce que reprend la v1 : repère propre à chaque œil, normalisation par la largeur, blendshapes, compensation de la tête, calibration 16 points avec régression ridge et validation croisée LOPO, filtre One Euro, gel pendant les clignements. Il corrige aussi le signe vertical : $y$ croît vers le bas dans l'image comme à l'écran, alors qu'une version antérieure calculait $1 - y$. \fichier{victor/utils.py} en est la version découpée en briques réutilisables et exécutée dans un thread.

\subsection{\texorpdfstring{\fichier{landmarks_visu.py}}{landmarks_visu.py}}

C'est un outil de visualisation qui dessine en direct les landmarks MediaPipe sur l'image de la webcam. Il sert à choisir les indices de points et à vérifier la détection.

\section{Ce qui fonctionne}

\begin{itemize}
  \item Chaîne complète webcam $\to$ regard $\to$ sélection par fixation $\to$ action, avec deux estimateurs de regard interchangeables.
  \item Compensation géométrique des mouvements de la tête (v2), avec un modèle 3D propre à l'utilisateur et une calibration en deux phases qui sépare effets de l'œil et effets de la distance.
  \item Recalibration continue à chaque clic, ancrée sur la calibration.
  \item Interface optimisée pour Raspberry Pi : rendu incrémental, caches, threads découplés.
  \item Configuration multi-format, entièrement validée au chargement, avec messages localisés.
  \item Actions GPIO et HTTP fonctionnelles, avec un mode simulé pour tester sur PC.
  \item Mode souris pour concevoir les menus sans caméra.
\end{itemize}

\section{Limites et incohérences relevées}

\begin{table}[H]
\centering\small
\begin{tabularx}{\textwidth}{@{}l L l@{}}
\toprule
Sujet & Constat & Section \\
\midrule
Clavier & \code{touches} et \code{ecrire} n'envoient rien au système (points d'accroche vides). & \ref{sec:clavier} \\
Calibration & \code{save_calibration} et \code{load_calibration} ne sont jamais appelées : recalibration à chaque lancement. & \ref{chap:recal} \\
Valeurs par défaut & Huit réglages ont des valeurs différentes dans \fichier{config.py} et \fichier{config.yaml}, contrairement à ce qu'affirme l'en-tête du YAML. & \ref{chap:config} \\
Séquences & Une exception imprévue n'interrompt pas la séquence, contrairement à une \code{ErreurAction}. & \ref{chap:actions} \\
En-tête & L'en-tête compte comme la première ligne de cases pour la détection. & \ref{sec:disposition} \\
Mise en page & Le dernier essai de taille de texte peut ne pas correspondre à la mise en page calculée. & \ref{chap:rendu} \\
Validation v2 & L'erreur LOPO est approchée (3 itérations à chaud) et légèrement optimiste. & \ref{chap:v2mod} \\
Code dupliqué & \fichier{main.py} et \fichier{main_v2.py} partagent \code{Navigation}, \code{Selecteur}, \code{TrackerSouris} et la boucle. & \ref{chap:interaction} \\
Outillage & Pas de fichier de dépendances, pas de tests automatisés, pas de mesures de performance. & \ref{sec:perf} \\
\bottomrule
\end{tabularx}
\caption{Limites connues.}
\end{table}

\paragraph{Hypothèses du modèle géométrique.} La v2 repose sur des valeurs moyennes : champ de vision de \ang{60}, distance interpupillaire de \qty{6.3}{\centi\metre}, globe de rayon \qty{1.15}{\centi\metre} centré \qty{1}{\centi\metre} derrière les coins, écran dans le plan de la caméra, pas de distorsion de l'objectif. La calibration en absorbe l'essentiel par les gains, les décalages, la correction de distance $k_f$, l'inclinaison et l'affine. Une webcam très grand-angle (forte distorsion), une caméra éloignée de l'écran (écran externe avec webcam posée à côté) ou des lunettes à fort reflet dégradent cependant le résultat. La profondeur $z$ de MediaPipe est aussi une estimation apprise, plus bruitée que $x$ et $y$.

\paragraph{Constantes empiriques.} Une grande partie des constantes (annexe~\ref{chap:constantes}) relève de choix de conception non mesurés. Les plus influentes, et les premières à étudier, sont :
\begin{itemize}
  \item \code{duree_clic}, compromis entre vitesse et clics involontaires ;
  \item \code{FIT_SIGMA}, qui fixe le poids des a priori ;
  \item \code{recalibration_part_centre} et \code{recalibration_poids}, qui fixent la réactivité de la recalibration ;
  \item les réglages du filtre One Euro.
\end{itemize}

\section{Perspectives}

Elles reprennent la feuille de route (\fichier{ROADMAP.md}) :
\begin{enumerate}
  \item \textbf{Backend clavier} : détection automatique \code{wtype}/\code{ydotool} (Wayland), \code{xdotool} (X11) ou \code{uinput}, gestion de l'AZERTY et des caractères accentués, message clair si aucun backend n'est disponible.
  \item \textbf{Persistance de la calibration} : chargement au démarrage et sauvegarde après chaque calibration réussie.
  \item \textbf{Qualité du suivi} : afficher l'erreur validée et proposer de recommencer si elle est mauvaise ; indicateur de visage perdu, de tête trop tournée ou de mauvais éclairage ; mode \enquote{pause} et mode \enquote{sûr} quand l'utilisateur ne regarde pas l'écran.
  \item \textbf{Mesures} : protocole fixe (cibles de test après calibration, tête immobile puis en mouvement) pour comparer v1 et v2, et pour régler \code{duree_clic}, $\alpha$ (en mesurant $\sigma_U$ et $\sigma_\varepsilon$, chapitre~\ref{chap:recal}), $f_{\min}$ et $\beta$ avec de vrais utilisateurs.
  \item \textbf{Nettoyage} : un seul point d'entrée (\code{--tracker v1|v2}), classes partagées dans un module, fichier de dépendances, tests unitaires de ce qui se teste sans caméra (configuration, analyse des actions, \code{umeyama}, pose, intersection rayon -- sphère).
  \item \textbf{Déploiement sur Raspberry Pi} : mesure de la cadence réelle, lancement automatique (service systemd), script d'installation (dépendances, droits \code{uinput} et GPIO, modèle).
\end{enumerate}


\appendix
\chapter{Récapitulatif des constantes}
\label{chap:constantes}

Origines : \litt{} (publication ou physiologie), \deriv{} (conséquence mathématique), \conv{} (convention ou valeur imposée par une bibliothèque), \empi{} (choix de conception réglé à l'usage, non mesuré). La dernière colonne renvoie au chapitre où la valeur est discutée.

\small
\begin{longtable}{@{}>{\raggedright\arraybackslash}p{4.6cm} >{\raggedright\arraybackslash}p{2.6cm} >{\raggedright\arraybackslash}p{2.3cm} >{\raggedright\arraybackslash}p{4.4cm} c@{}}
\toprule
Constante & Valeur & Origine & Rôle & Ch. \\
\midrule
\endfirsthead
\toprule
Constante & Valeur & Origine & Rôle & Ch. \\
\midrule
\endhead
\bottomrule
\endfoot
\multicolumn{5}{@{}l}{\textbf{Caméra et MediaPipe} (\fichier{utils.py}, \fichier{utils_v2.py})}\\
\code{CAM_WIDTH}, \code{CAM_HEIGHT} & $640\times480$ & \empi & compromis coût / précision sur Pi & \ref{chap:archi} \\
\code{CAM_FPS} & 30 & \conv & cadence standard des webcams & \ref{chap:archi} \\
\code{CAP_PROP_BUFFERSIZE} & 1 & \deriv & latence minimale : toujours l'image la plus récente & \ref{chap:archi} \\
confiances MediaPipe & \num{0.5} & \conv & valeurs par défaut de la bibliothèque & \ref{chap:mediapipe} \\
\code{num_faces} & 1 & \conv & un utilisateur ; lissage interne actif & \ref{chap:mediapipe} \\
indices \code{EYE_A}, \code{EYE_B}, \code{NOSE_TIP} & voir tableau & \conv & maillage canonique MediaPipe & \ref{chap:mediapipe} \\
\midrule
\multicolumn{5}{@{}l}{\textbf{Calibration commune}}\\
\code{CAL_COLS}, \code{CAL_ROWS} & $4\times4$ & \empi & 16 cibles & \ref{chap:v1} \\
\code{CAL_MARGIN} & \num{0.05} & \empi & cibles à \qty{5}{\percent} des bords & \ref{chap:v1} \\
\code{CAL_LEAD_S} (v1) & \qty{2.5}{\second} & \empi & lecture de la consigne & \ref{chap:v1} \\
\code{CAL_SETTLE_S} & \qty{0.8}{\second} & \empi{} / \litt & saccade et stabilisation & \ref{chap:v1} \\
\code{CAL_COLLECT_S} & \qty{2.0}{\second} & \empi & $\approx 60$ images par cible & \ref{chap:v1} \\
\code{CAL_MIN_SAMPLES} & 15 & \empi & \qty{25}{\percent} des images & \ref{chap:v1} \\
\code{CAL_MIN_POINTS} & 10 & \empi & couverture spatiale minimale & \ref{chap:v1} \\
reprises par cible & 2 & \empi & puis cible ignorée & \ref{chap:v1} \\
\midrule
\multicolumn{5}{@{}l}{\textbf{Tracker v1} (\fichier{utils.py})}\\
\code{RIDGE_LAMBDAS} & \num{0.002} \dots{} 1 & \empi & grille log, pas $\approx\times5$ & \ref{chap:v1} \\
\code{SD_FLOORS} & 13 valeurs & \empi & $\approx$ bruit de chaque feature & \ref{chap:v1} \\
\code{PENALTY} & 1 ; 4 (quadr.) & \empi & courbures plus pénalisées & \ref{chap:v1} \\
\code{HEAD_CLIP_Z} & 3 & \deriv & $\pm3\sigma$ : \qty{99.7}{\percent} d'une gaussienne & \ref{chap:v1} \\
\code{MAX_YAW_DEG}, \code{MAX_PITCH_DEG} & \ang{40}, \ang{35} & \empi & iris non fiable au-delà & \ref{chap:v1} \\
\code{CAL_HEAD_MIN_SD} & \ang{0.8} & \empi & alerte tête immobile & \ref{chap:v1} \\
seuil d'alerte LOPO & \qty{12}{\percent} & \empi & $\approx$ un tiers de case & \ref{chap:v1} \\
rejet par cible & 5 MAD, plancher \num{0.01} & \empi & aberrants francs seulement & \ref{chap:v1} \\
rejet sur résidus & \num{3.5} MAD, \qty{70}{\percent} & \litt{} / \empi & z-score modifié & \ref{chap:v1} \\
LOPO : entraînement minimal & 20 échantillons & \empi & & \ref{chap:v1} \\
\code{FALLBACK_GAIN_X}, \code{_Y} & 3, 6 & \empi & curseur approximatif & \ref{chap:v1} \\
repli clignement $a_p$ & \num{0.08} & \empi & œil presque fermé & \ref{chap:v1} \\
\midrule
\multicolumn{5}{@{}l}{\textbf{Statistiques robustes}}\\
constante de cohérence & \num{1.4826} & \deriv & $1/\Phi^{-1}(3/4)$ & \ref{chap:v1} \\
seuil d'aberrance & \num{3.5} & \litt & Iglewicz et Hoaglin & \ref{chap:v1} \\
\midrule
\multicolumn{5}{@{}l}{\textbf{Géométrie v2} (\fichier{utils_v2.py})}\\
\code{CAM_HFOV_DEG} & \ang{60} & \litt{} / \empi & webcam typique ; $f \approx \qty{554}{px}$ & \ref{chap:v2geo} \\
\code{IPD_CM} & \qty{6.3}{\centi\metre} & \litt & IPD adulte moyenne & \ref{chap:v2geo} \\
\code{EYE_RADIUS_CM} & \qty{1.15}{\centi\metre} & \litt{} / \empi & centre de rotation $\to$ pupille & \ref{chap:v2geo} \\
\code{EYE_DEPTH_CM} & \qty{1.0}{\centi\metre} & \litt{} / \empi & centre derrière les coins & \ref{chap:v2geo} \\
\code{MAX_HEAD_ANGLE_DEG} & \ang{50} & \empi & un œil masqué au-delà & \ref{chap:v2geo} \\
\code{RIGID_IDX} & 69 points & \empi & régions osseuses & \ref{chap:v2geo} \\
\code{MODEL_MAX_FRAMES} & 150 & \empi & coût du modèle de tête & \ref{chap:v2geo} \\
itérations Procrustes & 4 & \empi & convergence en 2--3 & \ref{chap:v2geo} \\
IRLS pose & 2 passes, $c = \num{2.5}\hat\sigma$ & \litt{} / \empi & poids de Cauchy ($\num{2.385}$ classique) & \ref{chap:v2geo} \\
\midrule
\multicolumn{5}{@{}l}{\textbf{Modèle de regard v2}}\\
\code{N_PARAMS} & 23 & \deriv & 6 + 7 + 7 + 2 + 1 & \ref{chap:v2mod} \\
\code{FIT_SIGMA} & \num{0.03} & \empi & bruit par échantillon & \ref{chap:v2mod} \\
\code{PRIOR_SIGMA} & voir tableau & \empi & a priori faibles & \ref{chap:v2mod} \\
borne du dénominateur & $-\num{0.01}$ & \empi & regard parallèle à l'écran & \ref{chap:v2mod} \\
LM : $\mu_0$, $\times4$, $/3$ & $10^{-3}$, 4, 3 & \empi & amortissement & \ref{chap:v2mod} \\
LM : $\mu_\text{max}$, $\mu_\text{min}$ & $10^{10}$, $10^{-9}$ & \empi & bornes & \ref{chap:v2mod} \\
LM : pas de dérivation & $10^{-6}\max(1,|p|)$ & \deriv & $\approx$ racine de la précision machine & \ref{chap:v2mod} \\
LM : tolérance, itérations & $10^{-8}$ ; 60 / 15 / 3 & \empi & calibration / recal / LOPO & \ref{chap:v2mod} \\
\code{CAL_HEAD_TARGETS} & 3 cibles & \empi & centre et diagonale & \ref{chap:v2mod} \\
\code{CAL_HEAD_SETTLE_S}, \code{_COLLECT_S} & \qty{1.0}{\second}, \qty{4.5}{\second} & \empi & mouvement complet de la tête & \ref{chap:v2mod} \\
\code{CAL_HEAD_MIN_SAMPLES} & 30 & \empi & & \ref{chap:v2mod} \\
\code{HEAD_GROUP_WEIGHT} & 3 & \empi & 3 cibles face à 16 & \ref{chap:v2mod} \\
rejet angles (cibles fixes) & \num{3.5} MAD, plancher \qty{0.003}{\radian} & \litt{} / \empi & saccades, clignements & \ref{chap:v2mod} \\
\midrule
\multicolumn{5}{@{}l}{\textbf{Recalibration}}\\
\code{RECAL_PARAMS} & [6, 13, 7, 15] & \deriv & décalages et gains $\theta$, $\varphi$ & \ref{chap:recal} \\
\code{RECAL_SIGMA} & [0,15 ; 0,15 ; 0,2 ; 0,2] & \empi & ancre sur la calibration & \ref{chap:recal} \\
\code{RECAL_WEIGHT} & 1 & \empi & poids du dernier clic & \ref{chap:recal} \\
\code{RECAL_DECAY} & \num{0.8} (YAML \num{0.85}) & \empi & 5 (6,7) clics effectifs & \ref{chap:recal} \\
\code{RECAL_MAX_POINTS} & 20 & \deriv{} / \empi & poids au-delà $< \qty{5}{\percent}$ & \ref{chap:recal} \\
\code{RECAL_WINDOW_S} & \qty{0.5}{\second} (YAML \num{0.4}) & \empi & cœur de la fixation & \ref{chap:recal} \\
images minimales par point & 3 & \empi & & \ref{chap:recal} \\
historique \code{_recent} & 300 & \empi & $\approx\qty{10}{\second}$ & \ref{chap:recal} \\
part du centre $\alpha$ & \num{0.5} (YAML \num{0.65}) & \empi & estimateur à rétrécissement & \ref{chap:recal} \\
\midrule
\multicolumn{5}{@{}l}{\textbf{Traitement du signal}}\\
\code{BLINK_THRESHOLD} & \num{0.6} & \empi & seuil sur \code{eyeBlink} & \ref{chap:traitements} \\
\code{BLINK_HOLD_S} & \qty{0.15}{\second} & \litt{} / \empi & réouverture de la paupière & \ref{chap:traitements} \\
\code{ONE_EURO_MIN_CUTOFF} & \qty{0.6}{\hertz} & \empi & stabilité au repos & \ref{chap:traitements} \\
\code{ONE_EURO_BETA} & 3 & \empi & réactivité & \ref{chap:traitements} \\
\code{ONE_EURO_D_CUTOFF} & \qty{1}{\hertz} & \litt & recommandation des auteurs & \ref{chap:traitements} \\
historique de recentrage & 20 prédictions & \empi & bruit $/\sqrt{20}$ & \ref{chap:traitements} \\
\code{delai_perte_regard} & \qty{0.5}{\second} & \deriv{} / \empi & $>$ clignement + maintien & \ref{chap:traitements} \\
moyenne de cadence & $0.9 / 0.1$ & \empi & $\approx 10$ images & \ref{chap:archi} \\
\midrule
\multicolumn{5}{@{}l}{\textbf{Interaction et rendu}}\\
\code{duree_clic} & \qty{2}{\second} & \empi{} / \litt & toucher de Midas & \ref{chap:interaction} \\
lissage du curseur $\tau_c$ & \qty{50}{\milli\second} & \empi & $\approx$ période caméra & \ref{chap:interaction} \\
marge, en-tête & \num{0.012}, \num{0.06} & \empi & proportions de l'écran & \ref{chap:interaction} \\
exposant super-ellipse & $2 + 8a$ & \empi & 2 (ellipse) à 10 (plat) & \ref{chap:interaction} \\
\code{DUREE_GLISSEMENT} & \qty{0.22}{\second} & \empi & animation de page & \ref{chap:rendu} \\
\code{DUREE_POP} & \qty{0.35}{\second} & \empi & confirmation du clic & \ref{chap:rendu} \\
bulle : apparition, disparition & \qty{0.12}{\second}, \qty{0.3}{\second} & \empi & & \ref{chap:rendu} \\
mélanges de couleur & \qty{30}{\percent}, \qty{18}{\percent} & \empi & bouton coloré, survol & \ref{chap:rendu} \\
réduction du texte & $\times\num{0.85}$, 6 essais & \empi & ajustement à la carte & \ref{chap:rendu} \\
limites des caches & 400, 24, 48 & \empi & mémoire bornée & \ref{chap:rendu} \\
segments de l'anneau & $\max(2, 64p)$ & \empi & arc lisse & \ref{chap:rendu} \\
\midrule
\multicolumn{5}{@{}l}{\textbf{Actions}}\\
\code{BROCHES_MAX} & 27 & \conv & GPIO 0--27 du connecteur 40 broches & \ref{chap:actions} \\
\code{HTTP_DELAI_S} & \qty{3}{\second} & \empi & appareil local & \ref{chap:actions} \\
\code{HTTP_TAILLE_MAX} & \qty{64}{\kilo\byte} & \empi & lecture d'une valeur & \ref{chap:actions} \\
troncature \code{http_lire} & 60 caractères & \empi & une ligne à l'écran & \ref{chap:actions} \\
\code{attendre} & 0--\qty{60}{\second} & \empi & & \ref{chap:actions} \\
\end{longtable}
\normalsize


\begin{thebibliography}{99}

\bibitem{ablavatski2020}
A.~Ablavatski, A.~Vakunov, I.~Grishchenko, K.~Raveendran, M.~Zhdanovich.
\newblock Real-time pupil tracking from monocular video for digital puppetry.
\newblock \emph{arXiv:2006.11341}, 2020.

\bibitem{casiez2012}
G.~Casiez, N.~Roussel, D.~Vogel.
\newblock 1\,€ filter: a simple speed-based low-pass filter for noisy input in interactive systems.
\newblock In \emph{Proceedings of CHI 2012}, ACM, 2012.

\bibitem{dodgson2004}
N.~A.~Dodgson.
\newblock Variation and extrema of human interpupillary distance.
\newblock In \emph{Stereoscopic Displays and Virtual Reality Systems XI}, Proc. SPIE 5291, 2004.

\bibitem{gower1975}
J.~C.~Gower.
\newblock Generalized Procrustes analysis.
\newblock \emph{Psychometrika}, 40(1), 1975.

\bibitem{guestrin2006}
E.~D.~Guestrin, M.~Eizenman.
\newblock General theory of remote gaze estimation using the pupil center and corneal reflections.
\newblock \emph{IEEE Transactions on Biomedical Engineering}, 53(6), 2006.

\bibitem{hansen2010}
D.~W.~Hansen, Q.~Ji.
\newblock In the eye of the beholder: a survey of models for eyes and gaze.
\newblock \emph{IEEE Transactions on Pattern Analysis and Machine Intelligence}, 32(3), 2010.

\bibitem{hoerl1970}
A.~E.~Hoerl, R.~W.~Kennard.
\newblock Ridge regression: biased estimation for nonorthogonal problems.
\newblock \emph{Technometrics}, 12(1), 1970.

\bibitem{holland1977}
P.~W.~Holland, R.~E.~Welsch.
\newblock Robust regression using iteratively reweighted least-squares.
\newblock \emph{Communications in Statistics -- Theory and Methods}, 6(9), 1977.

\bibitem{iglewicz1993}
B.~Iglewicz, D.~C.~Hoaglin.
\newblock \emph{How to Detect and Handle Outliers}.
\newblock ASQC Quality Press, 1993.

\bibitem{jacob1991}
R.~J.~K.~Jacob.
\newblock The use of eye movements in human-computer interaction techniques: what you look at is what you get.
\newblock \emph{ACM Transactions on Information Systems}, 9(2), 1991.

\bibitem{kartynnik2019}
Y.~Kartynnik, A.~Ablavatski, I.~Grishchenko, M.~Grundmann.
\newblock Real-time facial surface geometry from monocular video on mobile GPUs.
\newblock \emph{arXiv:1907.06724}, 2019.

\bibitem{levenberg1944}
K.~Levenberg.
\newblock A method for the solution of certain non-linear problems in least squares.
\newblock \emph{Quarterly of Applied Mathematics}, 2(2), 1944.

\bibitem{majaranta2002}
P.~Majaranta, K.-J.~Räihä.
\newblock Twenty years of eye typing: systems and design issues.
\newblock In \emph{Proceedings of ETRA 2002}, ACM, 2002.

\bibitem{marquardt1963}
D.~W.~Marquardt.
\newblock An algorithm for least-squares estimation of nonlinear parameters.
\newblock \emph{Journal of the Society for Industrial and Applied Mathematics}, 11(2), 1963.

\bibitem{rousseeuw1993}
P.~J.~Rousseeuw, C.~Croux.
\newblock Alternatives to the median absolute deviation.
\newblock \emph{Journal of the American Statistical Association}, 88(424), 1993.

\bibitem{stone1974}
M.~Stone.
\newblock Cross-validatory choice and assessment of statistical predictions.
\newblock \emph{Journal of the Royal Statistical Society, Series B}, 36(2), 1974.

\bibitem{umeyama1991}
S.~Umeyama.
\newblock Least-squares estimation of transformation parameters between two point patterns.
\newblock \emph{IEEE Transactions on Pattern Analysis and Machine Intelligence}, 13(4), 1991.

\end{thebibliography}


\end{document}
